\documentclass[12pt, ukrainian]{article}
\usepackage[a4paper]{geometry}
\setlength\parindent{0mm}
\geometry{top=1.1cm,bottom=1.1cm,left=1.0cm,right=1.0cm}
\pagestyle{empty}
\usepackage[T2A]{fontenc}
\usepackage[utf8]{inputenc}
\usepackage[ukrainian]{babel}
\usepackage{graphicx}
\usepackage{caption}
\DeclareCaptionFormat{nocaption}{}
\captionsetup{format=nocaption,aboveskip=0pt,belowskip=0pt}
\usepackage[Export]{adjustbox} 
\adjustboxset{max size={0.9\linewidth}{0.9\paperheight}}
\begin{document}
%begin
%\begin {Huge}
{{27.11.2024 Контрольна робота \No 2, варіант  \No \textbf { 1 }\par}}

%\begin{center}Частина 2
%\end{center}
%\vspace{10mm}
%\end {Huge}

{
Якість коду та економність обчислень оцінюється по всім завданням. Послідовності задаються масивами. Використовувати додаткові структури даних (у тому числі рядки) та функції стандартної бібліотеки заборонено.

\vspace{2mm}

\begin {large}
%beginitem
1. Написати функцію, що знаходить у скількох цілих точках інтервалу $(a; b]$ функція~$cos(x)$ приймає додатнє значення.

\vspace{2mm}

%enditem

%beginitem
2. Написати функцію, що замінює всі додатні елементи послідовності цілих чисел числом 15.

\vspace{2mm}

%enditem

%beginitem
3. Написати функцію, що перевіряє, чи містить С-рядок хоч одну велику латинську літеру.

\vspace{2mm}

%enditem

%beginitem
4. 
Написати функцію, що записує в масив дійсних послідовність чисел
$$\scalebox{1.2}{$C_{2n+k+1}^{n+k}$}, \quad k=0,1,\ldots,n$$
Використовувати вкладені цикли та виклики функцій \textbf{заборонено}.

\vspace{2mm}

%enditem

%beginitem
5. Написати функцію, що для цілого числа знаходить суму його простих дільників, що входять у його розклад не більш ніж у 4му степеню. \\ Наприклад, для числа $3^5{\cdot}7^9{\cdot}11{\cdot}23^{2}$ таких дільників два – $11$ та $23$, тому функція має повернути $34$.

\vspace{2mm}

%enditem
\end{large}
%end
\vspace{35mm}


%begin
%\begin {Huge}
{{27.11.2024 Контрольна робота \No 2, варіант  \No \textbf { 2 }\par}}

%\begin{center}Частина 2
%\end{center}
%\vspace{10mm}
%\end {Huge}

{
Якість коду та економність обчислень оцінюється по всім завданням. Послідовності задаються масивами. Використовувати додаткові структури даних (у тому числі рядки) та функції стандартної бібліотеки заборонено.

\vspace{2mm}

\begin {large}
%beginitem
1. Написати функцію, що знаходить у скількох цілих точках інтервалу $(a; b)$ значення функції~$sin(x)$ більше за значення функції~$cos(x)$.

\vspace{2mm}

%enditem

%beginitem
2. Написати функцію, що знаходить суму квадратів парних елементів послідовності цілих чисел.

\vspace{2mm}

%enditem

%beginitem
3. Написати функцію, що знаходить кількість входжень цифри 0 у С-рядок.

\vspace{2mm}

%enditem

%beginitem
4. 
Написати функцію, що записує в масив дійсних послідовність чисел
$$\scalebox{1.2}{$C_{n+k}^{2k-1}$}, \quad k=1,\ldots,n$$
Використовувати вкладені цикли та виклики функцій \textbf{заборонено}.

\vspace{2mm}

%enditem

%beginitem
5. Написати функцію, що для послідовності цілих знаходить найдовший відрізок, в якому всі сусідні числа відрізняються якнайменш на 3. Якщо таких відрізків кілька, то повернути перший (як індекс початку та довжину).

\vspace{2mm}

%enditem
\end{large}
%end
\newpage
%begin
%\begin {Huge}
{{27.11.2024 Контрольна робота \No 2, варіант  \No \textbf { 3 }\par}}

%\begin{center}Частина 2
%\end{center}
%\vspace{10mm}
%\end {Huge}

{
Якість коду та економність обчислень оцінюється по всім завданням. Послідовності задаються масивами. Використовувати додаткові структури даних (у тому числі рядки) та функції стандартної бібліотеки заборонено.

\vspace{2mm}

\begin {large}
%beginitem
1. Написати функцію, що знаходить у скількох цілих точках інтервалу $[a; b)$ значення функції~$sin(x)$ більше за значення функції~$cos(x)$.

\vspace{2mm}

%enditem

%beginitem
2. Написати функцію, що знаходить суму елементів послідовності дійсних чисел, що менші за свої індекси у масиві.

\vspace{2mm}

%enditem

%beginitem
3. Написати функцію, що перевіряє, чи містить С-рядок хоч одну малу латинську літеру.

\vspace{2mm}

%enditem

%beginitem
4. 
Написати функцію, що записує в масив дійсних послідовність чисел
$$\scalebox{1.2}{$C_{2k+1}^{k}$}, \quad k=2,\ldots,n$$
Використовувати вкладені цикли та виклики функцій \textbf{заборонено}.

\vspace{2mm}

%enditem

%beginitem
5. Написати функцію, що знаходить найбільше число з послідовності цілих, у десятковому запису якого нема цифри 3. Повертає індекс останнього входження цього числа або $-1$ (за відсутності).

\vspace{2mm}

%enditem
\end{large}
%end
\vspace{35mm}


%begin
%\begin {Huge}
{{27.11.2024 Контрольна робота \No 2, варіант  \No \textbf { 4 }\par}}

%\begin{center}Частина 2
%\end{center}
%\vspace{10mm}
%\end {Huge}

{
Якість коду та економність обчислень оцінюється по всім завданням. Послідовності задаються масивами. Використовувати додаткові структури даних (у тому числі рядки) та функції стандартної бібліотеки заборонено.

\vspace{2mm}

\begin {large}
%beginitem
1. Написати функцію, що знаходить найбільше значення функції~$sin(x)$ у цілих точках інтервалу $(a; b]$.

\vspace{2mm}

%enditem

%beginitem
2. Написати функцію, що збільшує всі парні елементи послідовності цілих чисел на 2.

\vspace{2mm}

%enditem

%beginitem
3. Написати функцію, що знаходить кількість входжень малих латинських літер у С-рядок.

\vspace{2mm}

%enditem

%beginitem
4. 
Написати функцію, що записує в масив дійсних послідовність чисел
$$\scalebox{1.2}{$C_{n+k}^{2k}$}, \quad k=1,\ldots,n$$
Використовувати вкладені цикли та виклики функцій \textbf{заборонено}.

\vspace{2mm}

%enditem

%beginitem
5. Написати функцію, що для послідовності цілих находить найдовший відрізок, в якому всі сусідні числа різні. Якщо таких відрізків кілька, то повернути останній (як індекс початку та довжину).

\vspace{2mm}

%enditem
\end{large}
%end
\newpage
%begin
%\begin {Huge}
{{27.11.2024 Контрольна робота \No 2, варіант  \No \textbf { 5 }\par}}

%\begin{center}Частина 2
%\end{center}
%\vspace{10mm}
%\end {Huge}

{
Якість коду та економність обчислень оцінюється по всім завданням. Послідовності задаються масивами. Використовувати додаткові структури даних (у тому числі рядки) та функції стандартної бібліотеки заборонено.

\vspace{2mm}

\begin {large}
%beginitem
1. Написати функцію, що знаходить у скількох цілих точках інтервалу $(a; b]$ функція~$sin(x)$ приймає від'ємне значення.

\vspace{2mm}

%enditem

%beginitem
2. Написати функцію, що замінює всі парні елементи послідовності цілих чисел числом 13.

\vspace{2mm}

%enditem

%beginitem
3. Написати функцію, що перевіряє, чи містить С-рядок онакову кількість входжень символів ( та ).

\vspace{2mm}

%enditem

%beginitem
4. 
Написати функцію, що записує в масив дійсних послідовність чисел
$$\scalebox{1.2}{$\frac{C_{2k}^k}{k+1}$}, \quad k=0,1,\ldots,n$$
Використовувати вкладені цикли та виклики функцій \textbf{заборонено}.

\vspace{2mm}

%enditem

%beginitem
5. Написати функцію, що для послідовності цілих знаходить найдовший відрізок, в якому всі сусідні числа різні. Якщо таких відрізків кілька, то повернути перший (як індекс початку та довжину).

\vspace{2mm}

%enditem
\end{large}
%end
\vspace{35mm}


%begin
%\begin {Huge}
{{27.11.2024 Контрольна робота \No 2, варіант  \No \textbf { 6 }\par}}

%\begin{center}Частина 2
%\end{center}
%\vspace{10mm}
%\end {Huge}

{
Якість коду та економність обчислень оцінюється по всім завданням. Послідовності задаються масивами. Використовувати додаткові структури даних (у тому числі рядки) та функції стандартної бібліотеки заборонено.

\vspace{2mm}

\begin {large}
%beginitem
1. Написати функцію, що знаходить у скількох цілих точках інтервалу $[a; b)$ функція~$cos(x)$ приймає додатнє значення.

\vspace{2mm}

%enditem

%beginitem
2. Написати функцію, що знаходить суму елементів послідовності дійсних чисел, що більші за свої індекси у масиві.

\vspace{2mm}

%enditem

%beginitem
3. Написати функцію, що перевіряє, яких символів у С-рядку більше: 0 чи 1.

\vspace{2mm}

%enditem

%beginitem
4. 
Написати функцію, що записує в масив дійсних послідовність чисел
$$\scalebox{1.2}{$(-1)^{k}C_{2n+k+1}^{n+k}$}, \quad k=0,1,\ldots,n$$
Використовувати вкладені цикли та виклики функцій \textbf{заборонено}.

\vspace{2mm}

%enditem

%beginitem
5. Написати функцію, що для послідовності цілих знаходить довжину найдовшого відрізку, в якому всі сусідні числа відрізняються якнайменш на 2.

\vspace{2mm}

%enditem
\end{large}
%end
\newpage
%begin
%\begin {Huge}
{{27.11.2024 Контрольна робота \No 2, варіант  \No \textbf { 7 }\par}}

%\begin{center}Частина 2
%\end{center}
%\vspace{10mm}
%\end {Huge}

{
Якість коду та економність обчислень оцінюється по всім завданням. Послідовності задаються масивами. Використовувати додаткові структури даних (у тому числі рядки) та функції стандартної бібліотеки заборонено.

\vspace{2mm}

\begin {large}
%beginitem
1. Написати функцію, що знаходить у скількох цілих точках інтервалу $(a; b)$ функція~$sin(x)$ приймає додатнє значення.

\vspace{2mm}

%enditem

%beginitem
2. Написати функцію, що збільшує всі непарні елементи послідовності цілих чисел на 4.

\vspace{2mm}

%enditem

%beginitem
3. Написати функцію, що знаходить кількість входжень великих латинських літер у С-рядок.

\vspace{2mm}

%enditem

%beginitem
4. 
Написати функцію, що записує в масив дійсних послідовність чисел
$$\scalebox{1.2}{$C_{n+k}^{2k}$}, \quad k=0,1,\ldots,n$$
Використовувати вкладені цикли та виклики функцій \textbf{заборонено}.

\vspace{2mm}

%enditem

%beginitem
5. Написати функцію, що для інтервалу $[a, b)$ знаходить знаходить найбільший степінь числа 3, на який ділиться хоча б одне ціле число з цього інтервалу.

\vspace{2mm}

%enditem
\end{large}
%end
\vspace{35mm}


%begin
%\begin {Huge}
{{27.11.2024 Контрольна робота \No 2, варіант  \No \textbf { 8 }\par}}

%\begin{center}Частина 2
%\end{center}
%\vspace{10mm}
%\end {Huge}

{
Якість коду та економність обчислень оцінюється по всім завданням. Послідовності задаються масивами. Використовувати додаткові структури даних (у тому числі рядки) та функції стандартної бібліотеки заборонено.

\vspace{2mm}

\begin {large}
%beginitem
1. Написати функцію, що знаходить найбільше значення функції~$sin(x)$ у цілих точках інтервалу $[a; b)$.

\vspace{2mm}

%enditem

%beginitem
2. Написати функцію, що знаходить суму елементів послідовності цілих чисел, що більші за свої індекси у масиві.

\vspace{2mm}

%enditem

%beginitem
3. Написати функцію, що знаходить кількість входжень літери Z у С-рядок.

\vspace{2mm}

%enditem

%beginitem
4. 
Написати функцію, що записує в масив дійсних послідовність чисел
$$\scalebox{1.2}{$C_{n+k}^{k+1}$}, \quad k=1,\ldots,n$$
Використовувати вкладені цикли та виклики функцій \textbf{заборонено}.

\vspace{2mm}

%enditem

%beginitem
5. Написати функцію, що для цілого числа знаходить кількість простих дільників, які входять у його розклад у найбільшому степені. \\ Наприклад, для числа $3^5{\cdot}7^9{\cdot}11{\cdot}23^{9}$ таких дільників два – $23$ та $7$, тому функція має повернути $2$.

\vspace{2mm}

%enditem
\end{large}
%end
\newpage
%begin
%\begin {Huge}
{{27.11.2024 Контрольна робота \No 2, варіант  \No \textbf { 9 }\par}}

%\begin{center}Частина 2
%\end{center}
%\vspace{10mm}
%\end {Huge}

{
Якість коду та економність обчислень оцінюється по всім завданням. Послідовності задаються масивами. Використовувати додаткові структури даних (у тому числі рядки) та функції стандартної бібліотеки заборонено.

\vspace{2mm}

\begin {large}
%beginitem
1. Написати функцію, що знаходить найбільше значення функції~$cos(x)$ у цілих точках інтервалу $(a; b]$.

\vspace{2mm}

%enditem

%beginitem
2. Написати функцію, що знаходить найменший непарний елемент послідовності цілих чисел.

\vspace{2mm}

%enditem

%beginitem
3. Написати функцію, що перевіряє, чи містить С-рядок хоч один символ підкреслення.

\vspace{2mm}

%enditem

%beginitem
4. 
Написати функцію, що записує в масив дійсних послідовність чисел
$$\scalebox{1.2}{$C_{3k}^{k}$}, \quad k=2,\ldots,n$$
Використовувати вкладені цикли та виклики функцій \textbf{заборонено}.

\vspace{2mm}

%enditem

%beginitem
5. Написати функцію, що для інтервалу $[a, b)$ знаходить найближчу пару чисел з цього інтервалу, що кожне з них має непарну кількість різних простих дільників. Якщо таких пар кілька, то цікавить пара з найменшим добутком.

\vspace{2mm}

%enditem
\end{large}
%end
\vspace{35mm}


%begin
%\begin {Huge}
{{27.11.2024 Контрольна робота \No 2, варіант  \No \textbf { 10 }\par}}

%\begin{center}Частина 2
%\end{center}
%\vspace{10mm}
%\end {Huge}

{
Якість коду та економність обчислень оцінюється по всім завданням. Послідовності задаються масивами. Використовувати додаткові структури даних (у тому числі рядки) та функції стандартної бібліотеки заборонено.

\vspace{2mm}

\begin {large}
%beginitem
1. Написати функцію, що знаходить у скількох цілих точках інтервалу $[a; b]$ функція~$cos(x)$ приймає додатнє значення.

\vspace{2mm}

%enditem

%beginitem
2. Написати функцію, що знаходить найбільший непарний елемент послідовності цілих чисел.

\vspace{2mm}

%enditem

%beginitem
3. Написати функцію, що перевіряє, чи містить С-рядок хоч один символ підкреслення.

\vspace{2mm}

%enditem

%beginitem
4. 
Написати функцію, що записує в масив дійсних послідовність чисел
$$\scalebox{1.2}{$(-1)^{k}C_{n+k}^{k+1}$}, \quad k=1,\ldots,n$$
Використовувати вкладені цикли та виклики функцій \textbf{заборонено}.

\vspace{2mm}

%enditem

%beginitem
5. Написати функцію, що для послідовності цілих знаходить довжину найдовшого відрізку, в якому всі сусідні числа різні.

\vspace{2mm}

%enditem
\end{large}
%end
\newpage
%begin
%\begin {Huge}
{{27.11.2024 Контрольна робота \No 2, варіант  \No \textbf { 11 }\par}}

%\begin{center}Частина 2
%\end{center}
%\vspace{10mm}
%\end {Huge}

{
Якість коду та економність обчислень оцінюється по всім завданням. Послідовності задаються масивами. Використовувати додаткові структури даних (у тому числі рядки) та функції стандартної бібліотеки заборонено.

\vspace{2mm}

\begin {large}
%beginitem
1. Написати функцію, що знаходить у скількох цілих точках інтервалу $[a; b]$ значення функції~$sin(x)$ менше за значення функції~$cos(x)$.

\vspace{2mm}

%enditem

%beginitem
2. Написати функцію, що замінює всі від'ємні елементи послідовності цілих чисел нулями.

\vspace{2mm}

%enditem

%beginitem
3. Написати функцію, що перевіряє, яких символів у С-рядку більше: 0 чи 1.

\vspace{2mm}

%enditem

%beginitem
4. 
Написати функцію, що записує в масив дійсних послідовність чисел
$$\scalebox{1.2}{$C_{n+2k}^{k+1}$}, \quad k=1,\ldots,n$$
Використовувати вкладені цикли та виклики функцій \textbf{заборонено}.

\vspace{2mm}

%enditem

%beginitem
5. Написати функцію, яка для цілого числа знаходить кількість його різних простих дільників, що мають у десятковому запису цифру 1.

\vspace{2mm}

%enditem
\end{large}
%end
\vspace{35mm}


%begin
%\begin {Huge}
{{27.11.2024 Контрольна робота \No 2, варіант  \No \textbf { 12 }\par}}

%\begin{center}Частина 2
%\end{center}
%\vspace{10mm}
%\end {Huge}

{
Якість коду та економність обчислень оцінюється по всім завданням. Послідовності задаються масивами. Використовувати додаткові структури даних (у тому числі рядки) та функції стандартної бібліотеки заборонено.

\vspace{2mm}

\begin {large}
%beginitem
1. Написати функцію, що знаходить у скількох цілих точках інтервалу $(a; b)$ функція~$sin(x)$ приймає від'ємне значення.

\vspace{2mm}

%enditem

%beginitem
2. Написати функцію, що знаходить суму квадратів елементів послідовності дійсних чисел, що менші 2.

\vspace{2mm}

%enditem

%beginitem
3. Написати функцію, що перевіряє, чи містить С-рядок хоч одну малу латинську літеру.

\vspace{2mm}

%enditem

%beginitem
4. 
Написати функцію, що записує в масив дійсних послідовність чисел
$$\scalebox{1.2}{$C_{2n+k}^{n+k}$}, \quad k=1,\ldots,n$$
Використовувати вкладені цикли та виклики функцій \textbf{заборонено}.

\vspace{2mm}

%enditem

%beginitem
5. Написати функцію, що для інтервалу $[a, b)$ знаходить найближчу пару чисел з цього інтервалу, що мають від трьох до п'яти (включно) простих дільників.

\vspace{2mm}

%enditem
\end{large}
%end
\newpage
%begin
%\begin {Huge}
{{27.11.2024 Контрольна робота \No 2, варіант  \No \textbf { 13 }\par}}

%\begin{center}Частина 2
%\end{center}
%\vspace{10mm}
%\end {Huge}

{
Якість коду та економність обчислень оцінюється по всім завданням. Послідовності задаються масивами. Використовувати додаткові структури даних (у тому числі рядки) та функції стандартної бібліотеки заборонено.

\vspace{2mm}

\begin {large}
%beginitem
1. Написати функцію, що знаходить у скількох цілих точках інтервалу $[a; b)$ функція~$sin(x)$ приймає від'ємне значення.

\vspace{2mm}

%enditem

%beginitem
2. Написати функцію, що замінює всі елементи на непарних позиціях послідовності цілих чисел (позиції нумеруються з 0) на число 25.

\vspace{2mm}

%enditem

%beginitem
3. Написати функцію, що знаходить кількість входжень літери Z у С-рядок.

\vspace{2mm}

%enditem

%beginitem
4. 
Написати функцію, що записує в масив дійсних послідовність чисел
$$\scalebox{1.2}{$(-1)^{k}C_{n+2k}^{k}$}, \quad k=0,1,\ldots,n-2$$
Використовувати вкладені цикли та виклики функцій \textbf{заборонено}.

\vspace{2mm}

%enditem

%beginitem
5. Написати функцію, що повертає число з послідовності цілих, що має парну кількість різних простих дільників. Якщо таких чисел кілька, то повернути найбільше.

\vspace{2mm}

%enditem
\end{large}
%end
\vspace{35mm}


%begin
%\begin {Huge}
{{27.11.2024 Контрольна робота \No 2, варіант  \No \textbf { 14 }\par}}

%\begin{center}Частина 2
%\end{center}
%\vspace{10mm}
%\end {Huge}

{
Якість коду та економність обчислень оцінюється по всім завданням. Послідовності задаються масивами. Використовувати додаткові структури даних (у тому числі рядки) та функції стандартної бібліотеки заборонено.

\vspace{2mm}

\begin {large}
%beginitem
1. Написати функцію, що знаходить у скількох цілих точках інтервалу $[a; b]$ функція~$cos(x)$ приймає від'ємне значення.

\vspace{2mm}

%enditem

%beginitem
2. Написати функцію, що замінює всі непарні елементи послідовності цілих чисел нулями.

\vspace{2mm}

%enditem

%beginitem
3. Написати функцію, що перевіряє, чи містить С-рядок хоч одну велику латинську літеру.

\vspace{2mm}

%enditem

%beginitem
4. 
Написати функцію, що записує в масив дійсних послідовність чисел
$$\scalebox{1.2}{$C_{n+k}^{k+1}$}, \quad k=0,1,\ldots,n-2$$
Використовувати вкладені цикли та виклики функцій \textbf{заборонено}.

\vspace{2mm}

%enditem

%beginitem
5. Написати функцію, що повертає число з послідовності цілих, що має найменшу кількість різних простих дільників. Якщо таких чисел кілька, то повернути найбільше.

\vspace{2mm}

%enditem
\end{large}
%end
\newpage
%begin
%\begin {Huge}
{{27.11.2024 Контрольна робота \No 2, варіант  \No \textbf { 15 }\par}}

%\begin{center}Частина 2
%\end{center}
%\vspace{10mm}
%\end {Huge}

{
Якість коду та економність обчислень оцінюється по всім завданням. Послідовності задаються масивами. Використовувати додаткові структури даних (у тому числі рядки) та функції стандартної бібліотеки заборонено.

\vspace{2mm}

\begin {large}
%beginitem
1. Написати функцію, що знаходить у скількох цілих точках інтервалу $[a; b]$ значення функції~$sin(x)$ більше за значення функції~$cos(x)$.

\vspace{2mm}

%enditem

%beginitem
2. Написати функцію, що знаходить суму елементів послідовності цілих чисел, що менші за свої індекси у масиві.

\vspace{2mm}

%enditem

%beginitem
3. Написати функцію, що знаходить кількість входжень цифри 0 у С-рядок.

\vspace{2mm}

%enditem

%beginitem
4. 
Написати функцію, що записує в масив дійсних послідовність чисел
$$\scalebox{1.2}{$(-1)^{k}C_{n+k}^{k+1}$}, \quad k=0,1,\ldots,n$$
Використовувати вкладені цикли та виклики функцій \textbf{заборонено}.

\vspace{2mm}

%enditem

%beginitem
5. Написати функцію, що для цілого числа знаходить кількість його простих дільників, що входять у його розклад рівно у другому степеню. \\ Наприклад, для числа $3^2{\cdot}7^9{\cdot}11{\cdot}23^{2}$ таких дільників два – $3$ та $23$, тому функція має повернути $2$.

\vspace{2mm}

%enditem
\end{large}
%end
\vspace{35mm}


%begin
%\begin {Huge}
{{27.11.2024 Контрольна робота \No 2, варіант  \No \textbf { 16 }\par}}

%\begin{center}Частина 2
%\end{center}
%\vspace{10mm}
%\end {Huge}

{
Якість коду та економність обчислень оцінюється по всім завданням. Послідовності задаються масивами. Використовувати додаткові структури даних (у тому числі рядки) та функції стандартної бібліотеки заборонено.

\vspace{2mm}

\begin {large}
%beginitem
1. Написати функцію, що знаходить найбільше значення функції~$cos(x)$ у цілих точках інтервалу $[a; b)$.

\vspace{2mm}

%enditem

%beginitem
2. Написати функцію, що замінює всі елементи на парних позиціях послідовності цілих чисел (позиції нумеруються з 0) на число 20.

\vspace{2mm}

%enditem

%beginitem
3. Написати функцію, що знаходить кількість входжень великих латинських літер у С-рядок.

\vspace{2mm}

%enditem

%beginitem
4. 
Написати функцію, що записує в масив дійсних послідовність чисел
$$\scalebox{1.2}{$(-1)^{k}C_{n+k}^{2k-1}$}, \quad k=1,\ldots,n$$
Використовувати вкладені цикли та виклики функцій \textbf{заборонено}.

\vspace{2mm}

%enditem

%beginitem
5. Написати функцію, що повертає найменше число з послідовності цілих, що має не більше трьох різних простих дільників. За відсутності має повертати $-1$.

\vspace{2mm}

%enditem
\end{large}
%end
\newpage
%begin
%\begin {Huge}
{{27.11.2024 Контрольна робота \No 2, варіант  \No \textbf { 17 }\par}}

%\begin{center}Частина 2
%\end{center}
%\vspace{10mm}
%\end {Huge}

{
Якість коду та економність обчислень оцінюється по всім завданням. Послідовності задаються масивами. Використовувати додаткові структури даних (у тому числі рядки) та функції стандартної бібліотеки заборонено.

\vspace{2mm}

\begin {large}
%beginitem
1. Написати функцію, що знаходить у скількох цілих точках інтервалу $(a; b]$ функція~$sin(x)$ приймає додатнє значення.

\vspace{2mm}

%enditem

%beginitem
2. Написати функцію, що знаходить середнє арифметичне елементів послідовності цілих чисел.

\vspace{2mm}

%enditem

%beginitem
3. Написати функцію, що знаходить кількість входжень малих латинських літер у С-рядок.

\vspace{2mm}

%enditem

%beginitem
4. 
Написати функцію, що записує в масив дійсних послідовність чисел
$$\scalebox{1.2}{$(-1)^{k}C_{n+2k}^{k+1}$}, \quad k=0,1,\ldots,n$$
Використовувати вкладені цикли та виклики функцій \textbf{заборонено}.

\vspace{2mm}

%enditem

%beginitem
5. Написати функцію, що для інтервалу $[a, b)$ знаходить найближчу пару чисел з цього інтервалу, що кожне з них має парну кількість різних простих дільників. Якщо таких пар кілька, то цікавить пара з найбільшою сумою.

\vspace{2mm}

%enditem
\end{large}
%end
\vspace{35mm}


%begin
%\begin {Huge}
{{27.11.2024 Контрольна робота \No 2, варіант  \No \textbf { 18 }\par}}

%\begin{center}Частина 2
%\end{center}
%\vspace{10mm}
%\end {Huge}

{
Якість коду та економність обчислень оцінюється по всім завданням. Послідовності задаються масивами. Використовувати додаткові структури даних (у тому числі рядки) та функції стандартної бібліотеки заборонено.

\vspace{2mm}

\begin {large}
%beginitem
1. Написати функцію, що знаходить у скількох цілих точках інтервалу $(a; b)$ функція~$cos(x)$ приймає додатнє значення.

\vspace{2mm}

%enditem

%beginitem
2. Написати функцію, що знаходить суму елементів послідовності цілих чисел, що дорівнюють своїм індексам у масиві.

\vspace{2mm}

%enditem

%beginitem
3. Написати функцію, що перевіряє, чи містить С-рядок онакову кількість входжень символів ( та ).

\vspace{2mm}

%enditem

%beginitem
4. 
Написати функцію, що записує в масив дійсних послідовність чисел
$$\scalebox{1.2}{$(-1)^{k}C_{2k-1}^{k}$}, \quad k=2,\ldots,n$$
Використовувати вкладені цикли та виклики функцій \textbf{заборонено}.

\vspace{2mm}

%enditem

%beginitem
5. Написати функцію, що для цілого числа знаходить простий дільник, який входить у його розклад ц найбільшому степені. Якщо таких кілька, то повернути найбільший. \\ Наприклад, для числа $2^{9}{\cdot}3^5{\cdot}7^9{\cdot}11$ функція має повернути $7$.

\vspace{2mm}

%enditem
\end{large}
%end
\newpage
%begin
%\begin {Huge}
{{27.11.2024 Контрольна робота \No 2, варіант  \No \textbf { 19 }\par}}

%\begin{center}Частина 2
%\end{center}
%\vspace{10mm}
%\end {Huge}

{
Якість коду та економність обчислень оцінюється по всім завданням. Послідовності задаються масивами. Використовувати додаткові структури даних (у тому числі рядки) та функції стандартної бібліотеки заборонено.

\vspace{2mm}

\begin {large}
%beginitem
1. Написати функцію, що знаходить у скількох цілих точках інтервалу $[a; b]$ функція~$sin(x)$ приймає від'ємне значення.

\vspace{2mm}

%enditem

%beginitem
2. Написати функцію, що зменшує всі непарні елементи послідовності цілих чисел на 1.

\vspace{2mm}

%enditem

%beginitem
3. Написати функцію, що перевіряє, чи містить С-рядок хоч один символ підкреслення.

\vspace{2mm}

%enditem

%beginitem
4. 
Написати функцію, що записує в масив дійсних послідовність чисел
$$\scalebox{1.2}{$(-1)^{k}C_{2n+k}^{n+k}$}, \quad k=0,1,\ldots,n$$
Використовувати вкладені цикли та виклики функцій \textbf{заборонено}.

\vspace{2mm}

%enditem

%beginitem
5. Написати функцію, що повертає число з послідовності цілих, що має найбільшу кількість різних простих дільників. Якщо таких чисел кілька, то повернути найменше.

\vspace{2mm}

%enditem
\end{large}
%end
\vspace{35mm}


%begin
%\begin {Huge}
{{27.11.2024 Контрольна робота \No 2, варіант  \No \textbf { 20 }\par}}

%\begin{center}Частина 2
%\end{center}
%\vspace{10mm}
%\end {Huge}

{
Якість коду та економність обчислень оцінюється по всім завданням. Послідовності задаються масивами. Використовувати додаткові структури даних (у тому числі рядки) та функції стандартної бібліотеки заборонено.

\vspace{2mm}

\begin {large}
%beginitem
1. Написати функцію, що знаходить у скількох цілих точках інтервалу $[a; b]$ функція~$sin(x)$ приймає додатнє значення.

\vspace{2mm}

%enditem

%beginitem
2. Написати функцію, що зменшує всі парні елементи послідовності цілих чисел на 3.

\vspace{2mm}

%enditem

%beginitem
3. Написати функцію, що знаходить кількість входжень великих латинських літер у С-рядок.

\vspace{2mm}

%enditem

%beginitem
4. 
Написати функцію, що записує в масив дійсних послідовність чисел
$$\scalebox{1.2}{$C_{2k-1}^{k}$}, \quad k=2,\ldots,n$$
Використовувати вкладені цикли та виклики функцій \textbf{заборонено}.

\vspace{2mm}

%enditem

%beginitem
5. Написати функцію, що в інтервалі $[a, b)$ знаходить ціле число, що ділиться на найбільший степінь числа 3. Якщо таких чисел декілька, повернути найбільше.

\vspace{2mm}

%enditem
\end{large}
%end
\newpage
%begin
%\begin {Huge}
{{27.11.2024 Контрольна робота \No 2, варіант  \No \textbf { 21 }\par}}

%\begin{center}Частина 2
%\end{center}
%\vspace{10mm}
%\end {Huge}

{
Якість коду та економність обчислень оцінюється по всім завданням. Послідовності задаються масивами. Використовувати додаткові структури даних (у тому числі рядки) та функції стандартної бібліотеки заборонено.

\vspace{2mm}

\begin {large}
%beginitem
1. Написати функцію, що знаходить найбільше значення функції~$cos(x)$ у цілих точках інтервалу $[a; b]$.

\vspace{2mm}

%enditem

%beginitem
2. Написати функцію, що замінює всі елементи на парних позиціях послідовності цілих чисел (позиції нумеруються з 0) на число 24.

\vspace{2mm}

%enditem

%beginitem
3. Написати функцію, що перевіряє, чи містить С-рядок хоч одну малу латинську літеру.

\vspace{2mm}

%enditem

%beginitem
4. 
Написати функцію, що записує в масив дійсних послідовність чисел
$$\scalebox{1.2}{$C_{n+k}^{k}$}, \quad k=1,\ldots,n$$
Використовувати вкладені цикли та виклики функцій \textbf{заборонено}.

\vspace{2mm}

%enditem

%beginitem
5. Написати функцію, що для послідовності цілих знаходить довжину найдовшого відрізку, що складається з чисел, у сімковому запису якого нема цифри 2. За відсутності має повертати $0$.

\vspace{2mm}

%enditem
\end{large}
%end
\vspace{35mm}


%begin
%\begin {Huge}
{{27.11.2024 Контрольна робота \No 2, варіант  \No \textbf { 22 }\par}}

%\begin{center}Частина 2
%\end{center}
%\vspace{10mm}
%\end {Huge}

{
Якість коду та економність обчислень оцінюється по всім завданням. Послідовності задаються масивами. Використовувати додаткові структури даних (у тому числі рядки) та функції стандартної бібліотеки заборонено.

\vspace{2mm}

\begin {large}
%beginitem
1. Написати функцію, що знаходить найбільше значення функції~$cos(x)$ у цілих точках інтервалу $(a; b)$.

\vspace{2mm}

%enditem

%beginitem
2. Написати функцію, що знаходить найбільший парний елемент послідовності цілих чисел.

\vspace{2mm}

%enditem

%beginitem
3. Написати функцію, що знаходить кількість входжень літери Z у С-рядок.

\vspace{2mm}

%enditem

%beginitem
4. 
Написати функцію, що записує в масив дійсних послідовність чисел
$$\scalebox{1.2}{$C_{2n+k}^{n+k}$}, \quad k=0,1,\ldots,n$$
Використовувати вкладені цикли та виклики функцій \textbf{заборонено}.

\vspace{2mm}

%enditem

%beginitem
5. Написати функцію, що знаходить найменше число з послідовності цілих, у десятковому запису якого нема цифри 7. Повертає індекс першого входження цього числа або $-1$ (за відсутності).

\vspace{2mm}

%enditem
\end{large}
%end
\newpage
%begin
%\begin {Huge}
{{27.11.2024 Контрольна робота \No 2, варіант  \No \textbf { 23 }\par}}

%\begin{center}Частина 2
%\end{center}
%\vspace{10mm}
%\end {Huge}

{
Якість коду та економність обчислень оцінюється по всім завданням. Послідовності задаються масивами. Використовувати додаткові структури даних (у тому числі рядки) та функції стандартної бібліотеки заборонено.

\vspace{2mm}

\begin {large}
%beginitem
1. Написати функцію, що знаходить у скількох цілих точках інтервалу $[a; b)$ значення функції~$sin(x)$ менше за значення функції~$cos(x)$.

\vspace{2mm}

%enditem

%beginitem
2. Написати функцію, що замінює всі непарні елементи послідовності цілих чисел числом 14.

\vspace{2mm}

%enditem

%beginitem
3. Написати функцію, що перевіряє, яких символів у С-рядку більше: 0 чи 1.

\vspace{2mm}

%enditem

%beginitem
4. 
Написати функцію, що записує в масив дійсних послідовність чисел
$$\scalebox{1.2}{$C_{n+2k}^{k}$}, \quad k=0,1,\ldots,n-2$$
Використовувати вкладені цикли та виклики функцій \textbf{заборонено}.

\vspace{2mm}

%enditem

%beginitem
5. Написати функцію, що для інтервалу $[a, b)$ знаходить найближчу пару чисел з цього інтервалу, що кожне з них має непарну кількість різних простих дільників. Якщо таких пар кілька, то цікавить пара з найбільшою сумою.

\vspace{2mm}

%enditem
\end{large}
%end
\vspace{35mm}


%begin
%\begin {Huge}
{{27.11.2024 Контрольна робота \No 2, варіант  \No \textbf { 24 }\par}}

%\begin{center}Частина 2
%\end{center}
%\vspace{10mm}
%\end {Huge}

{
Якість коду та економність обчислень оцінюється по всім завданням. Послідовності задаються масивами. Використовувати додаткові структури даних (у тому числі рядки) та функції стандартної бібліотеки заборонено.

\vspace{2mm}

\begin {large}
%beginitem
1. Написати функцію, що знаходить у скількох цілих точках інтервалу $[a; b)$ функція~$sin(x)$ приймає додатнє значення.

\vspace{2mm}

%enditem

%beginitem
2. Написати функцію, що замінює всі парні елементи послідовності цілих чисел нулями.

\vspace{2mm}

%enditem

%beginitem
3. Написати функцію, що знаходить кількість входжень малих латинських літер у С-рядок.

\vspace{2mm}

%enditem

%beginitem
4. 
Написати функцію, що записує в масив дійсних послідовність чисел
$$\scalebox{1.2}{$(-1)^{k}C_{3k}^{k}$}, \quad k=2,\ldots,n$$
Використовувати вкладені цикли та виклики функцій \textbf{заборонено}.

\vspace{2mm}

%enditem

%beginitem
5. Написати функцію, що знаходить найменше число з послідовності цілих, у десятковому запису якого нема цифри 5. Повертає індекс останнього входження цього числа або $-1$ (за відсутності).

\vspace{2mm}

%enditem
\end{large}
%end
\newpage
%begin
%\begin {Huge}
{{27.11.2024 Контрольна робота \No 2, варіант  \No \textbf { 25 }\par}}

%\begin{center}Частина 2
%\end{center}
%\vspace{10mm}
%\end {Huge}

{
Якість коду та економність обчислень оцінюється по всім завданням. Послідовності задаються масивами. Використовувати додаткові структури даних (у тому числі рядки) та функції стандартної бібліотеки заборонено.

\vspace{2mm}

\begin {large}
%beginitem
1. Написати функцію, що знаходить найбільше значення функції~$sin(x)$ у цілих точках інтервалу $(a; b)$.

\vspace{2mm}

%enditem

%beginitem
2. Написати функцію, що знаходить суму непарних елементів послідовності цілих чисел.

\vspace{2mm}

%enditem

%beginitem
3. Написати функцію, що перевіряє, чи містить С-рядок онакову кількість входжень символів ( та ).

\vspace{2mm}

%enditem

%beginitem
4. 
Написати функцію, що записує в масив дійсних послідовність чисел
$$\scalebox{1.2}{$(-1)^{k}C_{n+k}^{2k}$}, \quad k=1,\ldots,n$$
Використовувати вкладені цикли та виклики функцій \textbf{заборонено}.

\vspace{2mm}

%enditem

%beginitem
5. Написати функцію, що повертає найбільше число з послідовності цілих, що має не більше трьох різних простих дільників. За відсутності має повертати $-1$.

\vspace{2mm}

%enditem
\end{large}
%end
\vspace{35mm}


%begin
%\begin {Huge}
{{27.11.2024 Контрольна робота \No 2, варіант  \No \textbf { 26 }\par}}

%\begin{center}Частина 2
%\end{center}
%\vspace{10mm}
%\end {Huge}

{
Якість коду та економність обчислень оцінюється по всім завданням. Послідовності задаються масивами. Використовувати додаткові структури даних (у тому числі рядки) та функції стандартної бібліотеки заборонено.

\vspace{2mm}

\begin {large}
%beginitem
1. Написати функцію, що знаходить у скількох цілих точках інтервалу $(a; b]$ значення функції~$sin(x)$ менше за значення функції~$cos(x)$.

\vspace{2mm}

%enditem

%beginitem
2. Написати функцію, що знаходить найменший парний елемент послідовності цілих чисел.

\vspace{2mm}

%enditem

%beginitem
3. Написати функцію, що знаходить кількість входжень цифри 0 у С-рядок.

\vspace{2mm}

%enditem

%beginitem
4. 
Написати функцію, що записує в масив дійсних послідовність чисел
$$\scalebox{1.2}{$(-1)^{k}C_{n+k}^{k}$}, \quad k=1,\ldots,n-2$$
Використовувати вкладені цикли та виклики функцій \textbf{заборонено}.

\vspace{2mm}

%enditem

%beginitem
5. Написати функцію, що для цілого числа знаходить простий дільник, який входить у його розклад у найбільшому степені. Якщо таких кілька, то повернути найменший. \\ Наприклад, для числа $2^{9}{\cdot}3^5{\cdot}7^9{\cdot}11$ функція має повернути $2$.

\vspace{2mm}

%enditem
\end{large}
%end
\newpage
%begin
%\begin {Huge}
{{27.11.2024 Контрольна робота \No 2, варіант  \No \textbf { 27 }\par}}

%\begin{center}Частина 2
%\end{center}
%\vspace{10mm}
%\end {Huge}

{
Якість коду та економність обчислень оцінюється по всім завданням. Послідовності задаються масивами. Використовувати додаткові структури даних (у тому числі рядки) та функції стандартної бібліотеки заборонено.

\vspace{2mm}

\begin {large}
%beginitem
1. Написати функцію, що знаходить найбільше значення функції~$sin(x)$ у цілих точках інтервалу $[a; b]$.

\vspace{2mm}

%enditem

%beginitem
2. Написати функцію, що знаходить суму елементів послідовності дійсних чисел, що більші 1.

\vspace{2mm}

%enditem

%beginitem
3. Написати функцію, що перевіряє, чи містить С-рядок хоч одну велику латинську літеру.

\vspace{2mm}

%enditem

%beginitem
4. 
Написати функцію, що записує в масив дійсних послідовність чисел
$$\scalebox{1.2}{$C_{2n+k+1}^{n+k}$}, \quad k=1,\ldots,n$$
Використовувати вкладені цикли та виклики функцій \textbf{заборонено}.

\vspace{2mm}

%enditem

%beginitem
5. Написати функцію, що для цілого числа знаходить кількість його простих дільників, що входять у його розклад не більш ніж у 5му степеню. \\ Наприклад, для числа $3^8{\cdot}7^9{\cdot}11{\cdot}23^{2}$ таких дільників два – $11$ та $23$, тому функція має повернути $2$.

\vspace{2mm}

%enditem
\end{large}
%end
\vspace{35mm}


%begin
%\begin {Huge}
{{27.11.2024 Контрольна робота \No 2, варіант  \No \textbf { 28 }\par}}

%\begin{center}Частина 2
%\end{center}
%\vspace{10mm}
%\end {Huge}

{
Якість коду та економність обчислень оцінюється по всім завданням. Послідовності задаються масивами. Використовувати додаткові структури даних (у тому числі рядки) та функції стандартної бібліотеки заборонено.

\vspace{2mm}

\begin {large}
%beginitem
1. Написати функцію, що знаходить у скількох цілих точках інтервалу $(a; b)$ значення функції~$sin(x)$ менше за значення функції~$cos(x)$.

\vspace{2mm}

%enditem

%beginitem
2. Написати функцію, що знаходить середнє арифметичне елементів послідовності цілих чисел.

\vspace{2mm}

%enditem

%beginitem
3. Написати функцію, що перевіряє, чи містить С-рядок хоч одну малу латинську літеру.

\vspace{2mm}

%enditem

%beginitem
4. 
Написати функцію, що записує в масив дійсних послідовність чисел
$$\scalebox{1.2}{$(-1)^{k}C_{n+k}^{k}$}, \quad k=0,1,\ldots,n$$
Використовувати вкладені цикли та виклики функцій \textbf{заборонено}.

\vspace{2mm}

%enditem

%beginitem
5. Написати функцію, що для інтервалу $[a, b)$ знаходить кількість чисел, квадрати яких діляться на суму квадратів їх цифр. (Розглядаємо десятковий запис.)

\vspace{2mm}

%enditem
\end{large}
%end
\newpage
%begin
%\begin {Huge}
{{27.11.2024 Контрольна робота \No 2, варіант  \No \textbf { 29 }\par}}

%\begin{center}Частина 2
%\end{center}
%\vspace{10mm}
%\end {Huge}

{
Якість коду та економність обчислень оцінюється по всім завданням. Послідовності задаються масивами. Використовувати додаткові структури даних (у тому числі рядки) та функції стандартної бібліотеки заборонено.

\vspace{2mm}

\begin {large}
%beginitem
1. Написати функцію, що знаходить у скількох цілих точках інтервалу $(a; b]$ значення функції~$sin(x)$ більше за значення функції~$cos(x)$.

\vspace{2mm}

%enditem

%beginitem
2. Написати функцію, що зменшує всі непарні елементи послідовності цілих чисел на 1.

\vspace{2mm}

%enditem

%beginitem
3. Написати функцію, що знаходить кількість входжень великих латинських літер у С-рядок.

\vspace{2mm}

%enditem

%beginitem
4. 
Написати функцію, що записує в масив дійсних послідовність чисел
$$\scalebox{1.2}{$(-1)^{k}\frac{C_{2k}^k}{k+1}$}, \quad k=0,1,\ldots,n-2$$
Використовувати вкладені цикли та виклики функцій \textbf{заборонено}.

\vspace{2mm}

%enditem

%beginitem
5. Написати функцію, що для інтервалу $[a, b)$ знаходить найближчу пару чисел з цього інтервалу, що мають не менше трьох простих дільників та мають цифру 3 у своєму десятковому запису.

\vspace{2mm}

%enditem
\end{large}
%end
\vspace{35mm}


%begin
%\begin {Huge}
{{27.11.2024 Контрольна робота \No 2, варіант  \No \textbf { 30 }\par}}

%\begin{center}Частина 2
%\end{center}
%\vspace{10mm}
%\end {Huge}

{
Якість коду та економність обчислень оцінюється по всім завданням. Послідовності задаються масивами. Використовувати додаткові структури даних (у тому числі рядки) та функції стандартної бібліотеки заборонено.

\vspace{2mm}

\begin {large}
%beginitem
1. Написати функцію, що знаходить у скількох цілих точках інтервалу $[a; b)$ функція~$cos(x)$ приймає від'ємне значення.

\vspace{2mm}

%enditem

%beginitem
2. Написати функцію, що знаходить суму елементів послідовності дійсних чисел, що більші 1.

\vspace{2mm}

%enditem

%beginitem
3. Написати функцію, що перевіряє, чи містить С-рядок хоч один символ підкреслення.

\vspace{2mm}

%enditem

%beginitem
4. 
Написати функцію, що записує в масив дійсних послідовність чисел
$$\scalebox{1.2}{$(-1)^{k}C_{3k-1}^{k}$}, \quad k=2,\ldots,n$$
Використовувати вкладені цикли та виклики функцій \textbf{заборонено}.

\vspace{2mm}

%enditem

%beginitem
5. Написати функцію, що повертає найбільше число з послідовності цілих, у сімковому запису якого нема цифри 2. За відсутності має повертати $-1$.

\vspace{2mm}

%enditem
\end{large}
%end
\newpage
%begin
%\begin {Huge}
{{27.11.2024 Контрольна робота \No 2, варіант  \No \textbf { 31 }\par}}

%\begin{center}Частина 2
%\end{center}
%\vspace{10mm}
%\end {Huge}

{
Якість коду та економність обчислень оцінюється по всім завданням. Послідовності задаються масивами. Використовувати додаткові структури даних (у тому числі рядки) та функції стандартної бібліотеки заборонено.

\vspace{2mm}

\begin {large}
%beginitem
1. Написати функцію, що знаходить у скількох цілих точках інтервалу $(a; b)$ функція~$cos(x)$ приймає від'ємне значення.

\vspace{2mm}

%enditem

%beginitem
2. Написати функцію, що знаходить суму квадратів елементів послідовності дійсних чисел, що менші 2.

\vspace{2mm}

%enditem

%beginitem
3. Написати функцію, що знаходить кількість входжень літери Z у С-рядок.

\vspace{2mm}

%enditem

%beginitem
4. 
Написати функцію, що записує в масив дійсних послідовність чисел
$$\scalebox{1.2}{$C_{3k-1}^{k}$}, \quad k=2,\ldots,n$$
Використовувати вкладені цикли та виклики функцій \textbf{заборонено}.

\vspace{2mm}

%enditem

%beginitem
5. Написати функцію, що повертає найбільше число з послідовності цілих, що має рівно три різних простих дільники. За відсутності має повертати $-1$.

\vspace{2mm}

%enditem
\end{large}
%end
\vspace{35mm}


%begin
%\begin {Huge}
{{27.11.2024 Контрольна робота \No 2, варіант  \No \textbf { 32 }\par}}

%\begin{center}Частина 2
%\end{center}
%\vspace{10mm}
%\end {Huge}

{
Якість коду та економність обчислень оцінюється по всім завданням. Послідовності задаються масивами. Використовувати додаткові структури даних (у тому числі рядки) та функції стандартної бібліотеки заборонено.

\vspace{2mm}

\begin {large}
%beginitem
1. Написати функцію, що знаходить у скількох цілих точках інтервалу $(a; b]$ функція~$cos(x)$ приймає від'ємне значення.

\vspace{2mm}

%enditem

%beginitem
2. Написати функцію, що знаходить суму елементів послідовності цілих чисел, що менші за свої індекси у масиві.

\vspace{2mm}

%enditem

%beginitem
3. Написати функцію, що перевіряє, чи містить С-рядок онакову кількість входжень символів ( та ).

\vspace{2mm}

%enditem

%beginitem
4. 
Написати функцію, що записує в масив дійсних послідовність чисел
$$\scalebox{1.2}{$(-1)^{k}C_{n+2k}^{k+1}$}, \quad k=1,\ldots,n$$
Використовувати вкладені цикли та виклики функцій \textbf{заборонено}.

\vspace{2mm}

%enditem

%beginitem
5. Написати функцію, що для інтервалу $[a, b)$ знаходить найближчу пару чисел з цього інтервалу, що кожне з них має парну кількість різних простих дільників. Якщо таких пар кілька, то цікавить пара з найменшим добутком.

\vspace{2mm}

%enditem
\end{large}
%end
\newpage
%begin
%\begin {Huge}
{{27.11.2024 Контрольна робота \No 2, варіант  \No \textbf { 33 }\par}}

%\begin{center}Частина 2
%\end{center}
%\vspace{10mm}
%\end {Huge}

{
Якість коду та економність обчислень оцінюється по всім завданням. Послідовності задаються масивами. Використовувати додаткові структури даних (у тому числі рядки) та функції стандартної бібліотеки заборонено.

\vspace{2mm}

\begin {large}
%beginitem
1. Написати функцію, що знаходить найбільше значення функції~$cos(x)$ у цілих точках інтервалу $(a; b]$.

\vspace{2mm}

%enditem

%beginitem
2. Написати функцію, що замінює всі елементи на непарних позиціях послідовності цілих чисел (позиції нумеруються з 0) на число 25.

\vspace{2mm}

%enditem

%beginitem
3. Написати функцію, що перевіряє, чи містить С-рядок хоч одну велику латинську літеру.

\vspace{2mm}

%enditem

%beginitem
4. 
Написати функцію, що записує в масив дійсних послідовність чисел
$$\scalebox{1.2}{$(-1)^{k}C_{2n+k}^{n+k}$}, \quad k=1,\ldots,n$$
Використовувати вкладені цикли та виклики функцій \textbf{заборонено}.

\vspace{2mm}

%enditem

%beginitem
5. Написати функцію, що знаходить найбільше число з послідовності цілих, у десятковому запису якого нема цифри 4. Повертає індекс першого входження цього числа або $-1$ (за відсутності).

\vspace{2mm}

%enditem
\end{large}
%end
\vspace{35mm}


%begin
%\begin {Huge}
{{27.11.2024 Контрольна робота \No 2, варіант  \No \textbf { 34 }\par}}

%\begin{center}Частина 2
%\end{center}
%\vspace{10mm}
%\end {Huge}

{
Якість коду та економність обчислень оцінюється по всім завданням. Послідовності задаються масивами. Використовувати додаткові структури даних (у тому числі рядки) та функції стандартної бібліотеки заборонено.

\vspace{2mm}

\begin {large}
%beginitem
1. Написати функцію, що знаходить у скількох цілих точках інтервалу $(a; b)$ значення функції~$sin(x)$ менше за значення функції~$cos(x)$.

\vspace{2mm}

%enditem

%beginitem
2. Написати функцію, що знаходить найменший непарний елемент послідовності цілих чисел.

\vspace{2mm}

%enditem

%beginitem
3. Написати функцію, що знаходить кількість входжень малих латинських літер у С-рядок.

\vspace{2mm}

%enditem

%beginitem
4. 
Написати функцію, що записує в масив дійсних послідовність чисел
$$\scalebox{1.2}{$C_{n+2k}^{k+1}$}, \quad k=0,1,\ldots,n$$
Використовувати вкладені цикли та виклики функцій \textbf{заборонено}.

\vspace{2mm}

%enditem

%beginitem
5. Написати функцію, що в інтервалі $[a, b)$ знаходить ціле число, що ділиться на найбільший степінь числа 3. Якщо таких чисел декілька, повернути найменше.

\vspace{2mm}

%enditem
\end{large}
%end
\newpage
%begin
%\begin {Huge}
{{27.11.2024 Контрольна робота \No 2, варіант  \No \textbf { 35 }\par}}

%\begin{center}Частина 2
%\end{center}
%\vspace{10mm}
%\end {Huge}

{
Якість коду та економність обчислень оцінюється по всім завданням. Послідовності задаються масивами. Використовувати додаткові структури даних (у тому числі рядки) та функції стандартної бібліотеки заборонено.

\vspace{2mm}

\begin {large}
%beginitem
1. Написати функцію, що знаходить у скількох цілих точках інтервалу $(a; b]$ функція~$sin(x)$ приймає додатнє значення.

\vspace{2mm}

%enditem

%beginitem
2. Написати функцію, що знаходить суму елементів послідовності дійсних чисел, що менші за свої індекси у масиві.

\vspace{2mm}

%enditem

%beginitem
3. Написати функцію, що перевіряє, яких символів у С-рядку більше: 0 чи 1.

\vspace{2mm}

%enditem

%beginitem
4. 
Написати функцію, що записує в масив дійсних послідовність чисел
$$\scalebox{1.2}{$(-1)^{k}C_{2k+1}^{k}$}, \quad k=2,\ldots,n$$
Використовувати вкладені цикли та виклики функцій \textbf{заборонено}.

\vspace{2mm}

%enditem

%beginitem
5. Написати функцію, що знаходить найменше число з послідовності цілих, у десятковому запису якого нема цифри 7. Повертає індекс першого входження цього числа або $-1$ (за відсутності).

\vspace{2mm}

%enditem
\end{large}
%end
\vspace{35mm}


%begin
%\begin {Huge}
{{27.11.2024 Контрольна робота \No 2, варіант  \No \textbf { 36 }\par}}

%\begin{center}Частина 2
%\end{center}
%\vspace{10mm}
%\end {Huge}

{
Якість коду та економність обчислень оцінюється по всім завданням. Послідовності задаються масивами. Використовувати додаткові структури даних (у тому числі рядки) та функції стандартної бібліотеки заборонено.

\vspace{2mm}

\begin {large}
%beginitem
1. Написати функцію, що знаходить у скількох цілих точках інтервалу $[a; b]$ функція~$sin(x)$ приймає від'ємне значення.

\vspace{2mm}

%enditem

%beginitem
2. Написати функцію, що збільшує всі парні елементи послідовності цілих чисел на 2.

\vspace{2mm}

%enditem

%beginitem
3. Написати функцію, що знаходить кількість входжень цифри 0 у С-рядок.

\vspace{2mm}

%enditem

%beginitem
4. 
Написати функцію, що записує в масив дійсних послідовність чисел
$$\scalebox{1.2}{$(-1)^{k}C_{2n+k+1}^{n+k}$}, \quad k=1,\ldots,n$$
Використовувати вкладені цикли та виклики функцій \textbf{заборонено}.

\vspace{2mm}

%enditem

%beginitem
5. Написати функцію, що для інтервалу $[a, b)$ знаходить найближчу пару чисел з цього інтервалу, що мають від трьох до п'яти (включно) простих дільників.

\vspace{2mm}

%enditem
\end{large}
%end
\newpage
%begin
%\begin {Huge}
{{27.11.2024 Контрольна робота \No 2, варіант  \No \textbf { 37 }\par}}

%\begin{center}Частина 2
%\end{center}
%\vspace{10mm}
%\end {Huge}

{
Якість коду та економність обчислень оцінюється по всім завданням. Послідовності задаються масивами. Використовувати додаткові структури даних (у тому числі рядки) та функції стандартної бібліотеки заборонено.

\vspace{2mm}

\begin {large}
%beginitem
1. Написати функцію, що знаходить у скількох цілих точках інтервалу $[a; b)$ функція~$cos(x)$ приймає додатнє значення.

\vspace{2mm}

%enditem

%beginitem
2. Написати функцію, що знаходить суму квадратів парних елементів послідовності цілих чисел.

\vspace{2mm}

%enditem

%beginitem
3. Написати функцію, що перевіряє, чи містить С-рядок хоч одну велику латинську літеру.

\vspace{2mm}

%enditem

%beginitem
4. 
Написати функцію, що записує в масив дійсних послідовність чисел
$$\scalebox{1.2}{$(-1)^{k}C_{n+2k}^{k}$}, \quad k=1,\ldots,n$$
Використовувати вкладені цикли та виклики функцій \textbf{заборонено}.

\vspace{2mm}

%enditem

%beginitem
5. Написати функцію, що для послідовності цілих знаходить довжину найдовшого відрізку, в якому всі сусідні числа відрізняються якнайменш на 2.

\vspace{2mm}

%enditem
\end{large}
%end
\vspace{35mm}


%begin
%\begin {Huge}
{{27.11.2024 Контрольна робота \No 2, варіант  \No \textbf { 38 }\par}}

%\begin{center}Частина 2
%\end{center}
%\vspace{10mm}
%\end {Huge}

{
Якість коду та економність обчислень оцінюється по всім завданням. Послідовності задаються масивами. Використовувати додаткові структури даних (у тому числі рядки) та функції стандартної бібліотеки заборонено.

\vspace{2mm}

\begin {large}
%beginitem
1. Написати функцію, що знаходить у скількох цілих точках інтервалу $[a; b]$ значення функції~$sin(x)$ менше за значення функції~$cos(x)$.

\vspace{2mm}

%enditem

%beginitem
2. Написати функцію, що знаходить суму елементів послідовності цілих чисел, що більші за свої індекси у масиві.

\vspace{2mm}

%enditem

%beginitem
3. Написати функцію, що перевіряє, чи містить С-рядок хоч один символ підкреслення.

\vspace{2mm}

%enditem

%beginitem
4. 
Написати функцію, що записує в масив дійсних послідовність чисел
$$\scalebox{1.2}{$(-1)^{k}C_{n+k}^{2k}$}, \quad k=0,1,\ldots,n$$
Використовувати вкладені цикли та виклики функцій \textbf{заборонено}.

\vspace{2mm}

%enditem

%beginitem
5. Написати функцію, що в інтервалі $[a, b)$ знаходить ціле число, що ділиться на найбільший степінь числа 3. Якщо таких чисел декілька, повернути найбільше.

\vspace{2mm}

%enditem
\end{large}
%end
\newpage
%begin
%\begin {Huge}
{{27.11.2024 Контрольна робота \No 2, варіант  \No \textbf { 39 }\par}}

%\begin{center}Частина 2
%\end{center}
%\vspace{10mm}
%\end {Huge}

{
Якість коду та економність обчислень оцінюється по всім завданням. Послідовності задаються масивами. Використовувати додаткові структури даних (у тому числі рядки) та функції стандартної бібліотеки заборонено.

\vspace{2mm}

\begin {large}
%beginitem
1. Написати функцію, що знаходить у скількох цілих точках інтервалу $[a; b]$ функція~$cos(x)$ приймає додатнє значення.

\vspace{2mm}

%enditem

%beginitem
2. Написати функцію, що знаходить суму елементів послідовності дійсних чисел, що більші за свої індекси у масиві.

\vspace{2mm}

%enditem

%beginitem
3. Написати функцію, що знаходить кількість входжень цифри 0 у С-рядок.

\vspace{2mm}

%enditem

%beginitem
4. 
Написати функцію, що записує в масив дійсних послідовність чисел
$$\scalebox{1.2}{$C_{n+k}^{k}$}, \quad k=0,1,\ldots,n-2$$
Використовувати вкладені цикли та виклики функцій \textbf{заборонено}.

\vspace{2mm}

%enditem

%beginitem
5. Написати функцію, що для послідовності цілих знаходить довжину найдовшого відрізку, в якому всі сусідні числа різні.

\vspace{2mm}

%enditem
\end{large}
%end
\vspace{35mm}


%begin
%\begin {Huge}
{{27.11.2024 Контрольна робота \No 2, варіант  \No \textbf { 40 }\par}}

%\begin{center}Частина 2
%\end{center}
%\vspace{10mm}
%\end {Huge}

{
Якість коду та економність обчислень оцінюється по всім завданням. Послідовності задаються масивами. Використовувати додаткові структури даних (у тому числі рядки) та функції стандартної бібліотеки заборонено.

\vspace{2mm}

\begin {large}
%beginitem
1. Написати функцію, що знаходить у скількох цілих точках інтервалу $[a; b]$ функція~$cos(x)$ приймає від'ємне значення.

\vspace{2mm}

%enditem

%beginitem
2. Написати функцію, що замінює всі елементи на парних позиціях послідовності цілих чисел (позиції нумеруються з 0) на число 20.

\vspace{2mm}

%enditem

%beginitem
3. Написати функцію, що перевіряє, чи містить С-рядок онакову кількість входжень символів ( та ).

\vspace{2mm}

%enditem

%beginitem
4. 
Написати функцію, що записує в масив дійсних послідовність чисел
$$\scalebox{1.2}{$C_{n+2k}^{k}$}, \quad k=1,\ldots,n$$
Використовувати вкладені цикли та виклики функцій \textbf{заборонено}.

\vspace{2mm}

%enditem

%beginitem
5. Написати функцію, що повертає найбільше число з послідовності цілих, у сімковому запису якого нема цифри 2. За відсутності має повертати $-1$.

\vspace{2mm}

%enditem
\end{large}
%end
\newpage
%begin
%\begin {Huge}
{{27.11.2024 Контрольна робота \No 2, варіант  \No \textbf { 41 }\par}}

%\begin{center}Частина 2
%\end{center}
%\vspace{10mm}
%\end {Huge}

{
Якість коду та економність обчислень оцінюється по всім завданням. Послідовності задаються масивами. Використовувати додаткові структури даних (у тому числі рядки) та функції стандартної бібліотеки заборонено.

\vspace{2mm}

\begin {large}
%beginitem
1. Написати функцію, що знаходить у скількох цілих точках інтервалу $(a; b)$ функція~$cos(x)$ приймає від'ємне значення.

\vspace{2mm}

%enditem

%beginitem
2. Написати функцію, що знаходить найбільший парний елемент послідовності цілих чисел.

\vspace{2mm}

%enditem

%beginitem
3. Написати функцію, що перевіряє, чи містить С-рядок хоч одну малу латинську літеру.

\vspace{2mm}

%enditem

%beginitem
4. 
Написати функцію, що записує в масив дійсних послідовність чисел
$$\scalebox{1.2}{$C_{3k}^{k}$}, \quad k=2,\ldots,n$$
Використовувати вкладені цикли та виклики функцій \textbf{заборонено}.

\vspace{2mm}

%enditem

%beginitem
5. Написати функцію, що повертає число з послідовності цілих, що має найбільшу кількість різних простих дільників. Якщо таких чисел кілька, то повернути найменше.

\vspace{2mm}

%enditem
\end{large}
%end
\vspace{35mm}


%begin
%\begin {Huge}
{{27.11.2024 Контрольна робота \No 2, варіант  \No \textbf { 42 }\par}}

%\begin{center}Частина 2
%\end{center}
%\vspace{10mm}
%\end {Huge}

{
Якість коду та економність обчислень оцінюється по всім завданням. Послідовності задаються масивами. Використовувати додаткові структури даних (у тому числі рядки) та функції стандартної бібліотеки заборонено.

\vspace{2mm}

\begin {large}
%beginitem
1. Написати функцію, що знаходить у скількох цілих точках інтервалу $(a; b]$ функція~$sin(x)$ приймає від'ємне значення.

\vspace{2mm}

%enditem

%beginitem
2. Написати функцію, що замінює всі додатні елементи послідовності цілих чисел числом 15.

\vspace{2mm}

%enditem

%beginitem
3. Написати функцію, що знаходить кількість входжень літери Z у С-рядок.

\vspace{2mm}

%enditem

%beginitem
4. 
Написати функцію, що записує в масив дійсних послідовність чисел
$$\scalebox{1.2}{$(-1)^{k}C_{2n+k+1}^{n+k}$}, \quad k=0,1,\ldots,n$$
Використовувати вкладені цикли та виклики функцій \textbf{заборонено}.

\vspace{2mm}

%enditem

%beginitem
5. Написати функцію, що повертає число з послідовності цілих, що має парну кількість різних простих дільників. Якщо таких чисел кілька, то повернути найбільше.

\vspace{2mm}

%enditem
\end{large}
%end
\newpage
%begin
%\begin {Huge}
{{27.11.2024 Контрольна робота \No 2, варіант  \No \textbf { 43 }\par}}

%\begin{center}Частина 2
%\end{center}
%\vspace{10mm}
%\end {Huge}

{
Якість коду та економність обчислень оцінюється по всім завданням. Послідовності задаються масивами. Використовувати додаткові структури даних (у тому числі рядки) та функції стандартної бібліотеки заборонено.

\vspace{2mm}

\begin {large}
%beginitem
1. Написати функцію, що знаходить у скількох цілих точках інтервалу $(a; b)$ функція~$sin(x)$ приймає від'ємне значення.

\vspace{2mm}

%enditem

%beginitem
2. Написати функцію, що замінює всі парні елементи послідовності цілих чисел числом 13.

\vspace{2mm}

%enditem

%beginitem
3. Написати функцію, що перевіряє, яких символів у С-рядку більше: 0 чи 1.

\vspace{2mm}

%enditem

%beginitem
4. 
Написати функцію, що записує в масив дійсних послідовність чисел
$$\scalebox{1.2}{$C_{n+k}^{k}$}, \quad k=0,1,\ldots,n-2$$
Використовувати вкладені цикли та виклики функцій \textbf{заборонено}.

\vspace{2mm}

%enditem

%beginitem
5. Написати функцію, що для послідовності цілих знаходить найдовший відрізок, в якому всі сусідні числа відрізняються якнайменш на 3. Якщо таких відрізків кілька, то повернути перший (як індекс початку та довжину).

\vspace{2mm}

%enditem
\end{large}
%end
\vspace{35mm}


%begin
%\begin {Huge}
{{27.11.2024 Контрольна робота \No 2, варіант  \No \textbf { 44 }\par}}

%\begin{center}Частина 2
%\end{center}
%\vspace{10mm}
%\end {Huge}

{
Якість коду та економність обчислень оцінюється по всім завданням. Послідовності задаються масивами. Використовувати додаткові структури даних (у тому числі рядки) та функції стандартної бібліотеки заборонено.

\vspace{2mm}

\begin {large}
%beginitem
1. Написати функцію, що знаходить у скількох цілих точках інтервалу $[a; b)$ функція~$sin(x)$ приймає від'ємне значення.

\vspace{2mm}

%enditem

%beginitem
2. Написати функцію, що збільшує всі непарні елементи послідовності цілих чисел на 4.

\vspace{2mm}

%enditem

%beginitem
3. Написати функцію, що знаходить кількість входжень малих латинських літер у С-рядок.

\vspace{2mm}

%enditem

%beginitem
4. 
Написати функцію, що записує в масив дійсних послідовність чисел
$$\scalebox{1.2}{$(-1)^{k}C_{n+k}^{k+1}$}, \quad k=0,1,\ldots,n$$
Використовувати вкладені цикли та виклики функцій \textbf{заборонено}.

\vspace{2mm}

%enditem

%beginitem
5. Написати функцію, що для цілого числа знаходить простий дільник, який входить у його розклад у найбільшому степені. Якщо таких кілька, то повернути найменший. \\ Наприклад, для числа $2^{9}{\cdot}3^5{\cdot}7^9{\cdot}11$ функція має повернути $2$.

\vspace{2mm}

%enditem
\end{large}
%end
\newpage
%begin
%\begin {Huge}
{{27.11.2024 Контрольна робота \No 2, варіант  \No \textbf { 45 }\par}}

%\begin{center}Частина 2
%\end{center}
%\vspace{10mm}
%\end {Huge}

{
Якість коду та економність обчислень оцінюється по всім завданням. Послідовності задаються масивами. Використовувати додаткові структури даних (у тому числі рядки) та функції стандартної бібліотеки заборонено.

\vspace{2mm}

\begin {large}
%beginitem
1. Написати функцію, що знаходить у скількох цілих точках інтервалу $(a; b)$ значення функції~$sin(x)$ більше за значення функції~$cos(x)$.

\vspace{2mm}

%enditem

%beginitem
2. Написати функцію, що знаходить найбільший непарний елемент послідовності цілих чисел.

\vspace{2mm}

%enditem

%beginitem
3. Написати функцію, що знаходить кількість входжень великих латинських літер у С-рядок.

\vspace{2mm}

%enditem

%beginitem
4. 
Написати функцію, що записує в масив дійсних послідовність чисел
$$\scalebox{1.2}{$(-1)^{k}C_{n+2k}^{k+1}$}, \quad k=0,1,\ldots,n$$
Використовувати вкладені цикли та виклики функцій \textbf{заборонено}.

\vspace{2mm}

%enditem

%beginitem
5. Написати функцію, що знаходить найбільше число з послідовності цілих, у десятковому запису якого нема цифри 4. Повертає індекс першого входження цього числа або $-1$ (за відсутності).

\vspace{2mm}

%enditem
\end{large}
%end
\vspace{35mm}


%begin
%\begin {Huge}
{{27.11.2024 Контрольна робота \No 2, варіант  \No \textbf { 46 }\par}}

%\begin{center}Частина 2
%\end{center}
%\vspace{10mm}
%\end {Huge}

{
Якість коду та економність обчислень оцінюється по всім завданням. Послідовності задаються масивами. Використовувати додаткові структури даних (у тому числі рядки) та функції стандартної бібліотеки заборонено.

\vspace{2mm}

\begin {large}
%beginitem
1. Написати функцію, що знаходить у скількох цілих точках інтервалу $(a; b]$ функція~$cos(x)$ приймає додатнє значення.

\vspace{2mm}

%enditem

%beginitem
2. Написати функцію, що замінює всі від'ємні елементи послідовності цілих чисел нулями.

\vspace{2mm}

%enditem

%beginitem
3. Написати функцію, що перевіряє, чи містить С-рядок онакову кількість входжень символів ( та ).

\vspace{2mm}

%enditem

%beginitem
4. 
Написати функцію, що записує в масив дійсних послідовність чисел
$$\scalebox{1.2}{$(-1)^{k}C_{3k}^{k}$}, \quad k=2,\ldots,n$$
Використовувати вкладені цикли та виклики функцій \textbf{заборонено}.

\vspace{2mm}

%enditem

%beginitem
5. Написати функцію, що повертає найменше число з послідовності цілих, що має не більше трьох різних простих дільників. За відсутності має повертати $-1$.

\vspace{2mm}

%enditem
\end{large}
%end
\newpage
%begin
%\begin {Huge}
{{27.11.2024 Контрольна робота \No 2, варіант  \No \textbf { 47 }\par}}

%\begin{center}Частина 2
%\end{center}
%\vspace{10mm}
%\end {Huge}

{
Якість коду та економність обчислень оцінюється по всім завданням. Послідовності задаються масивами. Використовувати додаткові структури даних (у тому числі рядки) та функції стандартної бібліотеки заборонено.

\vspace{2mm}

\begin {large}
%beginitem
1. Написати функцію, що знаходить у скількох цілих точках інтервалу $(a; b)$ функція~$sin(x)$ приймає додатнє значення.

\vspace{2mm}

%enditem

%beginitem
2. Написати функцію, що знаходить суму непарних елементів послідовності цілих чисел.

\vspace{2mm}

%enditem

%beginitem
3. Написати функцію, що знаходить кількість входжень великих латинських літер у С-рядок.

\vspace{2mm}

%enditem

%beginitem
4. 
Написати функцію, що записує в масив дійсних послідовність чисел
$$\scalebox{1.2}{$C_{n+2k}^{k}$}, \quad k=0,1,\ldots,n-2$$
Використовувати вкладені цикли та виклики функцій \textbf{заборонено}.

\vspace{2mm}

%enditem

%beginitem
5. Написати функцію, що для інтервалу $[a, b)$ знаходить найближчу пару чисел з цього інтервалу, що кожне з них має парну кількість різних простих дільників. Якщо таких пар кілька, то цікавить пара з найбільшою сумою.

\vspace{2mm}

%enditem
\end{large}
%end
\vspace{35mm}


%begin
%\begin {Huge}
{{27.11.2024 Контрольна робота \No 2, варіант  \No \textbf { 48 }\par}}

%\begin{center}Частина 2
%\end{center}
%\vspace{10mm}
%\end {Huge}

{
Якість коду та економність обчислень оцінюється по всім завданням. Послідовності задаються масивами. Використовувати додаткові структури даних (у тому числі рядки) та функції стандартної бібліотеки заборонено.

\vspace{2mm}

\begin {large}
%beginitem
1. Написати функцію, що знаходить у скількох цілих точках інтервалу $(a; b]$ значення функції~$sin(x)$ більше за значення функції~$cos(x)$.

\vspace{2mm}

%enditem

%beginitem
2. Написати функцію, що знаходить найменший парний елемент послідовності цілих чисел.

\vspace{2mm}

%enditem

%beginitem
3. Написати функцію, що перевіряє, чи містить С-рядок хоч один символ підкреслення.

\vspace{2mm}

%enditem

%beginitem
4. 
Написати функцію, що записує в масив дійсних послідовність чисел
$$\scalebox{1.2}{$(-1)^{k}C_{2k+1}^{k}$}, \quad k=2,\ldots,n$$
Використовувати вкладені цикли та виклики функцій \textbf{заборонено}.

\vspace{2mm}

%enditem

%beginitem
5. Написати функцію, що для цілого числа знаходить суму його простих дільників, що входять у його розклад не більш ніж у 4му степеню. \\ Наприклад, для числа $3^5{\cdot}7^9{\cdot}11{\cdot}23^{2}$ таких дільників два – $11$ та $23$, тому функція має повернути $34$.

\vspace{2mm}

%enditem
\end{large}
%end
\newpage
%begin
%\begin {Huge}
{{27.11.2024 Контрольна робота \No 2, варіант  \No \textbf { 49 }\par}}

%\begin{center}Частина 2
%\end{center}
%\vspace{10mm}
%\end {Huge}

{
Якість коду та економність обчислень оцінюється по всім завданням. Послідовності задаються масивами. Використовувати додаткові структури даних (у тому числі рядки) та функції стандартної бібліотеки заборонено.

\vspace{2mm}

\begin {large}
%beginitem
1. Написати функцію, що знаходить найбільше значення функції~$cos(x)$ у цілих точках інтервалу $(a; b)$.

\vspace{2mm}

%enditem

%beginitem
2. Написати функцію, що замінює всі парні елементи послідовності цілих чисел нулями.

\vspace{2mm}

%enditem

%beginitem
3. Написати функцію, що перевіряє, чи містить С-рядок хоч одну малу латинську літеру.

\vspace{2mm}

%enditem

%beginitem
4. 
Написати функцію, що записує в масив дійсних послідовність чисел
$$\scalebox{1.2}{$(-1)^{k}C_{n+k}^{k+1}$}, \quad k=1,\ldots,n$$
Використовувати вкладені цикли та виклики функцій \textbf{заборонено}.

\vspace{2mm}

%enditem

%beginitem
5. Написати функцію, що для інтервалу $[a, b)$ знаходить кількість чисел, квадрати яких діляться на суму квадратів їх цифр. (Розглядаємо десятковий запис.)

\vspace{2mm}

%enditem
\end{large}
%end
\vspace{35mm}


%begin
%\begin {Huge}
{{27.11.2024 Контрольна робота \No 2, варіант  \No \textbf { 50 }\par}}

%\begin{center}Частина 2
%\end{center}
%\vspace{10mm}
%\end {Huge}

{
Якість коду та економність обчислень оцінюється по всім завданням. Послідовності задаються масивами. Використовувати додаткові структури даних (у тому числі рядки) та функції стандартної бібліотеки заборонено.

\vspace{2mm}

\begin {large}
%beginitem
1. Написати функцію, що знаходить у скількох цілих точках інтервалу $[a; b)$ значення функції~$sin(x)$ більше за значення функції~$cos(x)$.

\vspace{2mm}

%enditem

%beginitem
2. Написати функцію, що замінює всі непарні елементи послідовності цілих чисел числом 14.

\vspace{2mm}

%enditem

%beginitem
3. Написати функцію, що знаходить кількість входжень цифри 0 у С-рядок.

\vspace{2mm}

%enditem

%beginitem
4. 
Написати функцію, що записує в масив дійсних послідовність чисел
$$\scalebox{1.2}{$C_{n+k}^{2k-1}$}, \quad k=1,\ldots,n$$
Використовувати вкладені цикли та виклики функцій \textbf{заборонено}.

\vspace{2mm}

%enditem

%beginitem
5. Написати функцію, що в інтервалі $[a, b)$ знаходить ціле число, що ділиться на найбільший степінь числа 3. Якщо таких чисел декілька, повернути найменше.

\vspace{2mm}

%enditem
\end{large}
%end
\newpage
%begin
%\begin {Huge}
{{27.11.2024 Контрольна робота \No 2, варіант  \No \textbf { 51 }\par}}

%\begin{center}Частина 2
%\end{center}
%\vspace{10mm}
%\end {Huge}

{
Якість коду та економність обчислень оцінюється по всім завданням. Послідовності задаються масивами. Використовувати додаткові структури даних (у тому числі рядки) та функції стандартної бібліотеки заборонено.

\vspace{2mm}

\begin {large}
%beginitem
1. Написати функцію, що знаходить у скількох цілих точках інтервалу $[a; b)$ функція~$sin(x)$ приймає додатнє значення.

\vspace{2mm}

%enditem

%beginitem
2. Написати функцію, що знаходить суму елементів послідовності цілих чисел, що дорівнюють своїм індексам у масиві.

\vspace{2mm}

%enditem

%beginitem
3. Написати функцію, що знаходить кількість входжень літери Z у С-рядок.

\vspace{2mm}

%enditem

%beginitem
4. 
Написати функцію, що записує в масив дійсних послідовність чисел
$$\scalebox{1.2}{$C_{2n+k+1}^{n+k}$}, \quad k=1,\ldots,n$$
Використовувати вкладені цикли та виклики функцій \textbf{заборонено}.

\vspace{2mm}

%enditem

%beginitem
5. Написати функцію, що для інтервалу $[a, b)$ знаходить найближчу пару чисел з цього інтервалу, що кожне з них має непарну кількість різних простих дільників. Якщо таких пар кілька, то цікавить пара з найменшим добутком.

\vspace{2mm}

%enditem
\end{large}
%end
\vspace{35mm}


%begin
%\begin {Huge}
{{27.11.2024 Контрольна робота \No 2, варіант  \No \textbf { 52 }\par}}

%\begin{center}Частина 2
%\end{center}
%\vspace{10mm}
%\end {Huge}

{
Якість коду та економність обчислень оцінюється по всім завданням. Послідовності задаються масивами. Використовувати додаткові структури даних (у тому числі рядки) та функції стандартної бібліотеки заборонено.

\vspace{2mm}

\begin {large}
%beginitem
1. Написати функцію, що знаходить найбільше значення функції~$cos(x)$ у цілих точках інтервалу $[a; b]$.

\vspace{2mm}

%enditem

%beginitem
2. Написати функцію, що замінює всі непарні елементи послідовності цілих чисел нулями.

\vspace{2mm}

%enditem

%beginitem
3. Написати функцію, що знаходить кількість входжень малих латинських літер у С-рядок.

\vspace{2mm}

%enditem

%beginitem
4. 
Написати функцію, що записує в масив дійсних послідовність чисел
$$\scalebox{1.2}{$C_{n+k}^{k+1}$}, \quad k=0,1,\ldots,n-2$$
Використовувати вкладені цикли та виклики функцій \textbf{заборонено}.

\vspace{2mm}

%enditem

%beginitem
5. Написати функцію, що для цілого числа знаходить кількість простих дільників, які входять у його розклад у найбільшому степені. \\ Наприклад, для числа $3^5{\cdot}7^9{\cdot}11{\cdot}23^{9}$ таких дільників два – $23$ та $7$, тому функція має повернути $2$.

\vspace{2mm}

%enditem
\end{large}
%end
\newpage
%begin
%\begin {Huge}
{{27.11.2024 Контрольна робота \No 2, варіант  \No \textbf { 53 }\par}}

%\begin{center}Частина 2
%\end{center}
%\vspace{10mm}
%\end {Huge}

{
Якість коду та економність обчислень оцінюється по всім завданням. Послідовності задаються масивами. Використовувати додаткові структури даних (у тому числі рядки) та функції стандартної бібліотеки заборонено.

\vspace{2mm}

\begin {large}
%beginitem
1. Написати функцію, що знаходить у скількох цілих точках інтервалу $[a; b]$ функція~$sin(x)$ приймає додатнє значення.

\vspace{2mm}

%enditem

%beginitem
2. Написати функцію, що зменшує всі парні елементи послідовності цілих чисел на 3.

\vspace{2mm}

%enditem

%beginitem
3. Написати функцію, що перевіряє, чи містить С-рядок хоч одну велику латинську літеру.

\vspace{2mm}

%enditem

%beginitem
4. 
Написати функцію, що записує в масив дійсних послідовність чисел
$$\scalebox{1.2}{$(-1)^{k}C_{2n+k}^{n+k}$}, \quad k=0,1,\ldots,n$$
Використовувати вкладені цикли та виклики функцій \textbf{заборонено}.

\vspace{2mm}

%enditem

%beginitem
5. Написати функцію, що для цілого числа знаходить простий дільник, який входить у його розклад ц найбільшому степені. Якщо таких кілька, то повернути найбільший. \\ Наприклад, для числа $2^{9}{\cdot}3^5{\cdot}7^9{\cdot}11$ функція має повернути $7$.

\vspace{2mm}

%enditem
\end{large}
%end
\vspace{35mm}


%begin
%\begin {Huge}
{{27.11.2024 Контрольна робота \No 2, варіант  \No \textbf { 54 }\par}}

%\begin{center}Частина 2
%\end{center}
%\vspace{10mm}
%\end {Huge}

{
Якість коду та економність обчислень оцінюється по всім завданням. Послідовності задаються масивами. Використовувати додаткові структури даних (у тому числі рядки) та функції стандартної бібліотеки заборонено.

\vspace{2mm}

\begin {large}
%beginitem
1. Написати функцію, що знаходить найбільше значення функції~$sin(x)$ у цілих точках інтервалу $(a; b]$.

\vspace{2mm}

%enditem

%beginitem
2. Написати функцію, що замінює всі елементи на парних позиціях послідовності цілих чисел (позиції нумеруються з 0) на число 24.

\vspace{2mm}

%enditem

%beginitem
3. Написати функцію, що перевіряє, яких символів у С-рядку більше: 0 чи 1.

\vspace{2mm}

%enditem

%beginitem
4. 
Написати функцію, що записує в масив дійсних послідовність чисел
$$\scalebox{1.2}{$(-1)^{k}C_{n+k}^{2k}$}, \quad k=1,\ldots,n$$
Використовувати вкладені цикли та виклики функцій \textbf{заборонено}.

\vspace{2mm}

%enditem

%beginitem
5. Написати функцію, що для інтервалу $[a, b)$ знаходить знаходить найбільший степінь числа 3, на який ділиться хоча б одне ціле число з цього інтервалу.

\vspace{2mm}

%enditem
\end{large}
%end
\newpage
%begin
%\begin {Huge}
{{27.11.2024 Контрольна робота \No 2, варіант  \No \textbf { 55 }\par}}

%\begin{center}Частина 2
%\end{center}
%\vspace{10mm}
%\end {Huge}

{
Якість коду та економність обчислень оцінюється по всім завданням. Послідовності задаються масивами. Використовувати додаткові структури даних (у тому числі рядки) та функції стандартної бібліотеки заборонено.

\vspace{2mm}

\begin {large}
%beginitem
1. Написати функцію, що знаходить найбільше значення функції~$sin(x)$ у цілих точках інтервалу $[a; b)$.

\vspace{2mm}

%enditem

%beginitem
2. Написати функцію, що знаходить суму квадратів парних елементів послідовності цілих чисел.

\vspace{2mm}

%enditem

%beginitem
3. Написати функцію, що перевіряє, чи містить С-рядок онакову кількість входжень символів ( та ).

\vspace{2mm}

%enditem

%beginitem
4. 
Написати функцію, що записує в масив дійсних послідовність чисел
$$\scalebox{1.2}{$C_{n+k}^{2k}$}, \quad k=1,\ldots,n$$
Використовувати вкладені цикли та виклики функцій \textbf{заборонено}.

\vspace{2mm}

%enditem

%beginitem
5. Написати функцію, яка для цілого числа знаходить кількість його різних простих дільників, що мають у десятковому запису цифру 1.

\vspace{2mm}

%enditem
\end{large}
%end
\vspace{35mm}


%begin
%\begin {Huge}
{{27.11.2024 Контрольна робота \No 2, варіант  \No \textbf { 56 }\par}}

%\begin{center}Частина 2
%\end{center}
%\vspace{10mm}
%\end {Huge}

{
Якість коду та економність обчислень оцінюється по всім завданням. Послідовності задаються масивами. Використовувати додаткові структури даних (у тому числі рядки) та функції стандартної бібліотеки заборонено.

\vspace{2mm}

\begin {large}
%beginitem
1. Написати функцію, що знаходить найбільше значення функції~$sin(x)$ у цілих точках інтервалу $(a; b)$.

\vspace{2mm}

%enditem

%beginitem
2. Написати функцію, що знаходить суму елементів послідовності дійсних чисел, що менші за свої індекси у масиві.

\vspace{2mm}

%enditem

%beginitem
3. Написати функцію, що знаходить кількість входжень цифри 0 у С-рядок.

\vspace{2mm}

%enditem

%beginitem
4. 
Написати функцію, що записує в масив дійсних послідовність чисел
$$\scalebox{1.2}{$(-1)^{k}C_{2n+k+1}^{n+k}$}, \quad k=1,\ldots,n$$
Використовувати вкладені цикли та виклики функцій \textbf{заборонено}.

\vspace{2mm}

%enditem

%beginitem
5. Написати функцію, що для цілого числа знаходить кількість його простих дільників, що входять у його розклад рівно у другому степеню. \\ Наприклад, для числа $3^2{\cdot}7^9{\cdot}11{\cdot}23^{2}$ таких дільників два – $3$ та $23$, тому функція має повернути $2$.

\vspace{2mm}

%enditem
\end{large}
%end
\newpage
%begin
%\begin {Huge}
{{27.11.2024 Контрольна робота \No 2, варіант  \No \textbf { 57 }\par}}

%\begin{center}Частина 2
%\end{center}
%\vspace{10mm}
%\end {Huge}

{
Якість коду та економність обчислень оцінюється по всім завданням. Послідовності задаються масивами. Використовувати додаткові структури даних (у тому числі рядки) та функції стандартної бібліотеки заборонено.

\vspace{2mm}

\begin {large}
%beginitem
1. Написати функцію, що знаходить найбільше значення функції~$sin(x)$ у цілих точках інтервалу $[a; b]$.

\vspace{2mm}

%enditem

%beginitem
2. Написати функцію, що замінює всі від'ємні елементи послідовності цілих чисел нулями.

\vspace{2mm}

%enditem

%beginitem
3. Написати функцію, що перевіряє, чи містить С-рядок хоч один символ підкреслення.

\vspace{2mm}

%enditem

%beginitem
4. 
Написати функцію, що записує в масив дійсних послідовність чисел
$$\scalebox{1.2}{$C_{2k-1}^{k}$}, \quad k=2,\ldots,n$$
Використовувати вкладені цикли та виклики функцій \textbf{заборонено}.

\vspace{2mm}

%enditem

%beginitem
5. Написати функцію, що повертає найбільше число з послідовності цілих, що має не більше трьох різних простих дільників. За відсутності має повертати $-1$.

\vspace{2mm}

%enditem
\end{large}
%end
\vspace{35mm}


%begin
%\begin {Huge}
{{27.11.2024 Контрольна робота \No 2, варіант  \No \textbf { 58 }\par}}

%\begin{center}Частина 2
%\end{center}
%\vspace{10mm}
%\end {Huge}

{
Якість коду та економність обчислень оцінюється по всім завданням. Послідовності задаються масивами. Використовувати додаткові структури даних (у тому числі рядки) та функції стандартної бібліотеки заборонено.

\vspace{2mm}

\begin {large}
%beginitem
1. Написати функцію, що знаходить у скількох цілих точках інтервалу $[a; b)$ функція~$cos(x)$ приймає від'ємне значення.

\vspace{2mm}

%enditem

%beginitem
2. Написати функцію, що збільшує всі парні елементи послідовності цілих чисел на 2.

\vspace{2mm}

%enditem

%beginitem
3. Написати функцію, що знаходить кількість входжень літери Z у С-рядок.

\vspace{2mm}

%enditem

%beginitem
4. 
Написати функцію, що записує в масив дійсних послідовність чисел
$$\scalebox{1.2}{$C_{2n+k}^{n+k}$}, \quad k=1,\ldots,n$$
Використовувати вкладені цикли та виклики функцій \textbf{заборонено}.

\vspace{2mm}

%enditem

%beginitem
5. Написати функцію, що для інтервалу $[a, b)$ знаходить найближчу пару чисел з цього інтервалу, що кожне з них має парну кількість різних простих дільників. Якщо таких пар кілька, то цікавить пара з найменшим добутком.

\vspace{2mm}

%enditem
\end{large}
%end
\newpage
%begin
%\begin {Huge}
{{27.11.2024 Контрольна робота \No 2, варіант  \No \textbf { 59 }\par}}

%\begin{center}Частина 2
%\end{center}
%\vspace{10mm}
%\end {Huge}

{
Якість коду та економність обчислень оцінюється по всім завданням. Послідовності задаються масивами. Використовувати додаткові структури даних (у тому числі рядки) та функції стандартної бібліотеки заборонено.

\vspace{2mm}

\begin {large}
%beginitem
1. Написати функцію, що знаходить у скількох цілих точках інтервалу $(a; b)$ функція~$cos(x)$ приймає додатнє значення.

\vspace{2mm}

%enditem

%beginitem
2. Написати функцію, що знаходить суму елементів послідовності дійсних чисел, що більші 1.

\vspace{2mm}

%enditem

%beginitem
3. Написати функцію, що перевіряє, чи містить С-рядок хоч одну велику латинську літеру.

\vspace{2mm}

%enditem

%beginitem
4. 
Написати функцію, що записує в масив дійсних послідовність чисел
$$\scalebox{1.2}{$C_{n+2k}^{k+1}$}, \quad k=1,\ldots,n$$
Використовувати вкладені цикли та виклики функцій \textbf{заборонено}.

\vspace{2mm}

%enditem

%beginitem
5. Написати функцію, що знаходить найбільше число з послідовності цілих, у десятковому запису якого нема цифри 3. Повертає індекс останнього входження цього числа або $-1$ (за відсутності).

\vspace{2mm}

%enditem
\end{large}
%end
\vspace{35mm}


%begin
%\begin {Huge}
{{27.11.2024 Контрольна робота \No 2, варіант  \No \textbf { 60 }\par}}

%\begin{center}Частина 2
%\end{center}
%\vspace{10mm}
%\end {Huge}

{
Якість коду та економність обчислень оцінюється по всім завданням. Послідовності задаються масивами. Використовувати додаткові структури даних (у тому числі рядки) та функції стандартної бібліотеки заборонено.

\vspace{2mm}

\begin {large}
%beginitem
1. Написати функцію, що знаходить у скількох цілих точках інтервалу $(a; b]$ значення функції~$sin(x)$ менше за значення функції~$cos(x)$.

\vspace{2mm}

%enditem

%beginitem
2. Написати функцію, що знаходить найменший парний елемент послідовності цілих чисел.

\vspace{2mm}

%enditem

%beginitem
3. Написати функцію, що знаходить кількість входжень великих латинських літер у С-рядок.

\vspace{2mm}

%enditem

%beginitem
4. 
Написати функцію, що записує в масив дійсних послідовність чисел
$$\scalebox{1.2}{$C_{3k-1}^{k}$}, \quad k=2,\ldots,n$$
Використовувати вкладені цикли та виклики функцій \textbf{заборонено}.

\vspace{2mm}

%enditem

%beginitem
5. Написати функцію, що для послідовності цілих знаходить довжину найдовшого відрізку, що складається з чисел, у сімковому запису якого нема цифри 2. За відсутності має повертати $0$.

\vspace{2mm}

%enditem
\end{large}
%end
\newpage
%begin
%\begin {Huge}
{{27.11.2024 Контрольна робота \No 2, варіант  \No \textbf { 61 }\par}}

%\begin{center}Частина 2
%\end{center}
%\vspace{10mm}
%\end {Huge}

{
Якість коду та економність обчислень оцінюється по всім завданням. Послідовності задаються масивами. Використовувати додаткові структури даних (у тому числі рядки) та функції стандартної бібліотеки заборонено.

\vspace{2mm}

\begin {large}
%beginitem
1. Написати функцію, що знаходить найбільше значення функції~$cos(x)$ у цілих точках інтервалу $[a; b)$.

\vspace{2mm}

%enditem

%beginitem
2. Написати функцію, що зменшує всі непарні елементи послідовності цілих чисел на 1.

\vspace{2mm}

%enditem

%beginitem
3. Написати функцію, що знаходить кількість входжень малих латинських літер у С-рядок.

\vspace{2mm}

%enditem

%beginitem
4. 
Написати функцію, що записує в масив дійсних послідовність чисел
$$\scalebox{1.2}{$\frac{C_{2k}^k}{k+1}$}, \quad k=0,1,\ldots,n$$
Використовувати вкладені цикли та виклики функцій \textbf{заборонено}.

\vspace{2mm}

%enditem

%beginitem
5. Написати функцію, що для інтервалу $[a, b)$ знаходить найближчу пару чисел з цього інтервалу, що мають не менше трьох простих дільників та мають цифру 3 у своєму десятковому запису.

\vspace{2mm}

%enditem
\end{large}
%end
\vspace{35mm}


%begin
%\begin {Huge}
{{27.11.2024 Контрольна робота \No 2, варіант  \No \textbf { 62 }\par}}

%\begin{center}Частина 2
%\end{center}
%\vspace{10mm}
%\end {Huge}

{
Якість коду та економність обчислень оцінюється по всім завданням. Послідовності задаються масивами. Використовувати додаткові структури даних (у тому числі рядки) та функції стандартної бібліотеки заборонено.

\vspace{2mm}

\begin {large}
%beginitem
1. Написати функцію, що знаходить у скількох цілих точках інтервалу $[a; b)$ значення функції~$sin(x)$ менше за значення функції~$cos(x)$.

\vspace{2mm}

%enditem

%beginitem
2. Написати функцію, що знаходить найбільший непарний елемент послідовності цілих чисел.

\vspace{2mm}

%enditem

%beginitem
3. Написати функцію, що перевіряє, чи містить С-рядок хоч одну малу латинську літеру.

\vspace{2mm}

%enditem

%beginitem
4. 
Написати функцію, що записує в масив дійсних послідовність чисел
$$\scalebox{1.2}{$(-1)^{k}C_{2k-1}^{k}$}, \quad k=2,\ldots,n$$
Використовувати вкладені цикли та виклики функцій \textbf{заборонено}.

\vspace{2mm}

%enditem

%beginitem
5. Написати функцію, що для послідовності цілих находить найдовший відрізок, в якому всі сусідні числа різні. Якщо таких відрізків кілька, то повернути останній (як індекс початку та довжину).

\vspace{2mm}

%enditem
\end{large}
%end
\newpage
%begin
%\begin {Huge}
{{27.11.2024 Контрольна робота \No 2, варіант  \No \textbf { 63 }\par}}

%\begin{center}Частина 2
%\end{center}
%\vspace{10mm}
%\end {Huge}

{
Якість коду та економність обчислень оцінюється по всім завданням. Послідовності задаються масивами. Використовувати додаткові структури даних (у тому числі рядки) та функції стандартної бібліотеки заборонено.

\vspace{2mm}

\begin {large}
%beginitem
1. Написати функцію, що знаходить у скількох цілих точках інтервалу $(a; b]$ функція~$cos(x)$ приймає від'ємне значення.

\vspace{2mm}

%enditem

%beginitem
2. Написати функцію, що замінює всі елементи на непарних позиціях послідовності цілих чисел (позиції нумеруються з 0) на число 25.

\vspace{2mm}

%enditem

%beginitem
3. Написати функцію, що перевіряє, яких символів у С-рядку більше: 0 чи 1.

\vspace{2mm}

%enditem

%beginitem
4. 
Написати функцію, що записує в масив дійсних послідовність чисел
$$\scalebox{1.2}{$C_{n+k}^{k}$}, \quad k=1,\ldots,n$$
Використовувати вкладені цикли та виклики функцій \textbf{заборонено}.

\vspace{2mm}

%enditem

%beginitem
5. Написати функцію, що для послідовності цілих знаходить найдовший відрізок, в якому всі сусідні числа різні. Якщо таких відрізків кілька, то повернути перший (як індекс початку та довжину).

\vspace{2mm}

%enditem
\end{large}
%end
\vspace{35mm}


%begin
%\begin {Huge}
{{27.11.2024 Контрольна робота \No 2, варіант  \No \textbf { 64 }\par}}

%\begin{center}Частина 2
%\end{center}
%\vspace{10mm}
%\end {Huge}

{
Якість коду та економність обчислень оцінюється по всім завданням. Послідовності задаються масивами. Використовувати додаткові структури даних (у тому числі рядки) та функції стандартної бібліотеки заборонено.

\vspace{2mm}

\begin {large}
%beginitem
1. Написати функцію, що знаходить у скількох цілих точках інтервалу $[a; b]$ значення функції~$sin(x)$ більше за значення функції~$cos(x)$.

\vspace{2mm}

%enditem

%beginitem
2. Написати функцію, що зменшує всі парні елементи послідовності цілих чисел на 3.

\vspace{2mm}

%enditem

%beginitem
3. Написати функцію, що перевіряє, чи містить С-рядок хоч одну малу латинську літеру.

\vspace{2mm}

%enditem

%beginitem
4. 
Написати функцію, що записує в масив дійсних послідовність чисел
$$\scalebox{1.2}{$C_{2n+k+1}^{n+k}$}, \quad k=0,1,\ldots,n$$
Використовувати вкладені цикли та виклики функцій \textbf{заборонено}.

\vspace{2mm}

%enditem

%beginitem
5. Написати функцію, що знаходить найменше число з послідовності цілих, у десятковому запису якого нема цифри 5. Повертає індекс останнього входження цього числа або $-1$ (за відсутності).

\vspace{2mm}

%enditem
\end{large}
%end
\newpage
%begin
%\begin {Huge}
{{27.11.2024 Контрольна робота \No 2, варіант  \No \textbf { 65 }\par}}

%\begin{center}Частина 2
%\end{center}
%\vspace{10mm}
%\end {Huge}

{
Якість коду та економність обчислень оцінюється по всім завданням. Послідовності задаються масивами. Використовувати додаткові структури даних (у тому числі рядки) та функції стандартної бібліотеки заборонено.

\vspace{2mm}

\begin {large}
%beginitem
1. Написати функцію, що знаходить найбільше значення функції~$sin(x)$ у цілих точках інтервалу $(a; b]$.

\vspace{2mm}

%enditem

%beginitem
2. Написати функцію, що замінює всі елементи на парних позиціях послідовності цілих чисел (позиції нумеруються з 0) на число 20.

\vspace{2mm}

%enditem

%beginitem
3. Написати функцію, що знаходить кількість входжень цифри 0 у С-рядок.

\vspace{2mm}

%enditem

%beginitem
4. 
Написати функцію, що записує в масив дійсних послідовність чисел
$$\scalebox{1.2}{$(-1)^{k}C_{3k-1}^{k}$}, \quad k=2,\ldots,n$$
Використовувати вкладені цикли та виклики функцій \textbf{заборонено}.

\vspace{2mm}

%enditem

%beginitem
5. Написати функцію, що для інтервалу $[a, b)$ знаходить найближчу пару чисел з цього інтервалу, що кожне з них має непарну кількість різних простих дільників. Якщо таких пар кілька, то цікавить пара з найбільшою сумою.

\vspace{2mm}

%enditem
\end{large}
%end
\vspace{35mm}


%begin
%\begin {Huge}
{{27.11.2024 Контрольна робота \No 2, варіант  \No \textbf { 66 }\par}}

%\begin{center}Частина 2
%\end{center}
%\vspace{10mm}
%\end {Huge}

{
Якість коду та економність обчислень оцінюється по всім завданням. Послідовності задаються масивами. Використовувати додаткові структури даних (у тому числі рядки) та функції стандартної бібліотеки заборонено.

\vspace{2mm}

\begin {large}
%beginitem
1. Написати функцію, що знаходить у скількох цілих точках інтервалу $[a; b]$ функція~$cos(x)$ приймає від'ємне значення.

\vspace{2mm}

%enditem

%beginitem
2. Написати функцію, що замінює всі парні елементи послідовності цілих чисел нулями.

\vspace{2mm}

%enditem

%beginitem
3. Написати функцію, що знаходить кількість входжень літери Z у С-рядок.

\vspace{2mm}

%enditem

%beginitem
4. 
Написати функцію, що записує в масив дійсних послідовність чисел
$$\scalebox{1.2}{$(-1)^{k}C_{2n+k}^{n+k}$}, \quad k=1,\ldots,n$$
Використовувати вкладені цикли та виклики функцій \textbf{заборонено}.

\vspace{2mm}

%enditem

%beginitem
5. Написати функцію, що повертає число з послідовності цілих, що має найменшу кількість різних простих дільників. Якщо таких чисел кілька, то повернути найбільше.

\vspace{2mm}

%enditem
\end{large}
%end
\newpage
%begin
%\begin {Huge}
{{27.11.2024 Контрольна робота \No 2, варіант  \No \textbf { 67 }\par}}

%\begin{center}Частина 2
%\end{center}
%\vspace{10mm}
%\end {Huge}

{
Якість коду та економність обчислень оцінюється по всім завданням. Послідовності задаються масивами. Використовувати додаткові структури даних (у тому числі рядки) та функції стандартної бібліотеки заборонено.

\vspace{2mm}

\begin {large}
%beginitem
1. Написати функцію, що знаходить у скількох цілих точках інтервалу $(a; b]$ функція~$sin(x)$ приймає від'ємне значення.

\vspace{2mm}

%enditem

%beginitem
2. Написати функцію, що замінює всі елементи на парних позиціях послідовності цілих чисел (позиції нумеруються з 0) на число 24.

\vspace{2mm}

%enditem

%beginitem
3. Написати функцію, що перевіряє, яких символів у С-рядку більше: 0 чи 1.

\vspace{2mm}

%enditem

%beginitem
4. 
Написати функцію, що записує в масив дійсних послідовність чисел
$$\scalebox{1.2}{$(-1)^{k}C_{n+k}^{2k}$}, \quad k=0,1,\ldots,n$$
Використовувати вкладені цикли та виклики функцій \textbf{заборонено}.

\vspace{2mm}

%enditem

%beginitem
5. Написати функцію, що повертає найбільше число з послідовності цілих, що має рівно три різних простих дільники. За відсутності має повертати $-1$.

\vspace{2mm}

%enditem
\end{large}
%end
\vspace{35mm}


%begin
%\begin {Huge}
{{27.11.2024 Контрольна робота \No 2, варіант  \No \textbf { 68 }\par}}

%\begin{center}Частина 2
%\end{center}
%\vspace{10mm}
%\end {Huge}

{
Якість коду та економність обчислень оцінюється по всім завданням. Послідовності задаються масивами. Використовувати додаткові структури даних (у тому числі рядки) та функції стандартної бібліотеки заборонено.

\vspace{2mm}

\begin {large}
%beginitem
1. Написати функцію, що знаходить найбільше значення функції~$sin(x)$ у цілих точках інтервалу $[a; b)$.

\vspace{2mm}

%enditem

%beginitem
2. Написати функцію, що знаходить суму елементів послідовності цілих чисел, що дорівнюють своїм індексам у масиві.

\vspace{2mm}

%enditem

%beginitem
3. Написати функцію, що перевіряє, чи містить С-рядок хоч один символ підкреслення.

\vspace{2mm}

%enditem

%beginitem
4. 
Написати функцію, що записує в масив дійсних послідовність чисел
$$\scalebox{1.2}{$(-1)^{k}\frac{C_{2k}^k}{k+1}$}, \quad k=0,1,\ldots,n-2$$
Використовувати вкладені цикли та виклики функцій \textbf{заборонено}.

\vspace{2mm}

%enditem

%beginitem
5. Написати функцію, що для цілого числа знаходить кількість його простих дільників, що входять у його розклад не більш ніж у 5му степеню. \\ Наприклад, для числа $3^8{\cdot}7^9{\cdot}11{\cdot}23^{2}$ таких дільників два – $11$ та $23$, тому функція має повернути $2$.

\vspace{2mm}

%enditem
\end{large}
%end
\newpage
%begin
%\begin {Huge}
{{27.11.2024 Контрольна робота \No 2, варіант  \No \textbf { 69 }\par}}

%\begin{center}Частина 2
%\end{center}
%\vspace{10mm}
%\end {Huge}

{
Якість коду та економність обчислень оцінюється по всім завданням. Послідовності задаються масивами. Використовувати додаткові структури даних (у тому числі рядки) та функції стандартної бібліотеки заборонено.

\vspace{2mm}

\begin {large}
%beginitem
1. Написати функцію, що знаходить у скількох цілих точках інтервалу $(a; b)$ функція~$cos(x)$ приймає від'ємне значення.

\vspace{2mm}

%enditem

%beginitem
2. Написати функцію, що замінює всі додатні елементи послідовності цілих чисел числом 15.

\vspace{2mm}

%enditem

%beginitem
3. Написати функцію, що знаходить кількість входжень великих латинських літер у С-рядок.

\vspace{2mm}

%enditem

%beginitem
4. 
Написати функцію, що записує в масив дійсних послідовність чисел
$$\scalebox{1.2}{$(-1)^{k}C_{n+2k}^{k}$}, \quad k=1,\ldots,n$$
Використовувати вкладені цикли та виклики функцій \textbf{заборонено}.

\vspace{2mm}

%enditem

%beginitem
5. Написати функцію, що для цілого числа знаходить суму його простих дільників, що входять у його розклад не більш ніж у 4му степеню. \\ Наприклад, для числа $3^5{\cdot}7^9{\cdot}11{\cdot}23^{2}$ таких дільників два – $11$ та $23$, тому функція має повернути $34$.

\vspace{2mm}

%enditem
\end{large}
%end
\vspace{35mm}


%begin
%\begin {Huge}
{{27.11.2024 Контрольна робота \No 2, варіант  \No \textbf { 70 }\par}}

%\begin{center}Частина 2
%\end{center}
%\vspace{10mm}
%\end {Huge}

{
Якість коду та економність обчислень оцінюється по всім завданням. Послідовності задаються масивами. Використовувати додаткові структури даних (у тому числі рядки) та функції стандартної бібліотеки заборонено.

\vspace{2mm}

\begin {large}
%beginitem
1. Написати функцію, що знаходить у скількох цілих точках інтервалу $[a; b)$ функція~$sin(x)$ приймає від'ємне значення.

\vspace{2mm}

%enditem

%beginitem
2. Написати функцію, що замінює всі парні елементи послідовності цілих чисел числом 13.

\vspace{2mm}

%enditem

%beginitem
3. Написати функцію, що знаходить кількість входжень малих латинських літер у С-рядок.

\vspace{2mm}

%enditem

%beginitem
4. 
Написати функцію, що записує в масив дійсних послідовність чисел
$$\scalebox{1.2}{$(-1)^{k}C_{n+k}^{2k-1}$}, \quad k=1,\ldots,n$$
Використовувати вкладені цикли та виклики функцій \textbf{заборонено}.

\vspace{2mm}

%enditem

%beginitem
5. Написати функцію, що для послідовності цілих знаходить довжину найдовшого відрізку, що складається з чисел, у сімковому запису якого нема цифри 2. За відсутності має повертати $0$.

\vspace{2mm}

%enditem
\end{large}
%end
\newpage
%begin
%\begin {Huge}
{{27.11.2024 Контрольна робота \No 2, варіант  \No \textbf { 71 }\par}}

%\begin{center}Частина 2
%\end{center}
%\vspace{10mm}
%\end {Huge}

{
Якість коду та економність обчислень оцінюється по всім завданням. Послідовності задаються масивами. Використовувати додаткові структури даних (у тому числі рядки) та функції стандартної бібліотеки заборонено.

\vspace{2mm}

\begin {large}
%beginitem
1. Написати функцію, що знаходить у скількох цілих точках інтервалу $(a; b]$ значення функції~$sin(x)$ менше за значення функції~$cos(x)$.

\vspace{2mm}

%enditem

%beginitem
2. Написати функцію, що знаходить суму елементів послідовності цілих чисел, що більші за свої індекси у масиві.

\vspace{2mm}

%enditem

%beginitem
3. Написати функцію, що перевіряє, чи містить С-рядок онакову кількість входжень символів ( та ).

\vspace{2mm}

%enditem

%beginitem
4. 
Написати функцію, що записує в масив дійсних послідовність чисел
$$\scalebox{1.2}{$(-1)^{k}C_{n+k}^{k}$}, \quad k=1,\ldots,n-2$$
Використовувати вкладені цикли та виклики функцій \textbf{заборонено}.

\vspace{2mm}

%enditem

%beginitem
5. Написати функцію, що для інтервалу $[a, b)$ знаходить найближчу пару чисел з цього інтервалу, що кожне з них має парну кількість різних простих дільників. Якщо таких пар кілька, то цікавить пара з найменшим добутком.

\vspace{2mm}

%enditem
\end{large}
%end
\vspace{35mm}


%begin
%\begin {Huge}
{{27.11.2024 Контрольна робота \No 2, варіант  \No \textbf { 72 }\par}}

%\begin{center}Частина 2
%\end{center}
%\vspace{10mm}
%\end {Huge}

{
Якість коду та економність обчислень оцінюється по всім завданням. Послідовності задаються масивами. Використовувати додаткові структури даних (у тому числі рядки) та функції стандартної бібліотеки заборонено.

\vspace{2mm}

\begin {large}
%beginitem
1. Написати функцію, що знаходить найбільше значення функції~$sin(x)$ у цілих точках інтервалу $[a; b]$.

\vspace{2mm}

%enditem

%beginitem
2. Написати функцію, що знаходить суму непарних елементів послідовності цілих чисел.

\vspace{2mm}

%enditem

%beginitem
3. Написати функцію, що перевіряє, чи містить С-рядок хоч одну велику латинську літеру.

\vspace{2mm}

%enditem

%beginitem
4. 
Написати функцію, що записує в масив дійсних послідовність чисел
$$\scalebox{1.2}{$C_{n+k}^{k+1}$}, \quad k=1,\ldots,n$$
Використовувати вкладені цикли та виклики функцій \textbf{заборонено}.

\vspace{2mm}

%enditem

%beginitem
5. Написати функцію, що для послідовності цілих находить найдовший відрізок, в якому всі сусідні числа різні. Якщо таких відрізків кілька, то повернути останній (як індекс початку та довжину).

\vspace{2mm}

%enditem
\end{large}
%end
\newpage
%begin
%\begin {Huge}
{{27.11.2024 Контрольна робота \No 2, варіант  \No \textbf { 73 }\par}}

%\begin{center}Частина 2
%\end{center}
%\vspace{10mm}
%\end {Huge}

{
Якість коду та економність обчислень оцінюється по всім завданням. Послідовності задаються масивами. Використовувати додаткові структури даних (у тому числі рядки) та функції стандартної бібліотеки заборонено.

\vspace{2mm}

\begin {large}
%beginitem
1. Написати функцію, що знаходить найбільше значення функції~$cos(x)$ у цілих точках інтервалу $(a; b]$.

\vspace{2mm}

%enditem

%beginitem
2. Написати функцію, що знаходить середнє арифметичне елементів послідовності цілих чисел.

\vspace{2mm}

%enditem

%beginitem
3. Написати функцію, що знаходить кількість входжень літери Z у С-рядок.

\vspace{2mm}

%enditem

%beginitem
4. 
Написати функцію, що записує в масив дійсних послідовність чисел
$$\scalebox{1.2}{$C_{n+2k}^{k}$}, \quad k=1,\ldots,n$$
Використовувати вкладені цикли та виклики функцій \textbf{заборонено}.

\vspace{2mm}

%enditem

%beginitem
5. Написати функцію, що для послідовності цілих знаходить найдовший відрізок, в якому всі сусідні числа різні. Якщо таких відрізків кілька, то повернути перший (як індекс початку та довжину).

\vspace{2mm}

%enditem
\end{large}
%end
\vspace{35mm}


%begin
%\begin {Huge}
{{27.11.2024 Контрольна робота \No 2, варіант  \No \textbf { 74 }\par}}

%\begin{center}Частина 2
%\end{center}
%\vspace{10mm}
%\end {Huge}

{
Якість коду та економність обчислень оцінюється по всім завданням. Послідовності задаються масивами. Використовувати додаткові структури даних (у тому числі рядки) та функції стандартної бібліотеки заборонено.

\vspace{2mm}

\begin {large}
%beginitem
1. Написати функцію, що знаходить у скількох цілих точках інтервалу $[a; b)$ функція~$cos(x)$ приймає додатнє значення.

\vspace{2mm}

%enditem

%beginitem
2. Написати функцію, що знаходить суму квадратів елементів послідовності дійсних чисел, що менші 2.

\vspace{2mm}

%enditem

%beginitem
3. Написати функцію, що перевіряє, чи містить С-рядок онакову кількість входжень символів ( та ).

\vspace{2mm}

%enditem

%beginitem
4. 
Написати функцію, що записує в масив дійсних послідовність чисел
$$\scalebox{1.2}{$C_{n+2k}^{k+1}$}, \quad k=0,1,\ldots,n$$
Використовувати вкладені цикли та виклики функцій \textbf{заборонено}.

\vspace{2mm}

%enditem

%beginitem
5. Написати функцію, що для інтервалу $[a, b)$ знаходить знаходить найбільший степінь числа 3, на який ділиться хоча б одне ціле число з цього інтервалу.

\vspace{2mm}

%enditem
\end{large}
%end
\newpage
%begin
%\begin {Huge}
{{27.11.2024 Контрольна робота \No 2, варіант  \No \textbf { 75 }\par}}

%\begin{center}Частина 2
%\end{center}
%\vspace{10mm}
%\end {Huge}

{
Якість коду та економність обчислень оцінюється по всім завданням. Послідовності задаються масивами. Використовувати додаткові структури даних (у тому числі рядки) та функції стандартної бібліотеки заборонено.

\vspace{2mm}

\begin {large}
%beginitem
1. Написати функцію, що знаходить у скількох цілих точках інтервалу $[a; b)$ функція~$sin(x)$ приймає додатнє значення.

\vspace{2mm}

%enditem

%beginitem
2. Написати функцію, що знаходить найбільший парний елемент послідовності цілих чисел.

\vspace{2mm}

%enditem

%beginitem
3. Написати функцію, що знаходить кількість входжень великих латинських літер у С-рядок.

\vspace{2mm}

%enditem

%beginitem
4. 
Написати функцію, що записує в масив дійсних послідовність чисел
$$\scalebox{1.2}{$(-1)^{k}C_{n+2k}^{k}$}, \quad k=0,1,\ldots,n-2$$
Використовувати вкладені цикли та виклики функцій \textbf{заборонено}.

\vspace{2mm}

%enditem

%beginitem
5. Написати функцію, що знаходить найменше число з послідовності цілих, у десятковому запису якого нема цифри 5. Повертає індекс останнього входження цього числа або $-1$ (за відсутності).

\vspace{2mm}

%enditem
\end{large}
%end
\vspace{35mm}


%begin
%\begin {Huge}
{{27.11.2024 Контрольна робота \No 2, варіант  \No \textbf { 76 }\par}}

%\begin{center}Частина 2
%\end{center}
%\vspace{10mm}
%\end {Huge}

{
Якість коду та економність обчислень оцінюється по всім завданням. Послідовності задаються масивами. Використовувати додаткові структури даних (у тому числі рядки) та функції стандартної бібліотеки заборонено.

\vspace{2mm}

\begin {large}
%beginitem
1. Написати функцію, що знаходить у скількох цілих точках інтервалу $[a; b)$ значення функції~$sin(x)$ більше за значення функції~$cos(x)$.

\vspace{2mm}

%enditem

%beginitem
2. Написати функцію, що знаходить суму елементів послідовності цілих чисел, що менші за свої індекси у масиві.

\vspace{2mm}

%enditem

%beginitem
3. Написати функцію, що перевіряє, яких символів у С-рядку більше: 0 чи 1.

\vspace{2mm}

%enditem

%beginitem
4. 
Написати функцію, що записує в масив дійсних послідовність чисел
$$\scalebox{1.2}{$C_{2n+k}^{n+k}$}, \quad k=0,1,\ldots,n$$
Використовувати вкладені цикли та виклики функцій \textbf{заборонено}.

\vspace{2mm}

%enditem

%beginitem
5. Написати функцію, що для інтервалу $[a, b)$ знаходить кількість чисел, квадрати яких діляться на суму квадратів їх цифр. (Розглядаємо десятковий запис.)

\vspace{2mm}

%enditem
\end{large}
%end
\newpage
%begin
%\begin {Huge}
{{27.11.2024 Контрольна робота \No 2, варіант  \No \textbf { 77 }\par}}

%\begin{center}Частина 2
%\end{center}
%\vspace{10mm}
%\end {Huge}

{
Якість коду та економність обчислень оцінюється по всім завданням. Послідовності задаються масивами. Використовувати додаткові структури даних (у тому числі рядки) та функції стандартної бібліотеки заборонено.

\vspace{2mm}

\begin {large}
%beginitem
1. Написати функцію, що знаходить у скількох цілих точках інтервалу $(a; b)$ функція~$sin(x)$ приймає від'ємне значення.

\vspace{2mm}

%enditem

%beginitem
2. Написати функцію, що знаходить найменший непарний елемент послідовності цілих чисел.

\vspace{2mm}

%enditem

%beginitem
3. Написати функцію, що знаходить кількість входжень цифри 0 у С-рядок.

\vspace{2mm}

%enditem

%beginitem
4. 
Написати функцію, що записує в масив дійсних послідовність чисел
$$\scalebox{1.2}{$(-1)^{k}C_{n+k}^{k}$}, \quad k=0,1,\ldots,n$$
Використовувати вкладені цикли та виклики функцій \textbf{заборонено}.

\vspace{2mm}

%enditem

%beginitem
5. Написати функцію, що повертає число з послідовності цілих, що має найменшу кількість різних простих дільників. Якщо таких чисел кілька, то повернути найбільше.

\vspace{2mm}

%enditem
\end{large}
%end
\vspace{35mm}


%begin
%\begin {Huge}
{{27.11.2024 Контрольна робота \No 2, варіант  \No \textbf { 78 }\par}}

%\begin{center}Частина 2
%\end{center}
%\vspace{10mm}
%\end {Huge}

{
Якість коду та економність обчислень оцінюється по всім завданням. Послідовності задаються масивами. Використовувати додаткові структури даних (у тому числі рядки) та функції стандартної бібліотеки заборонено.

\vspace{2mm}

\begin {large}
%beginitem
1. Написати функцію, що знаходить у скількох цілих точках інтервалу $[a; b]$ функція~$sin(x)$ приймає від'ємне значення.

\vspace{2mm}

%enditem

%beginitem
2. Написати функцію, що замінює всі непарні елементи послідовності цілих чисел числом 14.

\vspace{2mm}

%enditem

%beginitem
3. Написати функцію, що перевіряє, чи містить С-рядок хоч один символ підкреслення.

\vspace{2mm}

%enditem

%beginitem
4. 
Написати функцію, що записує в масив дійсних послідовність чисел
$$\scalebox{1.2}{$(-1)^{k}C_{n+2k}^{k+1}$}, \quad k=1,\ldots,n$$
Використовувати вкладені цикли та виклики функцій \textbf{заборонено}.

\vspace{2mm}

%enditem

%beginitem
5. Написати функцію, що для цілого числа знаходить кількість його простих дільників, що входять у його розклад рівно у другому степеню. \\ Наприклад, для числа $3^2{\cdot}7^9{\cdot}11{\cdot}23^{2}$ таких дільників два – $3$ та $23$, тому функція має повернути $2$.

\vspace{2mm}

%enditem
\end{large}
%end
\newpage
%begin
%\begin {Huge}
{{27.11.2024 Контрольна робота \No 2, варіант  \No \textbf { 79 }\par}}

%\begin{center}Частина 2
%\end{center}
%\vspace{10mm}
%\end {Huge}

{
Якість коду та економність обчислень оцінюється по всім завданням. Послідовності задаються масивами. Використовувати додаткові структури даних (у тому числі рядки) та функції стандартної бібліотеки заборонено.

\vspace{2mm}

\begin {large}
%beginitem
1. Написати функцію, що знаходить у скількох цілих точках інтервалу $[a; b)$ значення функції~$sin(x)$ менше за значення функції~$cos(x)$.

\vspace{2mm}

%enditem

%beginitem
2. Написати функцію, що знаходить суму елементів послідовності дійсних чисел, що більші за свої індекси у масиві.

\vspace{2mm}

%enditem

%beginitem
3. Написати функцію, що перевіряє, чи містить С-рядок хоч одну велику латинську літеру.

\vspace{2mm}

%enditem

%beginitem
4. 
Написати функцію, що записує в масив дійсних послідовність чисел
$$\scalebox{1.2}{$C_{n+k}^{2k}$}, \quad k=0,1,\ldots,n$$
Використовувати вкладені цикли та виклики функцій \textbf{заборонено}.

\vspace{2mm}

%enditem

%beginitem
5. Написати функцію, яка для цілого числа знаходить кількість його різних простих дільників, що мають у десятковому запису цифру 1.

\vspace{2mm}

%enditem
\end{large}
%end
\vspace{35mm}


%begin
%\begin {Huge}
{{27.11.2024 Контрольна робота \No 2, варіант  \No \textbf { 80 }\par}}

%\begin{center}Частина 2
%\end{center}
%\vspace{10mm}
%\end {Huge}

{
Якість коду та економність обчислень оцінюється по всім завданням. Послідовності задаються масивами. Використовувати додаткові структури даних (у тому числі рядки) та функції стандартної бібліотеки заборонено.

\vspace{2mm}

\begin {large}
%beginitem
1. Написати функцію, що знаходить у скількох цілих точках інтервалу $[a; b]$ значення функції~$sin(x)$ більше за значення функції~$cos(x)$.

\vspace{2mm}

%enditem

%beginitem
2. Написати функцію, що збільшує всі непарні елементи послідовності цілих чисел на 4.

\vspace{2mm}

%enditem

%beginitem
3. Написати функцію, що знаходить кількість входжень малих латинських літер у С-рядок.

\vspace{2mm}

%enditem

%beginitem
4. 
Написати функцію, що записує в масив дійсних послідовність чисел
$$\scalebox{1.2}{$C_{2k+1}^{k}$}, \quad k=2,\ldots,n$$
Використовувати вкладені цикли та виклики функцій \textbf{заборонено}.

\vspace{2mm}

%enditem

%beginitem
5. Написати функцію, що для інтервалу $[a, b)$ знаходить найближчу пару чисел з цього інтервалу, що кожне з них має парну кількість різних простих дільників. Якщо таких пар кілька, то цікавить пара з найбільшою сумою.

\vspace{2mm}

%enditem
\end{large}
%end
\newpage
%begin
%\begin {Huge}
{{27.11.2024 Контрольна робота \No 2, варіант  \No \textbf { 81 }\par}}

%\begin{center}Частина 2
%\end{center}
%\vspace{10mm}
%\end {Huge}

{
Якість коду та економність обчислень оцінюється по всім завданням. Послідовності задаються масивами. Використовувати додаткові структури даних (у тому числі рядки) та функції стандартної бібліотеки заборонено.

\vspace{2mm}

\begin {large}
%beginitem
1. Написати функцію, що знаходить у скількох цілих точках інтервалу $[a; b)$ функція~$cos(x)$ приймає від'ємне значення.

\vspace{2mm}

%enditem

%beginitem
2. Написати функцію, що замінює всі непарні елементи послідовності цілих чисел нулями.

\vspace{2mm}

%enditem

%beginitem
3. Написати функцію, що перевіряє, чи містить С-рядок хоч одну малу латинську літеру.

\vspace{2mm}

%enditem

%beginitem
4. 
Написати функцію, що записує в масив дійсних послідовність чисел
$$\scalebox{1.2}{$(-1)^{k}C_{2n+k}^{n+k}$}, \quad k=0,1,\ldots,n$$
Використовувати вкладені цикли та виклики функцій \textbf{заборонено}.

\vspace{2mm}

%enditem

%beginitem
5. Написати функцію, що в інтервалі $[a, b)$ знаходить ціле число, що ділиться на найбільший степінь числа 3. Якщо таких чисел декілька, повернути найбільше.

\vspace{2mm}

%enditem
\end{large}
%end
\vspace{35mm}


%begin
%\begin {Huge}
{{27.11.2024 Контрольна робота \No 2, варіант  \No \textbf { 82 }\par}}

%\begin{center}Частина 2
%\end{center}
%\vspace{10mm}
%\end {Huge}

{
Якість коду та економність обчислень оцінюється по всім завданням. Послідовності задаються масивами. Використовувати додаткові структури даних (у тому числі рядки) та функції стандартної бібліотеки заборонено.

\vspace{2mm}

\begin {large}
%beginitem
1. Написати функцію, що знаходить у скількох цілих точках інтервалу $(a; b)$ значення функції~$sin(x)$ менше за значення функції~$cos(x)$.

\vspace{2mm}

%enditem

%beginitem
2. Написати функцію, що збільшує всі парні елементи послідовності цілих чисел на 2.

\vspace{2mm}

%enditem

%beginitem
3. Написати функцію, що знаходить кількість входжень літери Z у С-рядок.

\vspace{2mm}

%enditem

%beginitem
4. 
Написати функцію, що записує в масив дійсних послідовність чисел
$$\scalebox{1.2}{$C_{n+2k}^{k}$}, \quad k=0,1,\ldots,n-2$$
Використовувати вкладені цикли та виклики функцій \textbf{заборонено}.

\vspace{2mm}

%enditem

%beginitem
5. Написати функцію, що знаходить найменше число з послідовності цілих, у десятковому запису якого нема цифри 7. Повертає індекс першого входження цього числа або $-1$ (за відсутності).

\vspace{2mm}

%enditem
\end{large}
%end
\newpage
%begin
%\begin {Huge}
{{27.11.2024 Контрольна робота \No 2, варіант  \No \textbf { 83 }\par}}

%\begin{center}Частина 2
%\end{center}
%\vspace{10mm}
%\end {Huge}

{
Якість коду та економність обчислень оцінюється по всім завданням. Послідовності задаються масивами. Використовувати додаткові структури даних (у тому числі рядки) та функції стандартної бібліотеки заборонено.

\vspace{2mm}

\begin {large}
%beginitem
1. Написати функцію, що знаходить у скількох цілих точках інтервалу $[a; b]$ функція~$cos(x)$ приймає додатнє значення.

\vspace{2mm}

%enditem

%beginitem
2. Написати функцію, що знаходить суму квадратів парних елементів послідовності цілих чисел.

\vspace{2mm}

%enditem

%beginitem
3. Написати функцію, що перевіряє, чи містить С-рядок хоч одну велику латинську літеру.

\vspace{2mm}

%enditem

%beginitem
4. 
Написати функцію, що записує в масив дійсних послідовність чисел
$$\scalebox{1.2}{$C_{n+2k}^{k+1}$}, \quad k=1,\ldots,n$$
Використовувати вкладені цикли та виклики функцій \textbf{заборонено}.

\vspace{2mm}

%enditem

%beginitem
5. Написати функцію, що повертає найбільше число з послідовності цілих, у сімковому запису якого нема цифри 2. За відсутності має повертати $-1$.

\vspace{2mm}

%enditem
\end{large}
%end
\vspace{35mm}


%begin
%\begin {Huge}
{{27.11.2024 Контрольна робота \No 2, варіант  \No \textbf { 84 }\par}}

%\begin{center}Частина 2
%\end{center}
%\vspace{10mm}
%\end {Huge}

{
Якість коду та економність обчислень оцінюється по всім завданням. Послідовності задаються масивами. Використовувати додаткові структури даних (у тому числі рядки) та функції стандартної бібліотеки заборонено.

\vspace{2mm}

\begin {large}
%beginitem
1. Написати функцію, що знаходить у скількох цілих точках інтервалу $(a; b)$ функція~$sin(x)$ приймає додатнє значення.

\vspace{2mm}

%enditem

%beginitem
2. Написати функцію, що знаходить суму елементів послідовності дійсних чисел, що менші за свої індекси у масиві.

\vspace{2mm}

%enditem

%beginitem
3. Написати функцію, що знаходить кількість входжень великих латинських літер у С-рядок.

\vspace{2mm}

%enditem

%beginitem
4. 
Написати функцію, що записує в масив дійсних послідовність чисел
$$\scalebox{1.2}{$(-1)^{k}C_{2k+1}^{k}$}, \quad k=2,\ldots,n$$
Використовувати вкладені цикли та виклики функцій \textbf{заборонено}.

\vspace{2mm}

%enditem

%beginitem
5. Написати функцію, що в інтервалі $[a, b)$ знаходить ціле число, що ділиться на найбільший степінь числа 3. Якщо таких чисел декілька, повернути найменше.

\vspace{2mm}

%enditem
\end{large}
%end
\newpage
%begin
%\begin {Huge}
{{27.11.2024 Контрольна робота \No 2, варіант  \No \textbf { 85 }\par}}

%\begin{center}Частина 2
%\end{center}
%\vspace{10mm}
%\end {Huge}

{
Якість коду та економність обчислень оцінюється по всім завданням. Послідовності задаються масивами. Використовувати додаткові структури даних (у тому числі рядки) та функції стандартної бібліотеки заборонено.

\vspace{2mm}

\begin {large}
%beginitem
1. Написати функцію, що знаходить найбільше значення функції~$sin(x)$ у цілих точках інтервалу $(a; b)$.

\vspace{2mm}

%enditem

%beginitem
2. Написати функцію, що зменшує всі непарні елементи послідовності цілих чисел на 1.

\vspace{2mm}

%enditem

%beginitem
3. Написати функцію, що знаходить кількість входжень цифри 0 у С-рядок.

\vspace{2mm}

%enditem

%beginitem
4. 
Написати функцію, що записує в масив дійсних послідовність чисел
$$\scalebox{1.2}{$C_{n+k}^{k}$}, \quad k=1,\ldots,n$$
Використовувати вкладені цикли та виклики функцій \textbf{заборонено}.

\vspace{2mm}

%enditem

%beginitem
5. Написати функцію, що повертає число з послідовності цілих, що має парну кількість різних простих дільників. Якщо таких чисел кілька, то повернути найбільше.

\vspace{2mm}

%enditem
\end{large}
%end
\vspace{35mm}


%begin
%\begin {Huge}
{{27.11.2024 Контрольна робота \No 2, варіант  \No \textbf { 86 }\par}}

%\begin{center}Частина 2
%\end{center}
%\vspace{10mm}
%\end {Huge}

{
Якість коду та економність обчислень оцінюється по всім завданням. Послідовності задаються масивами. Використовувати додаткові структури даних (у тому числі рядки) та функції стандартної бібліотеки заборонено.

\vspace{2mm}

\begin {large}
%beginitem
1. Написати функцію, що знаходить у скількох цілих точках інтервалу $(a; b]$ функція~$cos(x)$ приймає додатнє значення.

\vspace{2mm}

%enditem

%beginitem
2. Написати функцію, що знаходить найбільший парний елемент послідовності цілих чисел.

\vspace{2mm}

%enditem

%beginitem
3. Написати функцію, що перевіряє, чи містить С-рядок онакову кількість входжень символів ( та ).

\vspace{2mm}

%enditem

%beginitem
4. 
Написати функцію, що записує в масив дійсних послідовність чисел
$$\scalebox{1.2}{$(-1)^{k}C_{n+k}^{k}$}, \quad k=0,1,\ldots,n$$
Використовувати вкладені цикли та виклики функцій \textbf{заборонено}.

\vspace{2mm}

%enditem

%beginitem
5. Написати функцію, що для інтервалу $[a, b)$ знаходить найближчу пару чисел з цього інтервалу, що кожне з них має непарну кількість різних простих дільників. Якщо таких пар кілька, то цікавить пара з найменшим добутком.

\vspace{2mm}

%enditem
\end{large}
%end
\newpage
%begin
%\begin {Huge}
{{27.11.2024 Контрольна робота \No 2, варіант  \No \textbf { 87 }\par}}

%\begin{center}Частина 2
%\end{center}
%\vspace{10mm}
%\end {Huge}

{
Якість коду та економність обчислень оцінюється по всім завданням. Послідовності задаються масивами. Використовувати додаткові структури даних (у тому числі рядки) та функції стандартної бібліотеки заборонено.

\vspace{2mm}

\begin {large}
%beginitem
1. Написати функцію, що знаходить у скількох цілих точках інтервалу $(a; b)$ функція~$cos(x)$ приймає додатнє значення.

\vspace{2mm}

%enditem

%beginitem
2. Написати функцію, що знаходить найменший непарний елемент послідовності цілих чисел.

\vspace{2mm}

%enditem

%beginitem
3. Написати функцію, що перевіряє, яких символів у С-рядку більше: 0 чи 1.

\vspace{2mm}

%enditem

%beginitem
4. 
Написати функцію, що записує в масив дійсних послідовність чисел
$$\scalebox{1.2}{$C_{n+k}^{k}$}, \quad k=0,1,\ldots,n-2$$
Використовувати вкладені цикли та виклики функцій \textbf{заборонено}.

\vspace{2mm}

%enditem

%beginitem
5. Написати функцію, що для цілого числа знаходить кількість простих дільників, які входять у його розклад у найбільшому степені. \\ Наприклад, для числа $3^5{\cdot}7^9{\cdot}11{\cdot}23^{9}$ таких дільників два – $23$ та $7$, тому функція має повернути $2$.

\vspace{2mm}

%enditem
\end{large}
%end
\vspace{35mm}


%begin
%\begin {Huge}
{{27.11.2024 Контрольна робота \No 2, варіант  \No \textbf { 88 }\par}}

%\begin{center}Частина 2
%\end{center}
%\vspace{10mm}
%\end {Huge}

{
Якість коду та економність обчислень оцінюється по всім завданням. Послідовності задаються масивами. Використовувати додаткові структури даних (у тому числі рядки) та функції стандартної бібліотеки заборонено.

\vspace{2mm}

\begin {large}
%beginitem
1. Написати функцію, що знаходить у скількох цілих точках інтервалу $[a; b]$ значення функції~$sin(x)$ менше за значення функції~$cos(x)$.

\vspace{2mm}

%enditem

%beginitem
2. Написати функцію, що знаходить суму елементів послідовності дійсних чисел, що більші за свої індекси у масиві.

\vspace{2mm}

%enditem

%beginitem
3. Написати функцію, що перевіряє, чи містить С-рядок хоч одну малу латинську літеру.

\vspace{2mm}

%enditem

%beginitem
4. 
Написати функцію, що записує в масив дійсних послідовність чисел
$$\scalebox{1.2}{$C_{3k-1}^{k}$}, \quad k=2,\ldots,n$$
Використовувати вкладені цикли та виклики функцій \textbf{заборонено}.

\vspace{2mm}

%enditem

%beginitem
5. Написати функцію, що повертає число з послідовності цілих, що має найбільшу кількість різних простих дільників. Якщо таких чисел кілька, то повернути найменше.

\vspace{2mm}

%enditem
\end{large}
%end
\newpage
%begin
%\begin {Huge}
{{27.11.2024 Контрольна робота \No 2, варіант  \No \textbf { 89 }\par}}

%\begin{center}Частина 2
%\end{center}
%\vspace{10mm}
%\end {Huge}

{
Якість коду та економність обчислень оцінюється по всім завданням. Послідовності задаються масивами. Використовувати додаткові структури даних (у тому числі рядки) та функції стандартної бібліотеки заборонено.

\vspace{2mm}

\begin {large}
%beginitem
1. Написати функцію, що знаходить у скількох цілих точках інтервалу $[a; b]$ функція~$sin(x)$ приймає додатнє значення.

\vspace{2mm}

%enditem

%beginitem
2. Написати функцію, що замінює всі елементи на парних позиціях послідовності цілих чисел (позиції нумеруються з 0) на число 24.

\vspace{2mm}

%enditem

%beginitem
3. Написати функцію, що знаходить кількість входжень малих латинських літер у С-рядок.

\vspace{2mm}

%enditem

%beginitem
4. 
Написати функцію, що записує в масив дійсних послідовність чисел
$$\scalebox{1.2}{$C_{2n+k}^{n+k}$}, \quad k=1,\ldots,n$$
Використовувати вкладені цикли та виклики функцій \textbf{заборонено}.

\vspace{2mm}

%enditem

%beginitem
5. Написати функцію, що для послідовності цілих знаходить довжину найдовшого відрізку, в якому всі сусідні числа відрізняються якнайменш на 2.

\vspace{2mm}

%enditem
\end{large}
%end
\vspace{35mm}


%begin
%\begin {Huge}
{{27.11.2024 Контрольна робота \No 2, варіант  \No \textbf { 90 }\par}}

%\begin{center}Частина 2
%\end{center}
%\vspace{10mm}
%\end {Huge}

{
Якість коду та економність обчислень оцінюється по всім завданням. Послідовності задаються масивами. Використовувати додаткові структури даних (у тому числі рядки) та функції стандартної бібліотеки заборонено.

\vspace{2mm}

\begin {large}
%beginitem
1. Написати функцію, що знаходить у скількох цілих точках інтервалу $(a; b]$ функція~$sin(x)$ приймає додатнє значення.

\vspace{2mm}

%enditem

%beginitem
2. Написати функцію, що знаходить суму елементів послідовності цілих чисел, що менші за свої індекси у масиві.

\vspace{2mm}

%enditem

%beginitem
3. Написати функцію, що перевіряє, чи містить С-рядок хоч один символ підкреслення.

\vspace{2mm}

%enditem

%beginitem
4. 
Написати функцію, що записує в масив дійсних послідовність чисел
$$\scalebox{1.2}{$(-1)^{k}C_{2n+k}^{n+k}$}, \quad k=1,\ldots,n$$
Використовувати вкладені цикли та виклики функцій \textbf{заборонено}.

\vspace{2mm}

%enditem

%beginitem
5. Написати функцію, що для послідовності цілих знаходить довжину найдовшого відрізку, в якому всі сусідні числа різні.

\vspace{2mm}

%enditem
\end{large}
%end
\newpage
%begin
%\begin {Huge}
{{27.11.2024 Контрольна робота \No 2, варіант  \No \textbf { 91 }\par}}

%\begin{center}Частина 2
%\end{center}
%\vspace{10mm}
%\end {Huge}

{
Якість коду та економність обчислень оцінюється по всім завданням. Послідовності задаються масивами. Використовувати додаткові структури даних (у тому числі рядки) та функції стандартної бібліотеки заборонено.

\vspace{2mm}

\begin {large}
%beginitem
1. Написати функцію, що знаходить найбільше значення функції~$cos(x)$ у цілих точках інтервалу $(a; b)$.

\vspace{2mm}

%enditem

%beginitem
2. Написати функцію, що замінює всі парні елементи послідовності цілих чисел нулями.

\vspace{2mm}

%enditem

%beginitem
3. Написати функцію, що перевіряє, чи містить С-рядок хоч одну велику латинську літеру.

\vspace{2mm}

%enditem

%beginitem
4. 
Написати функцію, що записує в масив дійсних послідовність чисел
$$\scalebox{1.2}{$(-1)^{k}C_{n+2k}^{k+1}$}, \quad k=1,\ldots,n$$
Використовувати вкладені цикли та виклики функцій \textbf{заборонено}.

\vspace{2mm}

%enditem

%beginitem
5. Написати функцію, що повертає найбільше число з послідовності цілих, що має рівно три різних простих дільники. За відсутності має повертати $-1$.

\vspace{2mm}

%enditem
\end{large}
%end
\vspace{35mm}


%begin
%\begin {Huge}
{{27.11.2024 Контрольна робота \No 2, варіант  \No \textbf { 92 }\par}}

%\begin{center}Частина 2
%\end{center}
%\vspace{10mm}
%\end {Huge}

{
Якість коду та економність обчислень оцінюється по всім завданням. Послідовності задаються масивами. Використовувати додаткові структури даних (у тому числі рядки) та функції стандартної бібліотеки заборонено.

\vspace{2mm}

\begin {large}
%beginitem
1. Написати функцію, що знаходить у скількох цілих точках інтервалу $(a; b]$ значення функції~$sin(x)$ більше за значення функції~$cos(x)$.

\vspace{2mm}

%enditem

%beginitem
2. Написати функцію, що замінює всі непарні елементи послідовності цілих чисел нулями.

\vspace{2mm}

%enditem

%beginitem
3. Написати функцію, що знаходить кількість входжень малих латинських літер у С-рядок.

\vspace{2mm}

%enditem

%beginitem
4. 
Написати функцію, що записує в масив дійсних послідовність чисел
$$\scalebox{1.2}{$(-1)^{k}C_{2n+k+1}^{n+k}$}, \quad k=1,\ldots,n$$
Використовувати вкладені цикли та виклики функцій \textbf{заборонено}.

\vspace{2mm}

%enditem

%beginitem
5. Написати функцію, що повертає найменше число з послідовності цілих, що має не більше трьох різних простих дільників. За відсутності має повертати $-1$.

\vspace{2mm}

%enditem
\end{large}
%end
\newpage
%begin
%\begin {Huge}
{{27.11.2024 Контрольна робота \No 2, варіант  \No \textbf { 93 }\par}}

%\begin{center}Частина 2
%\end{center}
%\vspace{10mm}
%\end {Huge}

{
Якість коду та економність обчислень оцінюється по всім завданням. Послідовності задаються масивами. Використовувати додаткові структури даних (у тому числі рядки) та функції стандартної бібліотеки заборонено.

\vspace{2mm}

\begin {large}
%beginitem
1. Написати функцію, що знаходить у скількох цілих точках інтервалу $(a; b]$ функція~$cos(x)$ приймає від'ємне значення.

\vspace{2mm}

%enditem

%beginitem
2. Написати функцію, що замінює всі парні елементи послідовності цілих чисел числом 13.

\vspace{2mm}

%enditem

%beginitem
3. Написати функцію, що перевіряє, яких символів у С-рядку більше: 0 чи 1.

\vspace{2mm}

%enditem

%beginitem
4. 
Написати функцію, що записує в масив дійсних послідовність чисел
$$\scalebox{1.2}{$(-1)^{k}C_{2k-1}^{k}$}, \quad k=2,\ldots,n$$
Використовувати вкладені цикли та виклики функцій \textbf{заборонено}.

\vspace{2mm}

%enditem

%beginitem
5. Написати функцію, що знаходить найбільше число з послідовності цілих, у десятковому запису якого нема цифри 3. Повертає індекс останнього входження цього числа або $-1$ (за відсутності).

\vspace{2mm}

%enditem
\end{large}
%end
\vspace{35mm}


%begin
%\begin {Huge}
{{27.11.2024 Контрольна робота \No 2, варіант  \No \textbf { 94 }\par}}

%\begin{center}Частина 2
%\end{center}
%\vspace{10mm}
%\end {Huge}

{
Якість коду та економність обчислень оцінюється по всім завданням. Послідовності задаються масивами. Використовувати додаткові структури даних (у тому числі рядки) та функції стандартної бібліотеки заборонено.

\vspace{2mm}

\begin {large}
%beginitem
1. Написати функцію, що знаходить найбільше значення функції~$cos(x)$ у цілих точках інтервалу $[a; b]$.

\vspace{2mm}

%enditem

%beginitem
2. Написати функцію, що замінює всі додатні елементи послідовності цілих чисел числом 15.

\vspace{2mm}

%enditem

%beginitem
3. Написати функцію, що знаходить кількість входжень літери Z у С-рядок.

\vspace{2mm}

%enditem

%beginitem
4. 
Написати функцію, що записує в масив дійсних послідовність чисел
$$\scalebox{1.2}{$(-1)^{k}C_{n+2k}^{k+1}$}, \quad k=0,1,\ldots,n$$
Використовувати вкладені цикли та виклики функцій \textbf{заборонено}.

\vspace{2mm}

%enditem

%beginitem
5. Написати функцію, що для цілого числа знаходить простий дільник, який входить у його розклад у найбільшому степені. Якщо таких кілька, то повернути найменший. \\ Наприклад, для числа $2^{9}{\cdot}3^5{\cdot}7^9{\cdot}11$ функція має повернути $2$.

\vspace{2mm}

%enditem
\end{large}
%end
\newpage
%begin
%\begin {Huge}
{{27.11.2024 Контрольна робота \No 2, варіант  \No \textbf { 95 }\par}}

%\begin{center}Частина 2
%\end{center}
%\vspace{10mm}
%\end {Huge}

{
Якість коду та економність обчислень оцінюється по всім завданням. Послідовності задаються масивами. Використовувати додаткові структури даних (у тому числі рядки) та функції стандартної бібліотеки заборонено.

\vspace{2mm}

\begin {large}
%beginitem
1. Написати функцію, що знаходить у скількох цілих точках інтервалу $(a; b)$ значення функції~$sin(x)$ більше за значення функції~$cos(x)$.

\vspace{2mm}

%enditem

%beginitem
2. Написати функцію, що знаходить найменший парний елемент послідовності цілих чисел.

\vspace{2mm}

%enditem

%beginitem
3. Написати функцію, що знаходить кількість входжень великих латинських літер у С-рядок.

\vspace{2mm}

%enditem

%beginitem
4. 
Написати функцію, що записує в масив дійсних послідовність чисел
$$\scalebox{1.2}{$C_{n+k}^{2k}$}, \quad k=0,1,\ldots,n$$
Використовувати вкладені цикли та виклики функцій \textbf{заборонено}.

\vspace{2mm}

%enditem

%beginitem
5. Написати функцію, що для послідовності цілих знаходить найдовший відрізок, в якому всі сусідні числа відрізняються якнайменш на 3. Якщо таких відрізків кілька, то повернути перший (як індекс початку та довжину).

\vspace{2mm}

%enditem
\end{large}
%end
\vspace{35mm}


%begin
%\begin {Huge}
{{27.11.2024 Контрольна робота \No 2, варіант  \No \textbf { 96 }\par}}

%\begin{center}Частина 2
%\end{center}
%\vspace{10mm}
%\end {Huge}

{
Якість коду та економність обчислень оцінюється по всім завданням. Послідовності задаються масивами. Використовувати додаткові структури даних (у тому числі рядки) та функції стандартної бібліотеки заборонено.

\vspace{2mm}

\begin {large}
%beginitem
1. Написати функцію, що знаходить найбільше значення функції~$cos(x)$ у цілих точках інтервалу $[a; b)$.

\vspace{2mm}

%enditem

%beginitem
2. Написати функцію, що знаходить суму елементів послідовності дійсних чисел, що більші 1.

\vspace{2mm}

%enditem

%beginitem
3. Написати функцію, що перевіряє, чи містить С-рядок хоч один символ підкреслення.

\vspace{2mm}

%enditem

%beginitem
4. 
Написати функцію, що записує в масив дійсних послідовність чисел
$$\scalebox{1.2}{$C_{2k+1}^{k}$}, \quad k=2,\ldots,n$$
Використовувати вкладені цикли та виклики функцій \textbf{заборонено}.

\vspace{2mm}

%enditem

%beginitem
5. Написати функцію, що для інтервалу $[a, b)$ знаходить найближчу пару чисел з цього інтервалу, що кожне з них має непарну кількість різних простих дільників. Якщо таких пар кілька, то цікавить пара з найбільшою сумою.

\vspace{2mm}

%enditem
\end{large}
%end
\newpage
%begin
%\begin {Huge}
{{27.11.2024 Контрольна робота \No 2, варіант  \No \textbf { 97 }\par}}

%\begin{center}Частина 2
%\end{center}
%\vspace{10mm}
%\end {Huge}

{
Якість коду та економність обчислень оцінюється по всім завданням. Послідовності задаються масивами. Використовувати додаткові структури даних (у тому числі рядки) та функції стандартної бібліотеки заборонено.

\vspace{2mm}

\begin {large}
%beginitem
1. Написати функцію, що знаходить у скількох цілих точках інтервалу $[a; b]$ значення функції~$sin(x)$ більше за значення функції~$cos(x)$.

\vspace{2mm}

%enditem

%beginitem
2. Написати функцію, що знаходить суму елементів послідовності цілих чисел, що більші за свої індекси у масиві.

\vspace{2mm}

%enditem

%beginitem
3. Написати функцію, що перевіряє, чи містить С-рядок хоч одну малу латинську літеру.

\vspace{2mm}

%enditem

%beginitem
4. 
Написати функцію, що записує в масив дійсних послідовність чисел
$$\scalebox{1.2}{$C_{2k-1}^{k}$}, \quad k=2,\ldots,n$$
Використовувати вкладені цикли та виклики функцій \textbf{заборонено}.

\vspace{2mm}

%enditem

%beginitem
5. Написати функцію, що повертає найбільше число з послідовності цілих, що має не більше трьох різних простих дільників. За відсутності має повертати $-1$.

\vspace{2mm}

%enditem
\end{large}
%end
\vspace{35mm}


%begin
%\begin {Huge}
{{27.11.2024 Контрольна робота \No 2, варіант  \No \textbf { 98 }\par}}

%\begin{center}Частина 2
%\end{center}
%\vspace{10mm}
%\end {Huge}

{
Якість коду та економність обчислень оцінюється по всім завданням. Послідовності задаються масивами. Використовувати додаткові структури даних (у тому числі рядки) та функції стандартної бібліотеки заборонено.

\vspace{2mm}

\begin {large}
%beginitem
1. Написати функцію, що знаходить у скількох цілих точках інтервалу $(a; b]$ значення функції~$sin(x)$ більше за значення функції~$cos(x)$.

\vspace{2mm}

%enditem

%beginitem
2. Написати функцію, що замінює всі від'ємні елементи послідовності цілих чисел нулями.

\vspace{2mm}

%enditem

%beginitem
3. Написати функцію, що знаходить кількість входжень цифри 0 у С-рядок.

\vspace{2mm}

%enditem

%beginitem
4. 
Написати функцію, що записує в масив дійсних послідовність чисел
$$\scalebox{1.2}{$C_{3k}^{k}$}, \quad k=2,\ldots,n$$
Використовувати вкладені цикли та виклики функцій \textbf{заборонено}.

\vspace{2mm}

%enditem

%beginitem
5. Написати функцію, що для інтервалу $[a, b)$ знаходить найближчу пару чисел з цього інтервалу, що мають від трьох до п'яти (включно) простих дільників.

\vspace{2mm}

%enditem
\end{large}
%end
\newpage
%begin
%\begin {Huge}
{{27.11.2024 Контрольна робота \No 2, варіант  \No \textbf { 99 }\par}}

%\begin{center}Частина 2
%\end{center}
%\vspace{10mm}
%\end {Huge}

{
Якість коду та економність обчислень оцінюється по всім завданням. Послідовності задаються масивами. Використовувати додаткові структури даних (у тому числі рядки) та функції стандартної бібліотеки заборонено.

\vspace{2mm}

\begin {large}
%beginitem
1. Написати функцію, що знаходить у скількох цілих точках інтервалу $[a; b]$ функція~$sin(x)$ приймає від'ємне значення.

\vspace{2mm}

%enditem

%beginitem
2. Написати функцію, що замінює всі непарні елементи послідовності цілих чисел числом 14.

\vspace{2mm}

%enditem

%beginitem
3. Написати функцію, що перевіряє, чи містить С-рядок онакову кількість входжень символів ( та ).

\vspace{2mm}

%enditem

%beginitem
4. 
Написати функцію, що записує в масив дійсних послідовність чисел
$$\scalebox{1.2}{$(-1)^{k}\frac{C_{2k}^k}{k+1}$}, \quad k=0,1,\ldots,n-2$$
Використовувати вкладені цикли та виклики функцій \textbf{заборонено}.

\vspace{2mm}

%enditem

%beginitem
5. Написати функцію, що для цілого числа знаходить кількість його простих дільників, що входять у його розклад не більш ніж у 5му степеню. \\ Наприклад, для числа $3^8{\cdot}7^9{\cdot}11{\cdot}23^{2}$ таких дільників два – $11$ та $23$, тому функція має повернути $2$.

\vspace{2mm}

%enditem
\end{large}
%end
\vspace{35mm}


%begin
%\begin {Huge}
{{27.11.2024 Контрольна робота \No 2, варіант  \No \textbf { 100 }\par}}

%\begin{center}Частина 2
%\end{center}
%\vspace{10mm}
%\end {Huge}

{
Якість коду та економність обчислень оцінюється по всім завданням. Послідовності задаються масивами. Використовувати додаткові структури даних (у тому числі рядки) та функції стандартної бібліотеки заборонено.

\vspace{2mm}

\begin {large}
%beginitem
1. Написати функцію, що знаходить у скількох цілих точках інтервалу $(a; b]$ функція~$sin(x)$ приймає від'ємне значення.

\vspace{2mm}

%enditem

%beginitem
2. Написати функцію, що збільшує всі непарні елементи послідовності цілих чисел на 4.

\vspace{2mm}

%enditem

%beginitem
3. Написати функцію, що знаходить кількість входжень літери Z у С-рядок.

\vspace{2mm}

%enditem

%beginitem
4. 
Написати функцію, що записує в масив дійсних послідовність чисел
$$\scalebox{1.2}{$(-1)^{k}C_{n+k}^{k}$}, \quad k=1,\ldots,n-2$$
Використовувати вкладені цикли та виклики функцій \textbf{заборонено}.

\vspace{2mm}

%enditem

%beginitem
5. Написати функцію, що для цілого числа знаходить простий дільник, який входить у його розклад ц найбільшому степені. Якщо таких кілька, то повернути найбільший. \\ Наприклад, для числа $2^{9}{\cdot}3^5{\cdot}7^9{\cdot}11$ функція має повернути $7$.

\vspace{2mm}

%enditem
\end{large}
%end
\newpage
%begin
%\begin {Huge}
{{27.11.2024 Контрольна робота \No 2, варіант  \No \textbf { 101 }\par}}

%\begin{center}Частина 2
%\end{center}
%\vspace{10mm}
%\end {Huge}

{
Якість коду та економність обчислень оцінюється по всім завданням. Послідовності задаються масивами. Використовувати додаткові структури даних (у тому числі рядки) та функції стандартної бібліотеки заборонено.

\vspace{2mm}

\begin {large}
%beginitem
1. Написати функцію, що знаходить у скількох цілих точках інтервалу $(a; b)$ функція~$sin(x)$ приймає додатнє значення.

\vspace{2mm}

%enditem

%beginitem
2. Написати функцію, що знаходить найбільший непарний елемент послідовності цілих чисел.

\vspace{2mm}

%enditem

%beginitem
3. Написати функцію, що перевіряє, чи містить С-рядок хоч одну малу латинську літеру.

\vspace{2mm}

%enditem

%beginitem
4. 
Написати функцію, що записує в масив дійсних послідовність чисел
$$\scalebox{1.2}{$C_{n+2k}^{k+1}$}, \quad k=0,1,\ldots,n$$
Використовувати вкладені цикли та виклики функцій \textbf{заборонено}.

\vspace{2mm}

%enditem

%beginitem
5. Написати функцію, що для інтервалу $[a, b)$ знаходить найближчу пару чисел з цього інтервалу, що мають не менше трьох простих дільників та мають цифру 3 у своєму десятковому запису.

\vspace{2mm}

%enditem
\end{large}
%end
\vspace{35mm}


%begin
%\begin {Huge}
{{27.11.2024 Контрольна робота \No 2, варіант  \No \textbf { 102 }\par}}

%\begin{center}Частина 2
%\end{center}
%\vspace{10mm}
%\end {Huge}

{
Якість коду та економність обчислень оцінюється по всім завданням. Послідовності задаються масивами. Використовувати додаткові структури даних (у тому числі рядки) та функції стандартної бібліотеки заборонено.

\vspace{2mm}

\begin {large}
%beginitem
1. Написати функцію, що знаходить у скількох цілих точках інтервалу $(a; b)$ функція~$cos(x)$ приймає від'ємне значення.

\vspace{2mm}

%enditem

%beginitem
2. Написати функцію, що знаходить середнє арифметичне елементів послідовності цілих чисел.

\vspace{2mm}

%enditem

%beginitem
3. Написати функцію, що перевіряє, чи містить С-рядок хоч один символ підкреслення.

\vspace{2mm}

%enditem

%beginitem
4. 
Написати функцію, що записує в масив дійсних послідовність чисел
$$\scalebox{1.2}{$(-1)^{k}C_{n+k}^{2k}$}, \quad k=1,\ldots,n$$
Використовувати вкладені цикли та виклики функцій \textbf{заборонено}.

\vspace{2mm}

%enditem

%beginitem
5. Написати функцію, що знаходить найбільше число з послідовності цілих, у десятковому запису якого нема цифри 4. Повертає індекс першого входження цього числа або $-1$ (за відсутності).

\vspace{2mm}

%enditem
\end{large}
%end
\newpage
%begin
%\begin {Huge}
{{27.11.2024 Контрольна робота \No 2, варіант  \No \textbf { 103 }\par}}

%\begin{center}Частина 2
%\end{center}
%\vspace{10mm}
%\end {Huge}

{
Якість коду та економність обчислень оцінюється по всім завданням. Послідовності задаються масивами. Використовувати додаткові структури даних (у тому числі рядки) та функції стандартної бібліотеки заборонено.

\vspace{2mm}

\begin {large}
%beginitem
1. Написати функцію, що знаходить у скількох цілих точках інтервалу $(a; b]$ значення функції~$sin(x)$ менше за значення функції~$cos(x)$.

\vspace{2mm}

%enditem

%beginitem
2. Написати функцію, що знаходить суму квадратів елементів послідовності дійсних чисел, що менші 2.

\vspace{2mm}

%enditem

%beginitem
3. Написати функцію, що знаходить кількість входжень цифри 0 у С-рядок.

\vspace{2mm}

%enditem

%beginitem
4. 
Написати функцію, що записує в масив дійсних послідовність чисел
$$\scalebox{1.2}{$\frac{C_{2k}^k}{k+1}$}, \quad k=0,1,\ldots,n$$
Використовувати вкладені цикли та виклики функцій \textbf{заборонено}.

\vspace{2mm}

%enditem

%beginitem
5. Написати функцію, що для інтервалу $[a, b)$ знаходить найближчу пару чисел з цього інтервалу, що кожне з них має непарну кількість різних простих дільників. Якщо таких пар кілька, то цікавить пара з найбільшою сумою.

\vspace{2mm}

%enditem
\end{large}
%end
\vspace{35mm}


%begin
%\begin {Huge}
{{27.11.2024 Контрольна робота \No 2, варіант  \No \textbf { 104 }\par}}

%\begin{center}Частина 2
%\end{center}
%\vspace{10mm}
%\end {Huge}

{
Якість коду та економність обчислень оцінюється по всім завданням. Послідовності задаються масивами. Використовувати додаткові структури даних (у тому числі рядки) та функції стандартної бібліотеки заборонено.

\vspace{2mm}

\begin {large}
%beginitem
1. Написати функцію, що знаходить у скількох цілих точках інтервалу $(a; b)$ значення функції~$sin(x)$ більше за значення функції~$cos(x)$.

\vspace{2mm}

%enditem

%beginitem
2. Написати функцію, що знаходить суму елементів послідовності цілих чисел, що дорівнюють своїм індексам у масиві.

\vspace{2mm}

%enditem

%beginitem
3. Написати функцію, що перевіряє, чи містить С-рядок хоч одну велику латинську літеру.

\vspace{2mm}

%enditem

%beginitem
4. 
Написати функцію, що записує в масив дійсних послідовність чисел
$$\scalebox{1.2}{$(-1)^{k}C_{2n+k+1}^{n+k}$}, \quad k=0,1,\ldots,n$$
Використовувати вкладені цикли та виклики функцій \textbf{заборонено}.

\vspace{2mm}

%enditem

%beginitem
5. Написати функцію, що для послідовності цілих знаходить найдовший відрізок, в якому всі сусідні числа різні. Якщо таких відрізків кілька, то повернути перший (як індекс початку та довжину).

\vspace{2mm}

%enditem
\end{large}
%end
\newpage
%begin
%\begin {Huge}
{{27.11.2024 Контрольна робота \No 2, варіант  \No \textbf { 105 }\par}}

%\begin{center}Частина 2
%\end{center}
%\vspace{10mm}
%\end {Huge}

{
Якість коду та економність обчислень оцінюється по всім завданням. Послідовності задаються масивами. Використовувати додаткові структури даних (у тому числі рядки) та функції стандартної бібліотеки заборонено.

\vspace{2mm}

\begin {large}
%beginitem
1. Написати функцію, що знаходить у скількох цілих точках інтервалу $[a; b)$ функція~$cos(x)$ приймає додатнє значення.

\vspace{2mm}

%enditem

%beginitem
2. Написати функцію, що замінює всі елементи на парних позиціях послідовності цілих чисел (позиції нумеруються з 0) на число 20.

\vspace{2mm}

%enditem

%beginitem
3. Написати функцію, що знаходить кількість входжень великих латинських літер у С-рядок.

\vspace{2mm}

%enditem

%beginitem
4. 
Написати функцію, що записує в масив дійсних послідовність чисел
$$\scalebox{1.2}{$(-1)^{k}C_{n+k}^{2k-1}$}, \quad k=1,\ldots,n$$
Використовувати вкладені цикли та виклики функцій \textbf{заборонено}.

\vspace{2mm}

%enditem

%beginitem
5. Написати функцію, що для цілого числа знаходить суму його простих дільників, що входять у його розклад не більш ніж у 4му степеню. \\ Наприклад, для числа $3^5{\cdot}7^9{\cdot}11{\cdot}23^{2}$ таких дільників два – $11$ та $23$, тому функція має повернути $34$.

\vspace{2mm}

%enditem
\end{large}
%end
\vspace{35mm}


%begin
%\begin {Huge}
{{27.11.2024 Контрольна робота \No 2, варіант  \No \textbf { 106 }\par}}

%\begin{center}Частина 2
%\end{center}
%\vspace{10mm}
%\end {Huge}

{
Якість коду та економність обчислень оцінюється по всім завданням. Послідовності задаються масивами. Використовувати додаткові структури даних (у тому числі рядки) та функції стандартної бібліотеки заборонено.

\vspace{2mm}

\begin {large}
%beginitem
1. Написати функцію, що знаходить у скількох цілих точках інтервалу $[a; b)$ функція~$cos(x)$ приймає від'ємне значення.

\vspace{2mm}

%enditem

%beginitem
2. Написати функцію, що замінює всі елементи на непарних позиціях послідовності цілих чисел (позиції нумеруються з 0) на число 25.

\vspace{2mm}

%enditem

%beginitem
3. Написати функцію, що знаходить кількість входжень малих латинських літер у С-рядок.

\vspace{2mm}

%enditem

%beginitem
4. 
Написати функцію, що записує в масив дійсних послідовність чисел
$$\scalebox{1.2}{$(-1)^{k}C_{n+2k}^{k}$}, \quad k=1,\ldots,n$$
Використовувати вкладені цикли та виклики функцій \textbf{заборонено}.

\vspace{2mm}

%enditem

%beginitem
5. Написати функцію, що для послідовності цілих знаходить довжину найдовшого відрізку, в якому всі сусідні числа відрізняються якнайменш на 2.

\vspace{2mm}

%enditem
\end{large}
%end
\newpage
%begin
%\begin {Huge}
{{27.11.2024 Контрольна робота \No 2, варіант  \No \textbf { 107 }\par}}

%\begin{center}Частина 2
%\end{center}
%\vspace{10mm}
%\end {Huge}

{
Якість коду та економність обчислень оцінюється по всім завданням. Послідовності задаються масивами. Використовувати додаткові структури даних (у тому числі рядки) та функції стандартної бібліотеки заборонено.

\vspace{2mm}

\begin {large}
%beginitem
1. Написати функцію, що знаходить найбільше значення функції~$cos(x)$ у цілих точках інтервалу $[a; b]$.

\vspace{2mm}

%enditem

%beginitem
2. Написати функцію, що зменшує всі парні елементи послідовності цілих чисел на 3.

\vspace{2mm}

%enditem

%beginitem
3. Написати функцію, що перевіряє, яких символів у С-рядку більше: 0 чи 1.

\vspace{2mm}

%enditem

%beginitem
4. 
Написати функцію, що записує в масив дійсних послідовність чисел
$$\scalebox{1.2}{$C_{2n+k+1}^{n+k}$}, \quad k=1,\ldots,n$$
Використовувати вкладені цикли та виклики функцій \textbf{заборонено}.

\vspace{2mm}

%enditem

%beginitem
5. Написати функцію, що для послідовності цілих знаходить довжину найдовшого відрізку, в якому всі сусідні числа різні.

\vspace{2mm}

%enditem
\end{large}
%end
\vspace{35mm}


%begin
%\begin {Huge}
{{27.11.2024 Контрольна робота \No 2, варіант  \No \textbf { 108 }\par}}

%\begin{center}Частина 2
%\end{center}
%\vspace{10mm}
%\end {Huge}

{
Якість коду та економність обчислень оцінюється по всім завданням. Послідовності задаються масивами. Використовувати додаткові структури даних (у тому числі рядки) та функції стандартної бібліотеки заборонено.

\vspace{2mm}

\begin {large}
%beginitem
1. Написати функцію, що знаходить найбільше значення функції~$sin(x)$ у цілих точках інтервалу $[a; b]$.

\vspace{2mm}

%enditem

%beginitem
2. Написати функцію, що знаходить суму непарних елементів послідовності цілих чисел.

\vspace{2mm}

%enditem

%beginitem
3. Написати функцію, що перевіряє, чи містить С-рядок онакову кількість входжень символів ( та ).

\vspace{2mm}

%enditem

%beginitem
4. 
Написати функцію, що записує в масив дійсних послідовність чисел
$$\scalebox{1.2}{$C_{n+k}^{2k-1}$}, \quad k=1,\ldots,n$$
Використовувати вкладені цикли та виклики функцій \textbf{заборонено}.

\vspace{2mm}

%enditem

%beginitem
5. Написати функцію, що для інтервалу $[a, b)$ знаходить знаходить найбільший степінь числа 3, на який ділиться хоча б одне ціле число з цього інтервалу.

\vspace{2mm}

%enditem
\end{large}
%end
\newpage
%begin
%\begin {Huge}
{{27.11.2024 Контрольна робота \No 2, варіант  \No \textbf { 109 }\par}}

%\begin{center}Частина 2
%\end{center}
%\vspace{10mm}
%\end {Huge}

{
Якість коду та економність обчислень оцінюється по всім завданням. Послідовності задаються масивами. Використовувати додаткові структури даних (у тому числі рядки) та функції стандартної бібліотеки заборонено.

\vspace{2mm}

\begin {large}
%beginitem
1. Написати функцію, що знаходить у скількох цілих точках інтервалу $[a; b]$ функція~$cos(x)$ приймає від'ємне значення.

\vspace{2mm}

%enditem

%beginitem
2. Написати функцію, що замінює всі непарні елементи послідовності цілих чисел числом 14.

\vspace{2mm}

%enditem

%beginitem
3. Написати функцію, що перевіряє, яких символів у С-рядку більше: 0 чи 1.

\vspace{2mm}

%enditem

%beginitem
4. 
Написати функцію, що записує в масив дійсних послідовність чисел
$$\scalebox{1.2}{$(-1)^{k}C_{n+k}^{2k}$}, \quad k=0,1,\ldots,n$$
Використовувати вкладені цикли та виклики функцій \textbf{заборонено}.

\vspace{2mm}

%enditem

%beginitem
5. Написати функцію, що для інтервалу $[a, b)$ знаходить найближчу пару чисел з цього інтервалу, що кожне з них має парну кількість різних простих дільників. Якщо таких пар кілька, то цікавить пара з найбільшою сумою.

\vspace{2mm}

%enditem
\end{large}
%end
\vspace{35mm}


%begin
%\begin {Huge}
{{27.11.2024 Контрольна робота \No 2, варіант  \No \textbf { 110 }\par}}

%\begin{center}Частина 2
%\end{center}
%\vspace{10mm}
%\end {Huge}

{
Якість коду та економність обчислень оцінюється по всім завданням. Послідовності задаються масивами. Використовувати додаткові структури даних (у тому числі рядки) та функції стандартної бібліотеки заборонено.

\vspace{2mm}

\begin {large}
%beginitem
1. Написати функцію, що знаходить у скількох цілих точках інтервалу $[a; b]$ функція~$sin(x)$ приймає додатнє значення.

\vspace{2mm}

%enditem

%beginitem
2. Написати функцію, що замінює всі елементи на парних позиціях послідовності цілих чисел (позиції нумеруються з 0) на число 24.

\vspace{2mm}

%enditem

%beginitem
3. Написати функцію, що знаходить кількість входжень малих латинських літер у С-рядок.

\vspace{2mm}

%enditem

%beginitem
4. 
Написати функцію, що записує в масив дійсних послідовність чисел
$$\scalebox{1.2}{$C_{2n+k}^{n+k}$}, \quad k=0,1,\ldots,n$$
Використовувати вкладені цикли та виклики функцій \textbf{заборонено}.

\vspace{2mm}

%enditem

%beginitem
5. Написати функцію, що знаходить найменше число з послідовності цілих, у десятковому запису якого нема цифри 7. Повертає індекс першого входження цього числа або $-1$ (за відсутності).

\vspace{2mm}

%enditem
\end{large}
%end
\newpage
%begin
%\begin {Huge}
{{27.11.2024 Контрольна робота \No 2, варіант  \No \textbf { 111 }\par}}

%\begin{center}Частина 2
%\end{center}
%\vspace{10mm}
%\end {Huge}

{
Якість коду та економність обчислень оцінюється по всім завданням. Послідовності задаються масивами. Використовувати додаткові структури даних (у тому числі рядки) та функції стандартної бібліотеки заборонено.

\vspace{2mm}

\begin {large}
%beginitem
1. Написати функцію, що знаходить у скількох цілих точках інтервалу $(a; b]$ функція~$cos(x)$ приймає від'ємне значення.

\vspace{2mm}

%enditem

%beginitem
2. Написати функцію, що замінює всі парні елементи послідовності цілих чисел нулями.

\vspace{2mm}

%enditem

%beginitem
3. Написати функцію, що знаходить кількість входжень цифри 0 у С-рядок.

\vspace{2mm}

%enditem

%beginitem
4. 
Написати функцію, що записує в масив дійсних послідовність чисел
$$\scalebox{1.2}{$C_{n+2k}^{k}$}, \quad k=1,\ldots,n$$
Використовувати вкладені цикли та виклики функцій \textbf{заборонено}.

\vspace{2mm}

%enditem

%beginitem
5. Написати функцію, що для інтервалу $[a, b)$ знаходить кількість чисел, квадрати яких діляться на суму квадратів їх цифр. (Розглядаємо десятковий запис.)

\vspace{2mm}

%enditem
\end{large}
%end
\vspace{35mm}


%begin
%\begin {Huge}
{{27.11.2024 Контрольна робота \No 2, варіант  \No \textbf { 112 }\par}}

%\begin{center}Частина 2
%\end{center}
%\vspace{10mm}
%\end {Huge}

{
Якість коду та економність обчислень оцінюється по всім завданням. Послідовності задаються масивами. Використовувати додаткові структури даних (у тому числі рядки) та функції стандартної бібліотеки заборонено.

\vspace{2mm}

\begin {large}
%beginitem
1. Написати функцію, що знаходить у скількох цілих точках інтервалу $[a; b)$ значення функції~$sin(x)$ більше за значення функції~$cos(x)$.

\vspace{2mm}

%enditem

%beginitem
2. Написати функцію, що замінює всі парні елементи послідовності цілих чисел числом 13.

\vspace{2mm}

%enditem

%beginitem
3. Написати функцію, що перевіряє, чи містить С-рядок хоч одну велику латинську літеру.

\vspace{2mm}

%enditem

%beginitem
4. 
Написати функцію, що записує в масив дійсних послідовність чисел
$$\scalebox{1.2}{$C_{n+k}^{k+1}$}, \quad k=1,\ldots,n$$
Використовувати вкладені цикли та виклики функцій \textbf{заборонено}.

\vspace{2mm}

%enditem

%beginitem
5. Написати функцію, що для інтервалу $[a, b)$ знаходить найближчу пару чисел з цього інтервалу, що кожне з них має непарну кількість різних простих дільників. Якщо таких пар кілька, то цікавить пара з найменшим добутком.

\vspace{2mm}

%enditem
\end{large}
%end
\newpage
%begin
%\begin {Huge}
{{27.11.2024 Контрольна робота \No 2, варіант  \No \textbf { 113 }\par}}

%\begin{center}Частина 2
%\end{center}
%\vspace{10mm}
%\end {Huge}

{
Якість коду та економність обчислень оцінюється по всім завданням. Послідовності задаються масивами. Використовувати додаткові структури даних (у тому числі рядки) та функції стандартної бібліотеки заборонено.

\vspace{2mm}

\begin {large}
%beginitem
1. Написати функцію, що знаходить найбільше значення функції~$cos(x)$ у цілих точках інтервалу $[a; b)$.

\vspace{2mm}

%enditem

%beginitem
2. Написати функцію, що знаходить найменший непарний елемент послідовності цілих чисел.

\vspace{2mm}

%enditem

%beginitem
3. Написати функцію, що перевіряє, чи містить С-рядок хоч один символ підкреслення.

\vspace{2mm}

%enditem

%beginitem
4. 
Написати функцію, що записує в масив дійсних послідовність чисел
$$\scalebox{1.2}{$C_{n+k}^{k+1}$}, \quad k=0,1,\ldots,n-2$$
Використовувати вкладені цикли та виклики функцій \textbf{заборонено}.

\vspace{2mm}

%enditem

%beginitem
5. Написати функцію, що повертає найбільше число з послідовності цілих, у сімковому запису якого нема цифри 2. За відсутності має повертати $-1$.

\vspace{2mm}

%enditem
\end{large}
%end
\vspace{35mm}


%begin
%\begin {Huge}
{{27.11.2024 Контрольна робота \No 2, варіант  \No \textbf { 114 }\par}}

%\begin{center}Частина 2
%\end{center}
%\vspace{10mm}
%\end {Huge}

{
Якість коду та економність обчислень оцінюється по всім завданням. Послідовності задаються масивами. Використовувати додаткові структури даних (у тому числі рядки) та функції стандартної бібліотеки заборонено.

\vspace{2mm}

\begin {large}
%beginitem
1. Написати функцію, що знаходить у скількох цілих точках інтервалу $[a; b]$ значення функції~$sin(x)$ менше за значення функції~$cos(x)$.

\vspace{2mm}

%enditem

%beginitem
2. Написати функцію, що знаходить суму елементів послідовності дійсних чисел, що більші за свої індекси у масиві.

\vspace{2mm}

%enditem

%beginitem
3. Написати функцію, що знаходить кількість входжень великих латинських літер у С-рядок.

\vspace{2mm}

%enditem

%beginitem
4. 
Написати функцію, що записує в масив дійсних послідовність чисел
$$\scalebox{1.2}{$C_{2n+k+1}^{n+k}$}, \quad k=0,1,\ldots,n$$
Використовувати вкладені цикли та виклики функцій \textbf{заборонено}.

\vspace{2mm}

%enditem

%beginitem
5. Написати функцію, що повертає найбільше число з послідовності цілих, що має рівно три різних простих дільники. За відсутності має повертати $-1$.

\vspace{2mm}

%enditem
\end{large}
%end
\newpage
%begin
%\begin {Huge}
{{27.11.2024 Контрольна робота \No 2, варіант  \No \textbf { 115 }\par}}

%\begin{center}Частина 2
%\end{center}
%\vspace{10mm}
%\end {Huge}

{
Якість коду та економність обчислень оцінюється по всім завданням. Послідовності задаються масивами. Використовувати додаткові структури даних (у тому числі рядки) та функції стандартної бібліотеки заборонено.

\vspace{2mm}

\begin {large}
%beginitem
1. Написати функцію, що знаходить у скількох цілих точках інтервалу $(a; b)$ функція~$sin(x)$ приймає від'ємне значення.

\vspace{2mm}

%enditem

%beginitem
2. Написати функцію, що знаходить суму квадратів елементів послідовності дійсних чисел, що менші 2.

\vspace{2mm}

%enditem

%beginitem
3. Написати функцію, що знаходить кількість входжень літери Z у С-рядок.

\vspace{2mm}

%enditem

%beginitem
4. 
Написати функцію, що записує в масив дійсних послідовність чисел
$$\scalebox{1.2}{$(-1)^{k}C_{n+k}^{k+1}$}, \quad k=0,1,\ldots,n$$
Використовувати вкладені цикли та виклики функцій \textbf{заборонено}.

\vspace{2mm}

%enditem

%beginitem
5. Написати функцію, що для інтервалу $[a, b)$ знаходить найближчу пару чисел з цього інтервалу, що мають від трьох до п'яти (включно) простих дільників.

\vspace{2mm}

%enditem
\end{large}
%end
\vspace{35mm}


%begin
%\begin {Huge}
{{27.11.2024 Контрольна робота \No 2, варіант  \No \textbf { 116 }\par}}

%\begin{center}Частина 2
%\end{center}
%\vspace{10mm}
%\end {Huge}

{
Якість коду та економність обчислень оцінюється по всім завданням. Послідовності задаються масивами. Використовувати додаткові структури даних (у тому числі рядки) та функції стандартної бібліотеки заборонено.

\vspace{2mm}

\begin {large}
%beginitem
1. Написати функцію, що знаходить найбільше значення функції~$cos(x)$ у цілих точках інтервалу $(a; b)$.

\vspace{2mm}

%enditem

%beginitem
2. Написати функцію, що замінює всі елементи на непарних позиціях послідовності цілих чисел (позиції нумеруються з 0) на число 25.

\vspace{2mm}

%enditem

%beginitem
3. Написати функцію, що перевіряє, чи містить С-рядок хоч одну малу латинську літеру.

\vspace{2mm}

%enditem

%beginitem
4. 
Написати функцію, що записує в масив дійсних послідовність чисел
$$\scalebox{1.2}{$(-1)^{k}C_{n+k}^{k+1}$}, \quad k=1,\ldots,n$$
Використовувати вкладені цикли та виклики функцій \textbf{заборонено}.

\vspace{2mm}

%enditem

%beginitem
5. Написати функцію, що повертає найменше число з послідовності цілих, що має не більше трьох різних простих дільників. За відсутності має повертати $-1$.

\vspace{2mm}

%enditem
\end{large}
%end
\newpage
%begin
%\begin {Huge}
{{27.11.2024 Контрольна робота \No 2, варіант  \No \textbf { 117 }\par}}

%\begin{center}Частина 2
%\end{center}
%\vspace{10mm}
%\end {Huge}

{
Якість коду та економність обчислень оцінюється по всім завданням. Послідовності задаються масивами. Використовувати додаткові структури даних (у тому числі рядки) та функції стандартної бібліотеки заборонено.

\vspace{2mm}

\begin {large}
%beginitem
1. Написати функцію, що знаходить у скількох цілих точках інтервалу $(a; b]$ функція~$sin(x)$ приймає додатнє значення.

\vspace{2mm}

%enditem

%beginitem
2. Написати функцію, що знаходить суму елементів послідовності цілих чисел, що менші за свої індекси у масиві.

\vspace{2mm}

%enditem

%beginitem
3. Написати функцію, що перевіряє, чи містить С-рядок онакову кількість входжень символів ( та ).

\vspace{2mm}

%enditem

%beginitem
4. 
Написати функцію, що записує в масив дійсних послідовність чисел
$$\scalebox{1.2}{$(-1)^{k}C_{3k}^{k}$}, \quad k=2,\ldots,n$$
Використовувати вкладені цикли та виклики функцій \textbf{заборонено}.

\vspace{2mm}

%enditem

%beginitem
5. Написати функцію, що повертає найбільше число з послідовності цілих, що має не більше трьох різних простих дільників. За відсутності має повертати $-1$.

\vspace{2mm}

%enditem
\end{large}
%end
\vspace{35mm}


%begin
%\begin {Huge}
{{27.11.2024 Контрольна робота \No 2, варіант  \No \textbf { 118 }\par}}

%\begin{center}Частина 2
%\end{center}
%\vspace{10mm}
%\end {Huge}

{
Якість коду та економність обчислень оцінюється по всім завданням. Послідовності задаються масивами. Використовувати додаткові структури даних (у тому числі рядки) та функції стандартної бібліотеки заборонено.

\vspace{2mm}

\begin {large}
%beginitem
1. Написати функцію, що знаходить найбільше значення функції~$sin(x)$ у цілих точках інтервалу $(a; b]$.

\vspace{2mm}

%enditem

%beginitem
2. Написати функцію, що знаходить суму елементів послідовності цілих чисел, що більші за свої індекси у масиві.

\vspace{2mm}

%enditem

%beginitem
3. Написати функцію, що знаходить кількість входжень малих латинських літер у С-рядок.

\vspace{2mm}

%enditem

%beginitem
4. 
Написати функцію, що записує в масив дійсних послідовність чисел
$$\scalebox{1.2}{$(-1)^{k}C_{3k-1}^{k}$}, \quad k=2,\ldots,n$$
Використовувати вкладені цикли та виклики функцій \textbf{заборонено}.

\vspace{2mm}

%enditem

%beginitem
5. Написати функцію, що для цілого числа знаходить кількість його простих дільників, що входять у його розклад не більш ніж у 5му степеню. \\ Наприклад, для числа $3^8{\cdot}7^9{\cdot}11{\cdot}23^{2}$ таких дільників два – $11$ та $23$, тому функція має повернути $2$.

\vspace{2mm}

%enditem
\end{large}
%end
\newpage
%begin
%\begin {Huge}
{{27.11.2024 Контрольна робота \No 2, варіант  \No \textbf { 119 }\par}}

%\begin{center}Частина 2
%\end{center}
%\vspace{10mm}
%\end {Huge}

{
Якість коду та економність обчислень оцінюється по всім завданням. Послідовності задаються масивами. Використовувати додаткові структури даних (у тому числі рядки) та функції стандартної бібліотеки заборонено.

\vspace{2mm}

\begin {large}
%beginitem
1. Написати функцію, що знаходить у скількох цілих точках інтервалу $[a; b)$ функція~$sin(x)$ приймає від'ємне значення.

\vspace{2mm}

%enditem

%beginitem
2. Написати функцію, що замінює всі елементи на парних позиціях послідовності цілих чисел (позиції нумеруються з 0) на число 20.

\vspace{2mm}

%enditem

%beginitem
3. Написати функцію, що знаходить кількість входжень літери Z у С-рядок.

\vspace{2mm}

%enditem

%beginitem
4. 
Написати функцію, що записує в масив дійсних послідовність чисел
$$\scalebox{1.2}{$(-1)^{k}C_{n+2k}^{k}$}, \quad k=0,1,\ldots,n-2$$
Використовувати вкладені цикли та виклики функцій \textbf{заборонено}.

\vspace{2mm}

%enditem

%beginitem
5. Написати функцію, що знаходить найбільше число з послідовності цілих, у десятковому запису якого нема цифри 3. Повертає індекс останнього входження цього числа або $-1$ (за відсутності).

\vspace{2mm}

%enditem
\end{large}
%end
\vspace{35mm}


%begin
%\begin {Huge}
{{27.11.2024 Контрольна робота \No 2, варіант  \No \textbf { 120 }\par}}

%\begin{center}Частина 2
%\end{center}
%\vspace{10mm}
%\end {Huge}

{
Якість коду та економність обчислень оцінюється по всім завданням. Послідовності задаються масивами. Використовувати додаткові структури даних (у тому числі рядки) та функції стандартної бібліотеки заборонено.

\vspace{2mm}

\begin {large}
%beginitem
1. Написати функцію, що знаходить у скількох цілих точках інтервалу $[a; b]$ функція~$cos(x)$ приймає додатнє значення.

\vspace{2mm}

%enditem

%beginitem
2. Написати функцію, що знаходить суму елементів послідовності дійсних чисел, що менші за свої індекси у масиві.

\vspace{2mm}

%enditem

%beginitem
3. Написати функцію, що перевіряє, чи містить С-рядок хоч один символ підкреслення.

\vspace{2mm}

%enditem

%beginitem
4. 
Написати функцію, що записує в масив дійсних послідовність чисел
$$\scalebox{1.2}{$C_{n+k}^{2k}$}, \quad k=1,\ldots,n$$
Використовувати вкладені цикли та виклики функцій \textbf{заборонено}.

\vspace{2mm}

%enditem

%beginitem
5. Написати функцію, що для послідовності цілих знаходить найдовший відрізок, в якому всі сусідні числа відрізняються якнайменш на 3. Якщо таких відрізків кілька, то повернути перший (як індекс початку та довжину).

\vspace{2mm}

%enditem
\end{large}
%end
\newpage
%begin
%\begin {Huge}
{{27.11.2024 Контрольна робота \No 2, варіант  \No \textbf { 121 }\par}}

%\begin{center}Частина 2
%\end{center}
%\vspace{10mm}
%\end {Huge}

{
Якість коду та економність обчислень оцінюється по всім завданням. Послідовності задаються масивами. Використовувати додаткові структури даних (у тому числі рядки) та функції стандартної бібліотеки заборонено.

\vspace{2mm}

\begin {large}
%beginitem
1. Написати функцію, що знаходить найбільше значення функції~$sin(x)$ у цілих точках інтервалу $[a; b)$.

\vspace{2mm}

%enditem

%beginitem
2. Написати функцію, що знаходить найбільший непарний елемент послідовності цілих чисел.

\vspace{2mm}

%enditem

%beginitem
3. Написати функцію, що знаходить кількість входжень цифри 0 у С-рядок.

\vspace{2mm}

%enditem

%beginitem
4. 
Написати функцію, що записує в масив дійсних послідовність чисел
$$\scalebox{1.2}{$(-1)^{k}C_{n+2k}^{k}$}, \quad k=0,1,\ldots,n-2$$
Використовувати вкладені цикли та виклики функцій \textbf{заборонено}.

\vspace{2mm}

%enditem

%beginitem
5. Написати функцію, що знаходить найменше число з послідовності цілих, у десятковому запису якого нема цифри 5. Повертає індекс останнього входження цього числа або $-1$ (за відсутності).

\vspace{2mm}

%enditem
\end{large}
%end
\vspace{35mm}


%begin
%\begin {Huge}
{{27.11.2024 Контрольна робота \No 2, варіант  \No \textbf { 122 }\par}}

%\begin{center}Частина 2
%\end{center}
%\vspace{10mm}
%\end {Huge}

{
Якість коду та економність обчислень оцінюється по всім завданням. Послідовності задаються масивами. Використовувати додаткові структури даних (у тому числі рядки) та функції стандартної бібліотеки заборонено.

\vspace{2mm}

\begin {large}
%beginitem
1. Написати функцію, що знаходить найбільше значення функції~$sin(x)$ у цілих точках інтервалу $(a; b)$.

\vspace{2mm}

%enditem

%beginitem
2. Написати функцію, що замінює всі непарні елементи послідовності цілих чисел нулями.

\vspace{2mm}

%enditem

%beginitem
3. Написати функцію, що перевіряє, яких символів у С-рядку більше: 0 чи 1.

\vspace{2mm}

%enditem

%beginitem
4. 
Написати функцію, що записує в масив дійсних послідовність чисел
$$\scalebox{1.2}{$C_{n+k}^{k+1}$}, \quad k=1,\ldots,n$$
Використовувати вкладені цикли та виклики функцій \textbf{заборонено}.

\vspace{2mm}

%enditem

%beginitem
5. Написати функцію, що для послідовності цілих находить найдовший відрізок, в якому всі сусідні числа різні. Якщо таких відрізків кілька, то повернути останній (як індекс початку та довжину).

\vspace{2mm}

%enditem
\end{large}
%end
\newpage
%begin
%\begin {Huge}
{{27.11.2024 Контрольна робота \No 2, варіант  \No \textbf { 123 }\par}}

%\begin{center}Частина 2
%\end{center}
%\vspace{10mm}
%\end {Huge}

{
Якість коду та економність обчислень оцінюється по всім завданням. Послідовності задаються масивами. Використовувати додаткові структури даних (у тому числі рядки) та функції стандартної бібліотеки заборонено.

\vspace{2mm}

\begin {large}
%beginitem
1. Написати функцію, що знаходить у скількох цілих точках інтервалу $(a; b]$ функція~$cos(x)$ приймає додатнє значення.

\vspace{2mm}

%enditem

%beginitem
2. Написати функцію, що знаходить найменший парний елемент послідовності цілих чисел.

\vspace{2mm}

%enditem

%beginitem
3. Написати функцію, що перевіряє, чи містить С-рядок хоч одну велику латинську літеру.

\vspace{2mm}

%enditem

%beginitem
4. 
Написати функцію, що записує в масив дійсних послідовність чисел
$$\scalebox{1.2}{$C_{3k-1}^{k}$}, \quad k=2,\ldots,n$$
Використовувати вкладені цикли та виклики функцій \textbf{заборонено}.

\vspace{2mm}

%enditem

%beginitem
5. Написати функцію, що для цілого числа знаходить кількість простих дільників, які входять у його розклад у найбільшому степені. \\ Наприклад, для числа $3^5{\cdot}7^9{\cdot}11{\cdot}23^{9}$ таких дільників два – $23$ та $7$, тому функція має повернути $2$.

\vspace{2mm}

%enditem
\end{large}
%end
\vspace{35mm}


%begin
%\begin {Huge}
{{27.11.2024 Контрольна робота \No 2, варіант  \No \textbf { 124 }\par}}

%\begin{center}Частина 2
%\end{center}
%\vspace{10mm}
%\end {Huge}

{
Якість коду та економність обчислень оцінюється по всім завданням. Послідовності задаються масивами. Використовувати додаткові структури даних (у тому числі рядки) та функції стандартної бібліотеки заборонено.

\vspace{2mm}

\begin {large}
%beginitem
1. Написати функцію, що знаходить у скількох цілих точках інтервалу $[a; b)$ функція~$sin(x)$ приймає додатнє значення.

\vspace{2mm}

%enditem

%beginitem
2. Написати функцію, що збільшує всі парні елементи послідовності цілих чисел на 2.

\vspace{2mm}

%enditem

%beginitem
3. Написати функцію, що перевіряє, чи містить С-рядок онакову кількість входжень символів ( та ).

\vspace{2mm}

%enditem

%beginitem
4. 
Написати функцію, що записує в масив дійсних послідовність чисел
$$\scalebox{1.2}{$(-1)^{k}C_{2k-1}^{k}$}, \quad k=2,\ldots,n$$
Використовувати вкладені цикли та виклики функцій \textbf{заборонено}.

\vspace{2mm}

%enditem

%beginitem
5. Написати функцію, що повертає число з послідовності цілих, що має найменшу кількість різних простих дільників. Якщо таких чисел кілька, то повернути найбільше.

\vspace{2mm}

%enditem
\end{large}
%end
\newpage
%begin
%\begin {Huge}
{{27.11.2024 Контрольна робота \No 2, варіант  \No \textbf { 125 }\par}}

%\begin{center}Частина 2
%\end{center}
%\vspace{10mm}
%\end {Huge}

{
Якість коду та економність обчислень оцінюється по всім завданням. Послідовності задаються масивами. Використовувати додаткові структури даних (у тому числі рядки) та функції стандартної бібліотеки заборонено.

\vspace{2mm}

\begin {large}
%beginitem
1. Написати функцію, що знаходить у скількох цілих точках інтервалу $(a; b)$ функція~$cos(x)$ приймає додатнє значення.

\vspace{2mm}

%enditem

%beginitem
2. Написати функцію, що зменшує всі парні елементи послідовності цілих чисел на 3.

\vspace{2mm}

%enditem

%beginitem
3. Написати функцію, що перевіряє, чи містить С-рядок хоч одну малу латинську літеру.

\vspace{2mm}

%enditem

%beginitem
4. 
Написати функцію, що записує в масив дійсних послідовність чисел
$$\scalebox{1.2}{$(-1)^{k}C_{3k}^{k}$}, \quad k=2,\ldots,n$$
Використовувати вкладені цикли та виклики функцій \textbf{заборонено}.

\vspace{2mm}

%enditem

%beginitem
5. Написати функцію, що знаходить найбільше число з послідовності цілих, у десятковому запису якого нема цифри 4. Повертає індекс першого входження цього числа або $-1$ (за відсутності).

\vspace{2mm}

%enditem
\end{large}
%end
\vspace{35mm}


%begin
%\begin {Huge}
{{27.11.2024 Контрольна робота \No 2, варіант  \No \textbf { 126 }\par}}

%\begin{center}Частина 2
%\end{center}
%\vspace{10mm}
%\end {Huge}

{
Якість коду та економність обчислень оцінюється по всім завданням. Послідовності задаються масивами. Використовувати додаткові структури даних (у тому числі рядки) та функції стандартної бібліотеки заборонено.

\vspace{2mm}

\begin {large}
%beginitem
1. Написати функцію, що знаходить найбільше значення функції~$cos(x)$ у цілих точках інтервалу $(a; b]$.

\vspace{2mm}

%enditem

%beginitem
2. Написати функцію, що зменшує всі непарні елементи послідовності цілих чисел на 1.

\vspace{2mm}

%enditem

%beginitem
3. Написати функцію, що знаходить кількість входжень великих латинських літер у С-рядок.

\vspace{2mm}

%enditem

%beginitem
4. 
Написати функцію, що записує в масив дійсних послідовність чисел
$$\scalebox{1.2}{$C_{n+k}^{2k}$}, \quad k=1,\ldots,n$$
Використовувати вкладені цикли та виклики функцій \textbf{заборонено}.

\vspace{2mm}

%enditem

%beginitem
5. Написати функцію, що для послідовності цілих знаходить довжину найдовшого відрізку, що складається з чисел, у сімковому запису якого нема цифри 2. За відсутності має повертати $0$.

\vspace{2mm}

%enditem
\end{large}
%end
\newpage
%begin
%\begin {Huge}
{{27.11.2024 Контрольна робота \No 2, варіант  \No \textbf { 127 }\par}}

%\begin{center}Частина 2
%\end{center}
%\vspace{10mm}
%\end {Huge}

{
Якість коду та економність обчислень оцінюється по всім завданням. Послідовності задаються масивами. Використовувати додаткові структури даних (у тому числі рядки) та функції стандартної бібліотеки заборонено.

\vspace{2mm}

\begin {large}
%beginitem
1. Написати функцію, що знаходить у скількох цілих точках інтервалу $[a; b)$ значення функції~$sin(x)$ менше за значення функції~$cos(x)$.

\vspace{2mm}

%enditem

%beginitem
2. Написати функцію, що знаходить суму квадратів парних елементів послідовності цілих чисел.

\vspace{2mm}

%enditem

%beginitem
3. Написати функцію, що знаходить кількість входжень великих латинських літер у С-рядок.

\vspace{2mm}

%enditem

%beginitem
4. 
Написати функцію, що записує в масив дійсних послідовність чисел
$$\scalebox{1.2}{$(-1)^{k}C_{n+k}^{k+1}$}, \quad k=1,\ldots,n$$
Використовувати вкладені цикли та виклики функцій \textbf{заборонено}.

\vspace{2mm}

%enditem

%beginitem
5. Написати функцію, що для інтервалу $[a, b)$ знаходить найближчу пару чисел з цього інтервалу, що кожне з них має парну кількість різних простих дільників. Якщо таких пар кілька, то цікавить пара з найменшим добутком.

\vspace{2mm}

%enditem
\end{large}
%end
\vspace{35mm}


%begin
%\begin {Huge}
{{27.11.2024 Контрольна робота \No 2, варіант  \No \textbf { 128 }\par}}

%\begin{center}Частина 2
%\end{center}
%\vspace{10mm}
%\end {Huge}

{
Якість коду та економність обчислень оцінюється по всім завданням. Послідовності задаються масивами. Використовувати додаткові структури даних (у тому числі рядки) та функції стандартної бібліотеки заборонено.

\vspace{2mm}

\begin {large}
%beginitem
1. Написати функцію, що знаходить у скількох цілих точках інтервалу $(a; b)$ значення функції~$sin(x)$ менше за значення функції~$cos(x)$.

\vspace{2mm}

%enditem

%beginitem
2. Написати функцію, що збільшує всі непарні елементи послідовності цілих чисел на 4.

\vspace{2mm}

%enditem

%beginitem
3. Написати функцію, що знаходить кількість входжень малих латинських літер у С-рядок.

\vspace{2mm}

%enditem

%beginitem
4. 
Написати функцію, що записує в масив дійсних послідовність чисел
$$\scalebox{1.2}{$C_{n+k}^{k}$}, \quad k=1,\ldots,n$$
Використовувати вкладені цикли та виклики функцій \textbf{заборонено}.

\vspace{2mm}

%enditem

%beginitem
5. Написати функцію, що повертає число з послідовності цілих, що має парну кількість різних простих дільників. Якщо таких чисел кілька, то повернути найбільше.

\vspace{2mm}

%enditem
\end{large}
%end
\newpage
%begin
%\begin {Huge}
{{27.11.2024 Контрольна робота \No 2, варіант  \No \textbf { 129 }\par}}

%\begin{center}Частина 2
%\end{center}
%\vspace{10mm}
%\end {Huge}

{
Якість коду та економність обчислень оцінюється по всім завданням. Послідовності задаються масивами. Використовувати додаткові структури даних (у тому числі рядки) та функції стандартної бібліотеки заборонено.

\vspace{2mm}

\begin {large}
%beginitem
1. Написати функцію, що знаходить найбільше значення функції~$sin(x)$ у цілих точках інтервалу $[a; b]$.

\vspace{2mm}

%enditem

%beginitem
2. Написати функцію, що знаходить суму непарних елементів послідовності цілих чисел.

\vspace{2mm}

%enditem

%beginitem
3. Написати функцію, що знаходить кількість входжень цифри 0 у С-рядок.

\vspace{2mm}

%enditem

%beginitem
4. 
Написати функцію, що записує в масив дійсних послідовність чисел
$$\scalebox{1.2}{$(-1)^{k}C_{2n+k+1}^{n+k}$}, \quad k=0,1,\ldots,n$$
Використовувати вкладені цикли та виклики функцій \textbf{заборонено}.

\vspace{2mm}

%enditem

%beginitem
5. Написати функцію, що в інтервалі $[a, b)$ знаходить ціле число, що ділиться на найбільший степінь числа 3. Якщо таких чисел декілька, повернути найбільше.

\vspace{2mm}

%enditem
\end{large}
%end
\vspace{35mm}


%begin
%\begin {Huge}
{{27.11.2024 Контрольна робота \No 2, варіант  \No \textbf { 130 }\par}}

%\begin{center}Частина 2
%\end{center}
%\vspace{10mm}
%\end {Huge}

{
Якість коду та економність обчислень оцінюється по всім завданням. Послідовності задаються масивами. Використовувати додаткові структури даних (у тому числі рядки) та функції стандартної бібліотеки заборонено.

\vspace{2mm}

\begin {large}
%beginitem
1. Написати функцію, що знаходить у скількох цілих точках інтервалу $[a; b)$ функція~$sin(x)$ приймає від'ємне значення.

\vspace{2mm}

%enditem

%beginitem
2. Написати функцію, що знаходить найбільший парний елемент послідовності цілих чисел.

\vspace{2mm}

%enditem

%beginitem
3. Написати функцію, що перевіряє, чи містить С-рядок хоч одну малу латинську літеру.

\vspace{2mm}

%enditem

%beginitem
4. 
Написати функцію, що записує в масив дійсних послідовність чисел
$$\scalebox{1.2}{$C_{n+k}^{2k-1}$}, \quad k=1,\ldots,n$$
Використовувати вкладені цикли та виклики функцій \textbf{заборонено}.

\vspace{2mm}

%enditem

%beginitem
5. Написати функцію, що повертає число з послідовності цілих, що має найбільшу кількість різних простих дільників. Якщо таких чисел кілька, то повернути найменше.

\vspace{2mm}

%enditem
\end{large}
%end
\newpage
%begin
%\begin {Huge}
{{27.11.2024 Контрольна робота \No 2, варіант  \No \textbf { 131 }\par}}

%\begin{center}Частина 2
%\end{center}
%\vspace{10mm}
%\end {Huge}

{
Якість коду та економність обчислень оцінюється по всім завданням. Послідовності задаються масивами. Використовувати додаткові структури даних (у тому числі рядки) та функції стандартної бібліотеки заборонено.

\vspace{2mm}

\begin {large}
%beginitem
1. Написати функцію, що знаходить найбільше значення функції~$sin(x)$ у цілих точках інтервалу $[a; b)$.

\vspace{2mm}

%enditem

%beginitem
2. Написати функцію, що знаходить суму елементів послідовності цілих чисел, що дорівнюють своїм індексам у масиві.

\vspace{2mm}

%enditem

%beginitem
3. Написати функцію, що перевіряє, яких символів у С-рядку більше: 0 чи 1.

\vspace{2mm}

%enditem

%beginitem
4. 
Написати функцію, що записує в масив дійсних послідовність чисел
$$\scalebox{1.2}{$(-1)^{k}C_{n+2k}^{k+1}$}, \quad k=1,\ldots,n$$
Використовувати вкладені цикли та виклики функцій \textbf{заборонено}.

\vspace{2mm}

%enditem

%beginitem
5. Написати функцію, яка для цілого числа знаходить кількість його різних простих дільників, що мають у десятковому запису цифру 1.

\vspace{2mm}

%enditem
\end{large}
%end
\vspace{35mm}


%begin
%\begin {Huge}
{{27.11.2024 Контрольна робота \No 2, варіант  \No \textbf { 132 }\par}}

%\begin{center}Частина 2
%\end{center}
%\vspace{10mm}
%\end {Huge}

{
Якість коду та економність обчислень оцінюється по всім завданням. Послідовності задаються масивами. Використовувати додаткові структури даних (у тому числі рядки) та функції стандартної бібліотеки заборонено.

\vspace{2mm}

\begin {large}
%beginitem
1. Написати функцію, що знаходить у скількох цілих точках інтервалу $(a; b]$ функція~$cos(x)$ приймає від'ємне значення.

\vspace{2mm}

%enditem

%beginitem
2. Написати функцію, що замінює всі від'ємні елементи послідовності цілих чисел нулями.

\vspace{2mm}

%enditem

%beginitem
3. Написати функцію, що перевіряє, чи містить С-рядок онакову кількість входжень символів ( та ).

\vspace{2mm}

%enditem

%beginitem
4. 
Написати функцію, що записує в масив дійсних послідовність чисел
$$\scalebox{1.2}{$(-1)^{k}C_{n+k}^{2k-1}$}, \quad k=1,\ldots,n$$
Використовувати вкладені цикли та виклики функцій \textbf{заборонено}.

\vspace{2mm}

%enditem

%beginitem
5. Написати функцію, що для цілого числа знаходить простий дільник, який входить у його розклад ц найбільшому степені. Якщо таких кілька, то повернути найбільший. \\ Наприклад, для числа $2^{9}{\cdot}3^5{\cdot}7^9{\cdot}11$ функція має повернути $7$.

\vspace{2mm}

%enditem
\end{large}
%end
\newpage
%begin
%\begin {Huge}
{{27.11.2024 Контрольна робота \No 2, варіант  \No \textbf { 133 }\par}}

%\begin{center}Частина 2
%\end{center}
%\vspace{10mm}
%\end {Huge}

{
Якість коду та економність обчислень оцінюється по всім завданням. Послідовності задаються масивами. Використовувати додаткові структури даних (у тому числі рядки) та функції стандартної бібліотеки заборонено.

\vspace{2mm}

\begin {large}
%beginitem
1. Написати функцію, що знаходить у скількох цілих точках інтервалу $(a; b)$ функція~$cos(x)$ приймає від'ємне значення.

\vspace{2mm}

%enditem

%beginitem
2. Написати функцію, що знаходить суму елементів послідовності дійсних чисел, що більші 1.

\vspace{2mm}

%enditem

%beginitem
3. Написати функцію, що перевіряє, чи містить С-рядок хоч одну велику латинську літеру.

\vspace{2mm}

%enditem

%beginitem
4. 
Написати функцію, що записує в масив дійсних послідовність чисел
$$\scalebox{1.2}{$(-1)^{k}\frac{C_{2k}^k}{k+1}$}, \quad k=0,1,\ldots,n-2$$
Використовувати вкладені цикли та виклики функцій \textbf{заборонено}.

\vspace{2mm}

%enditem

%beginitem
5. Написати функцію, що для цілого числа знаходить кількість його простих дільників, що входять у його розклад рівно у другому степеню. \\ Наприклад, для числа $3^2{\cdot}7^9{\cdot}11{\cdot}23^{2}$ таких дільників два – $3$ та $23$, тому функція має повернути $2$.

\vspace{2mm}

%enditem
\end{large}
%end
\vspace{35mm}


%begin
%\begin {Huge}
{{27.11.2024 Контрольна робота \No 2, варіант  \No \textbf { 134 }\par}}

%\begin{center}Частина 2
%\end{center}
%\vspace{10mm}
%\end {Huge}

{
Якість коду та економність обчислень оцінюється по всім завданням. Послідовності задаються масивами. Використовувати додаткові структури даних (у тому числі рядки) та функції стандартної бібліотеки заборонено.

\vspace{2mm}

\begin {large}
%beginitem
1. Написати функцію, що знаходить у скількох цілих точках інтервалу $[a; b]$ функція~$cos(x)$ приймає додатнє значення.

\vspace{2mm}

%enditem

%beginitem
2. Написати функцію, що знаходить середнє арифметичне елементів послідовності цілих чисел.

\vspace{2mm}

%enditem

%beginitem
3. Написати функцію, що перевіряє, чи містить С-рядок хоч один символ підкреслення.

\vspace{2mm}

%enditem

%beginitem
4. 
Написати функцію, що записує в масив дійсних послідовність чисел
$$\scalebox{1.2}{$(-1)^{k}C_{n+k}^{k+1}$}, \quad k=0,1,\ldots,n$$
Використовувати вкладені цикли та виклики функцій \textbf{заборонено}.

\vspace{2mm}

%enditem

%beginitem
5. Написати функцію, що в інтервалі $[a, b)$ знаходить ціле число, що ділиться на найбільший степінь числа 3. Якщо таких чисел декілька, повернути найменше.

\vspace{2mm}

%enditem
\end{large}
%end
\newpage
%begin
%\begin {Huge}
{{27.11.2024 Контрольна робота \No 2, варіант  \No \textbf { 135 }\par}}

%\begin{center}Частина 2
%\end{center}
%\vspace{10mm}
%\end {Huge}

{
Якість коду та економність обчислень оцінюється по всім завданням. Послідовності задаються масивами. Використовувати додаткові структури даних (у тому числі рядки) та функції стандартної бібліотеки заборонено.

\vspace{2mm}

\begin {large}
%beginitem
1. Написати функцію, що знаходить у скількох цілих точках інтервалу $(a; b]$ значення функції~$sin(x)$ менше за значення функції~$cos(x)$.

\vspace{2mm}

%enditem

%beginitem
2. Написати функцію, що замінює всі додатні елементи послідовності цілих чисел числом 15.

\vspace{2mm}

%enditem

%beginitem
3. Написати функцію, що знаходить кількість входжень літери Z у С-рядок.

\vspace{2mm}

%enditem

%beginitem
4. 
Написати функцію, що записує в масив дійсних послідовність чисел
$$\scalebox{1.2}{$C_{n+2k}^{k}$}, \quad k=1,\ldots,n$$
Використовувати вкладені цикли та виклики функцій \textbf{заборонено}.

\vspace{2mm}

%enditem

%beginitem
5. Написати функцію, що для інтервалу $[a, b)$ знаходить найближчу пару чисел з цього інтервалу, що мають не менше трьох простих дільників та мають цифру 3 у своєму десятковому запису.

\vspace{2mm}

%enditem
\end{large}
%end
\vspace{35mm}


%begin
%\begin {Huge}
{{27.11.2024 Контрольна робота \No 2, варіант  \No \textbf { 136 }\par}}

%\begin{center}Частина 2
%\end{center}
%\vspace{10mm}
%\end {Huge}

{
Якість коду та економність обчислень оцінюється по всім завданням. Послідовності задаються масивами. Використовувати додаткові структури даних (у тому числі рядки) та функції стандартної бібліотеки заборонено.

\vspace{2mm}

\begin {large}
%beginitem
1. Написати функцію, що знаходить у скількох цілих точках інтервалу $(a; b)$ функція~$sin(x)$ приймає від'ємне значення.

\vspace{2mm}

%enditem

%beginitem
2. Написати функцію, що замінює всі елементи на непарних позиціях послідовності цілих чисел (позиції нумеруються з 0) на число 25.

\vspace{2mm}

%enditem

%beginitem
3. Написати функцію, що перевіряє, чи містить С-рядок онакову кількість входжень символів ( та ).

\vspace{2mm}

%enditem

%beginitem
4. 
Написати функцію, що записує в масив дійсних послідовність чисел
$$\scalebox{1.2}{$C_{2k+1}^{k}$}, \quad k=2,\ldots,n$$
Використовувати вкладені цикли та виклики функцій \textbf{заборонено}.

\vspace{2mm}

%enditem

%beginitem
5. Написати функцію, що для цілого числа знаходить простий дільник, який входить у його розклад у найбільшому степені. Якщо таких кілька, то повернути найменший. \\ Наприклад, для числа $2^{9}{\cdot}3^5{\cdot}7^9{\cdot}11$ функція має повернути $2$.

\vspace{2mm}

%enditem
\end{large}
%end
\newpage
%begin
%\begin {Huge}
{{27.11.2024 Контрольна робота \No 2, варіант  \No \textbf { 137 }\par}}

%\begin{center}Частина 2
%\end{center}
%\vspace{10mm}
%\end {Huge}

{
Якість коду та економність обчислень оцінюється по всім завданням. Послідовності задаються масивами. Використовувати додаткові структури даних (у тому числі рядки) та функції стандартної бібліотеки заборонено.

\vspace{2mm}

\begin {large}
%beginitem
1. Написати функцію, що знаходить у скількох цілих точках інтервалу $(a; b)$ значення функції~$sin(x)$ більше за значення функції~$cos(x)$.

\vspace{2mm}

%enditem

%beginitem
2. Написати функцію, що замінює всі парні елементи послідовності цілих чисел числом 13.

\vspace{2mm}

%enditem

%beginitem
3. Написати функцію, що перевіряє, яких символів у С-рядку більше: 0 чи 1.

\vspace{2mm}

%enditem

%beginitem
4. 
Написати функцію, що записує в масив дійсних послідовність чисел
$$\scalebox{1.2}{$(-1)^{k}C_{n+k}^{k}$}, \quad k=1,\ldots,n-2$$
Використовувати вкладені цикли та виклики функцій \textbf{заборонено}.

\vspace{2mm}

%enditem

%beginitem
5. Написати функцію, що для цілого числа знаходить кількість його простих дільників, що входять у його розклад не більш ніж у 5му степеню. \\ Наприклад, для числа $3^8{\cdot}7^9{\cdot}11{\cdot}23^{2}$ таких дільників два – $11$ та $23$, тому функція має повернути $2$.

\vspace{2mm}

%enditem
\end{large}
%end
\vspace{35mm}


%begin
%\begin {Huge}
{{27.11.2024 Контрольна робота \No 2, варіант  \No \textbf { 138 }\par}}

%\begin{center}Частина 2
%\end{center}
%\vspace{10mm}
%\end {Huge}

{
Якість коду та економність обчислень оцінюється по всім завданням. Послідовності задаються масивами. Використовувати додаткові структури даних (у тому числі рядки) та функції стандартної бібліотеки заборонено.

\vspace{2mm}

\begin {large}
%beginitem
1. Написати функцію, що знаходить у скількох цілих точках інтервалу $(a; b]$ функція~$sin(x)$ приймає від'ємне значення.

\vspace{2mm}

%enditem

%beginitem
2. Написати функцію, що замінює всі елементи на парних позиціях послідовності цілих чисел (позиції нумеруються з 0) на число 20.

\vspace{2mm}

%enditem

%beginitem
3. Написати функцію, що знаходить кількість входжень великих латинських літер у С-рядок.

\vspace{2mm}

%enditem

%beginitem
4. 
Написати функцію, що записує в масив дійсних послідовність чисел
$$\scalebox{1.2}{$C_{n+2k}^{k+1}$}, \quad k=1,\ldots,n$$
Використовувати вкладені цикли та виклики функцій \textbf{заборонено}.

\vspace{2mm}

%enditem

%beginitem
5. Написати функцію, що для послідовності цілих знаходить довжину найдовшого відрізку, в якому всі сусідні числа різні.

\vspace{2mm}

%enditem
\end{large}
%end
\newpage
%begin
%\begin {Huge}
{{27.11.2024 Контрольна робота \No 2, варіант  \No \textbf { 139 }\par}}

%\begin{center}Частина 2
%\end{center}
%\vspace{10mm}
%\end {Huge}

{
Якість коду та економність обчислень оцінюється по всім завданням. Послідовності задаються масивами. Використовувати додаткові структури даних (у тому числі рядки) та функції стандартної бібліотеки заборонено.

\vspace{2mm}

\begin {large}
%beginitem
1. Написати функцію, що знаходить у скількох цілих точках інтервалу $[a; b]$ значення функції~$sin(x)$ менше за значення функції~$cos(x)$.

\vspace{2mm}

%enditem

%beginitem
2. Написати функцію, що знаходить найбільший непарний елемент послідовності цілих чисел.

\vspace{2mm}

%enditem

%beginitem
3. Написати функцію, що перевіряє, чи містить С-рядок хоч одну малу латинську літеру.

\vspace{2mm}

%enditem

%beginitem
4. 
Написати функцію, що записує в масив дійсних послідовність чисел
$$\scalebox{1.2}{$(-1)^{k}C_{3k-1}^{k}$}, \quad k=2,\ldots,n$$
Використовувати вкладені цикли та виклики функцій \textbf{заборонено}.

\vspace{2mm}

%enditem

%beginitem
5. Написати функцію, яка для цілого числа знаходить кількість його різних простих дільників, що мають у десятковому запису цифру 1.

\vspace{2mm}

%enditem
\end{large}
%end
\vspace{35mm}


%begin
%\begin {Huge}
{{27.11.2024 Контрольна робота \No 2, варіант  \No \textbf { 140 }\par}}

%\begin{center}Частина 2
%\end{center}
%\vspace{10mm}
%\end {Huge}

{
Якість коду та економність обчислень оцінюється по всім завданням. Послідовності задаються масивами. Використовувати додаткові структури даних (у тому числі рядки) та функції стандартної бібліотеки заборонено.

\vspace{2mm}

\begin {large}
%beginitem
1. Написати функцію, що знаходить у скількох цілих точках інтервалу $[a; b)$ функція~$sin(x)$ приймає додатнє значення.

\vspace{2mm}

%enditem

%beginitem
2. Написати функцію, що знаходить суму елементів послідовності дійсних чисел, що більші за свої індекси у масиві.

\vspace{2mm}

%enditem

%beginitem
3. Написати функцію, що перевіряє, чи містить С-рядок хоч одну велику латинську літеру.

\vspace{2mm}

%enditem

%beginitem
4. 
Написати функцію, що записує в масив дійсних послідовність чисел
$$\scalebox{1.2}{$\frac{C_{2k}^k}{k+1}$}, \quad k=0,1,\ldots,n$$
Використовувати вкладені цикли та виклики функцій \textbf{заборонено}.

\vspace{2mm}

%enditem

%beginitem
5. Написати функцію, що для інтервалу $[a, b)$ знаходить найближчу пару чисел з цього інтервалу, що кожне з них має парну кількість різних простих дільників. Якщо таких пар кілька, то цікавить пара з найбільшою сумою.

\vspace{2mm}

%enditem
\end{large}
%end
\newpage
\end{document}

